\fancyhead[C]{{\LARGE Proper of Saints}}
\fancyhead[RO,LE]{}
\fancyhead[RE,LO]{}
\phantomsection
\addcontentsline{toc}{section}{Through the Year}

\begin{rubric}
	Propers for Memorials are indicated by parentheses: ().
\end{rubric}

\begin{rubric}
	When a Common is specified, it is used for the propers, except for anything provided in the Feast Day itself.
	
	\textsc{Note,} Unless otherwise specified, the first options within a Common are used.
\end{rubric}

%A. Apostles
%Ab. Abbot
%E. Evangelist
%M. Martyr(ess)
%V. Virgin
%B. Bishop
%C. Confessor
%D. Doctor
%W. Widow
%Ma. Matron
%Em. Empress

\festival{Common Vigil}{29 November}{Vigil of St. Andrew}{A.}{LentenViolet}
\feastday{St. Andrew Vigil}
\feastdate{29}{November}

\lett{W}{e} beseech thee, Almighty God: that the blessed Apostle Andrew, whose festival we prevent, may implore thy help for us; that we, being absolved from our offences, may likewise be delivered from all dangers. Through.

\begin{rubric}
	Commemoration of the Feria (if it be Advent) \& \blackrubric{of St. Saturninus} (p. \pageref{SaturninusCollect}).
\end{rubric}


\subbymemorial{29 November}{St. Saturninus}{B.M.}{MartyrRed}
\feastday{{St. Saturninus}}
\feastdate{29}{November}

\memorialCollect\label{SaturninusCollect}
\lett{O}{God,} who vouchsafest unto us to rejoice in the birthday of thy blessed Martyr Saturninus: grant, we pray thee; that we may be succoured by his merits. Through.


\festival{Second Class Double}{30 November}{St. Andrew}{A.}{FestiveWhite}
\feastday{St. Andrew}
\feastdate{30}{November}

\begin{rubric}
	Common of Apostles (p. \pageref{CommonApostles}).
\end{rubric}

\openingSentence{The Lord loved Andrew in the odour of sweetness.}

%\begin{paracol}{2}[]
%\sloppy
\firstEvensong
\psalmAntiphon{Hail, precious Cross, {\dag} receive the disciple of him who hanged on thee, my Master Christ.}

\firstEvensongAntiphon{One of the two {\dag} which followed the Lord was Andrew, Simon Peter's brother, alleluia.}

%\switchcolumn

\mattins
\psalmAntiphon{Blessed Andrew {\dag} prayed saying: O Lord King of eternal glory, sustain me, as I hang upon the tree.}

\mattinsAntiphon{Yield up to us {\dag} a man so righteous, restore to us a man so holy: destroy not a man so dear to God, righteous, dutiful, and gentle.}

%\fussy
%\end{paracol}

\secondEvensong
\psalmAntiphon{Thou hast overwhelmed in hell, O Lord, {\dag} them that persecute the righteous: and thou wast leader on the wood of the Cross.}

\secondEvensongAntiphon{When blessed Andrew {\dag} came to the place where the Cross had been prepared, he cried out and said: O goodly Cross, so long desired, and now made ready for my eager spirit; fearless and joyful do I come to thee: therefore do thou also receive me gladly, as his disciple, who did hang upon thee.}

\collect
\lett{W}{e} humbly entreat thy Majesty, O Lord: that as blessed Andrew the Apostle was to thy Church a preacher and governor; so he may be a perpetual intercessor for us in thy sight. Through.

\begin{rubric}
    In Advent, Commemoration of the Feria.
\end{rubric}


\festival{Double (m.t.v.)}{2 December}{St. Peter Chrysologus}{B.C.D.}{FestiveWhite}
\feastday{{St. Peter Chrysologus}}
\feastdate{2}{December}

\begin{rubric}
	Common of a Confessor Bishop (p. \pageref{CommonConfessorBishop}).
\end{rubric}

\lett{O}{God,} who by divine foreshewing wast pleased to choose blessed Peter Chrysologus thy illustrious Doctor to rule and instruct thy Church: grant, we beseech thee; that, as we have had him for a Doctor of life on earth, so we may be found worthy to have him for an intercessor in heaven. Through.

\begin{rubric}
	Commemoration \blackrubric{of St. Bibiana} \& of the Advent Feria.
\end{rubric}


\subbymemorial{2 December}{St. Bibiana}{V.M.}{MartyrRed}
\feastday{{St. Bibiana}}
\feastdate{2}{December}

\memorialCollect
\lett{O}{God,} the giver of all good gifts, who in thy handmaid Bibiana didst unite the palm of martyrdom with the flower of virginity: unite by her intercession our hearts in charity with thee; that all perils being done away, we may attain unto everlasting rewards. Through.


\subbymemorial{4 December}{St. Barbara}{V.M.}{MartyrRed}
\feastday{{St. Barbara}}
\feastdate{4}{December}

\begin{rubric}
	Common of Virgins (p. \pageref{VirginMartyressCollect}).
\end{rubric}


%MANUAL ADJUSTMENT:
\clearpage
\subbymemorial{5 December}{St. Sabbas of Jud{\ae}a}{Ab.}{FestiveWhite}
\feastday{{St. Sabbas}}
\feastdate{5}{December}

\begin{rubric}
	Common of a Confessor not a Bishop (p. \pageref{AbbotCollect}).
\end{rubric}


\festival{Double}{6 December}{St. Nicholas}{B.C.}{FestiveWhite}
\feastday{{St. Nicholas}}
\feastdate{6}{December}

\begin{rubric}
	Common of a Confessor Bishop (p. \pageref{CommonConfessorBishop}).
\end{rubric}

\lett{O}{God,} who didst adorn thy blessed Bishop Nicholas with innumerable miracles: grant, we beseech thee; that by his merits and prayers we may be delivered from the fires of hell. Through.


\festival{Greater Double (m.t.v.)}{7 December}{St. Ambrose}{B.C.D.}{FestiveWhite}
\feastday{{St. Ambrose}}
\feastdate{7}{December}

\begin{rubric}
	Common of a Confessor Bishop (p. \pageref{CommonConfessorBishop}).
\end{rubric}

\lett{O}{God,} who didst give blessed Ambrose unto thy people to be a minister of everlasting salvation: grant, we beseech thee; that as we have learned of him the doctrine of life on earth, so we may be found worthy to have him for our advocate in heaven. Through.

\begin{rubric}
	Commemoration of the Advent Feria.
\end{rubric}


\festivalLL{Second Class Double with a Common Octave}{8 December}{Conception of the Blessed Virgin Mary}\label{ConceptionBVM}
\feastday{Conception B.V.M.}
\feastdate{8}{December}

\begin{rubric}
	Common of the Blessed Virgin Mary (p. \pageref{CommonBVM}).
\end{rubric}

\begin{paracol}{2}[]
\sloppy

\firstEvensong
\psalmAntiphon{Thou art all fair, O Mary; {\dag} there is no spot in thee.}\\

℣. To-day is the Conception of the holy Virgin Mary.

℟. Whose glorious life illumineth all the churches.

\firstEvensongAntiphon{All generations {\dag} shall call me blessed: for he that is mighty hath magnified me, alleluia.}

\switchcolumn

\mattins
\psalmAntiphon{Thou art the exaltation of Jerusalem, {\dag} thou art the glory of Israel, thou art the great rejoicing of our nation.}\\

℣. To-day is the Conception of the holy Virgin Mary.

℟. Whose glorious life illumineth all the churches.

\properantiphon{Ben.}{The Lord God said {\dag} unto the serpent, I will put enmity between thee and the woman, and between thy seed and her seed; it shall bruise thy head, alleluia.}

\fussy	
\end{paracol}

\secondEvensong
\psalmAntiphon{Blessed art thou, {\dag} O Virgin Mary, of the Most High God, above all the women upon the earth.}\\

℣. To-day is the Conception of the holy Virgin Mary.

℟. Whose glorious life illumineth all the churches.

\secondEvensongAntiphon{Let us celebrate {\dag} the worshipful Conception of the blessed and glorious Virgin Mary, whose lowliness the Lord regarded when at the word of an Angel she conceived the world's Redeemer, alleluia.}


\collect
\lett{O}{God,} mercifully hear the supplication of thy servants; that we who are assembled together on the Conception of the Virgin Mother of God, may at her intercession be delivered by thee from the dangers which beset us. Through.
\begin{rubric}
    Commemoration of the Advent Feria.
\end{rubric}


\subbymemorial{10 December}{Pope St. Melchiades of Rome}{B.M.}{MartyrRed}
\feastday{{Pope St. Melchiades}}
\feastdate{10}{December}

\begin{rubric}
	Common of a Martyr (p. \pageref{CommonMartyrNotBishopCollect}).
\end{rubric}


\subbymemorial{11 December}{Pope St. Damasus of Rome}{B.C.}{FestiveWhite}
\feastday{{Pope St. Damasus}}
\feastdate{11}{December}

\memorialCollect
\lett{G}{raciously} hear our prayers, O Lord: and at the intercession of blessed Damasus, thy Confessor and Bishop, mercifully grant us pardon and peace. Through.

\lettspace

\begin{rubric}
	Commemoration of the Octave \& Advent Feria.
\end{rubric}


%MANUAL ADJUSTMENT:
\clearpage
\festival{Greater Double}{13 December}{St. Lucy}{V.M.}{MartyrRed}
\feastday{{St. Lucy}}
\feastdate{13}{December}

\begin{rubric}
	Common of Virgins (p. \pageref{CommonVirgin}).
\end{rubric}


\begin{paracol}{2}[]
\sloppy

\firstEvensong
\psalmAntiphon{While holy Lucy {\dag} prayed, there appeared unto her the blessed Agatha, who consoled the handmaid of Christ.}\\

\firstEvensongAntiphon{In thy patience {\dag} thou didst possess thy soul, O Lucy, spouse of Christ: thou didst hate the things which are in the world, and thou shinest among the Angels: resisting unto blood, thou didst vanquish the enemy.}

\switchcolumn

\mattins
\psalmAntiphon{Through thee, O Virgin Lucy, {\dag} the Syracusan city shall be glorified by the Lord Jesus Christ.}\\

\mattinsAntiphon{A pillar art thou {\dag} that may not be moved, O Lucy, spouse of Christ: and all the people are waiting until thou receive the crown of life, alleluia.}

\end{paracol}

\secondEvensong
\psalmAntiphon{I bless thee, {\dag} O Father of my Lord Jesus Christ: for by thy Son the fire is quenched round about thine handmaid.}\\

℣. Full of grace are thy lips.

℟. Because God hath blessed thee for ever.

\secondEvensongAntiphon{With such gravity {\dag} was she endued by the Holy Spirit, that the virgin of the Lord remained unmoved.}

\collect
\lett{G}{raciously} hear us, O God of our salvation: that, like as we do rejoice in the festival of blessed Lucy thy Virgin; so we may be instructed in all godly and devout affection. Through.

\begin{rubric}
	Commemoration \blackrubric{of St. Herman of Alaska} from the Common of a Confessor not a Bishop (p. \pageref{ConfessorNotBishopCollect}), the Octave, \& the Advent Feria.
\end{rubric}


\subbymemorial{13 December}{St. Herman of Alaska (m.t.v.)}{C.}{FestiveWhite}
\feastday{{St. Herman}}
\feastdate{13}{December}

\begin{rubric}
	Common of a Confessor not a Bishop (p. \pageref{ConfessorNotBishopCollect}).
\end{rubric}


%MANUAL ADJUSTMENT:
\clearpage
\festivalLL{Semidouble}{14 December}{Day VII within the Octave of the Conception}
\feastday{{Conception Octave}}
\feastdate{14}{December}

\begin{rubric}
	Of the Octave, as on 9 December. But if Ember Wednesday fall on this day, of the Ember Day, with a Commemoration of the Octave.
\end{rubric}

%\begin{rubric}
%	If an Ember Day occur on any of the following Feasts, the Commemoration of the Feria is omitted.
%\end{rubric}


\festivalLL{Greater Double}{15 December}{Octave Day of the Conception of the Blessed Virgin Mary}
\feastday{{Conception Octave Day}}
\feastdate{15}{December}

\begin{rubric}
	As on the Feast (p. \pageref{ConceptionBVM}), Commemoration of the Advent Feria.
\end{rubric}


\subbymemorial{16 December}{St. Eusebius of Vercelli}{B.M.}{MartyrRed}
\feastday{{St. Eusebius}}
\feastdate{16}{December}

\begin{rubric}
	Common of a Martyr (p. \pageref{CommonMartyrBishopCollect}), second Collect.
\end{rubric}


\festivalLL{Greater Double}{18 December}{Expectation of the Blessed Virgin Mary}
\feastday{{Expectation B.V.M.}}
\feastdate{18}{December}

\begin{rubric}
	Common of the Blessed Virgin Mary (p. \pageref{CommonBVM}).
\end{rubric}

\openingSentence{Send ye the lamb to the ruler of the land from Sela to the wilderness, unto the mount of the daughter of Zion.\novr{Is. 16:1}}

%MANUAL ADJUSTMENT:
\vspace{-0.5\baselineskip}

\begin{paracol}{2}[]
\sloppy

\firstEvensong

\psalmAntiphon{Hail Mary, full of grace;~{\dag} The Lord is with thee; Blessed art thou amongst women. Alleluia.}\\

℣. Hail Mary, full of grace.

℟. The Lord is with thee.

\secondEvensong

\psalmAntiphon{The angel Gabriel was sent to the Virgin Mary, espoused to Joseph.}\\

℣. Hail Mary, full of grace.

℟. The Lord is with thee.

\switchcolumn

\mattins

\psalmAntiphon{Fear not, Mary:~{\dag} for thou hast found favour with God. Behold, thou shalt conceive in thy womb, and bring forth a son. Alleluia.}\\

℣. The Holy Ghost shall come upon thee.

℟. And the power of the Highest shall overshadow thee.

\antiphon{Ben.}{He shall sit upon the throne of David, {\dag} and of his kingdom, for ever.}

\fussy
\end{paracol}


\collect
\lett{O}{God,} who wast pleased that thy Word should take flesh of the womb of the Blessed Virgin Mary at the message of an Angel: grant to thy humble servants; that we who believe her to be truly the Mother of God may be aided by her intercession in thy sight. Through the same.


\festival{Vigil}{20 December}{Vigil of St. Thomas}{A.}{LentenViolet}
\feastday{{St. Thomas Vigil}}
\feastdate{20}{December}

%\begin{rubric}
%	If today be Sunday, in the Ember Saturday Mass, Commemoration is made of the anticipated Vigil of St. Thomas and the last Gospel of the Vigil is read at the end, the \nth{3} Collect is \blackrubric{of St. Mary in Advent} (p. \pageref{SPMaryInAdvent}).
%\end{rubric}

\begin{rubric}
	Common of Vigils of the Apostles (p. \pageref{CommonVigilApostles}), Commemoration of the Advent Feria.
\end{rubric}


\festival{Second Class Double}{21 December}{St. Thomas}{A.}{FestiveWhite}
\feastday{St. Thomas}
\feastdate{21}{December}

\begin{rubric}
	Common of Apostles (p. \pageref{CommonApostles}).
\end{rubric}

\properantiphon{Mag. \& Ben.}{Because thou hast seen me, {\dag} Thomas, thou hast believed: blessed are they that have not seen, and yet have believed, alleluia.}

\collect
\lett{G}{rant} us, we beseech thee, O Lord, to glory in the solemnity of thy blessed Apostle Thomas: that we may ever be succoured by his protection; and follow his faith with worthy devotion. Through.

\begin{rubric}
    Commemoration of the Advent Feria.
\end{rubric}


\subbymemorial{10 January}{St. Paul the First Hermit}{C.}{FestiveWhite}
\feastday{{St. Paul Hermit}}
\feastdate{10}{January}

\memorialCollect
\lett{O}{God,} who makest us glad with the yearly solemnity of blessed Paul, thy Confessor: mercifully grant; that, as we now celebrate his birthday, so we may follow the example of his life. Through.


\subbymemorial{11 January}{Pope St. Hyginus of Rome}{B.M.}{MartyrRed}
\feastday{{Pope St. Hyginus Rome}}
\feastdate{11}{January}

\begin{rubric}
	Common of a Martyr (p. \pageref{CommonMartyrBishopCollect}).
\end{rubric}


%MANUAL ADJUSTMENT:
\clearpage
\subbymemorial{12 January}{St. Benedict Biscop}{Ab.}{FestiveWhite}
\feastday{{St. Benedict Biscop}}
\feastdate{12}{January}

\memorialCollect
\lett{O}{God,} by whose gift the blessed Abbot Benedict left all things that he might be made perfect: grant unto all those who have entered upon the path of evangelical perfection, that they may neither look back nor linger in the way; but hastening to thee without stumbling, may lay hold upon eternal life. Through.


\festival{Double}{14 January}{St. Hilary}{B.C.D.}{FestiveWhite}
\feastday{{St. Hilary}}
\feastdate{14}{January}

\begin{rubric}
	Common of a Confessor Bishop (p. \pageref{CommonConfessorBishop}).
\end{rubric}

\begin{rubric}
	\textsc{Note,} Commemoration \blackrubric{of St. Felix} (p. \pageref{FelixCollect}).
\end{rubric}


\subbymemorial{14 January}{St. Felix of Nola}{M.}{MartyrRed}
\feastday{{St. Felix Nola}}
\feastdate{14}{January}

%\begin{rubric}
%	The Common of a Martyr (p. \pageref{CommonMartyr}).\par
%\end{rubric}

\memorialCollect\label{FelixCollect}
\lett{G}{rant,} we beseech thee, Almighty God, that the examples of thy Saints may provoke us to a better life: that as we celebrate their festival so we may imitate their actions. Through.


\festival{Greater Double}{15 January}{St. Maurus}{Ab.}{FestiveWhite}
\feastday{{St. Maurus}}
\feastdate{15}{January}

\begin{paracol}{2}[]
\sloppy

%MANUAL:
\vspace{-1\baselineskip}

\firstEvensong

\psalmAntiphon{The blessed Maurus, {\dag} born of a renowned and noble house, from his boyhood esteemed the reproach of Christ greater riches than the treasures of the world.}

\firstEvensongHymn\label{MaurusEvensong}

\hymnLett{D}{efender,} leader true, thine own companions deem'd

\hymnSecondLine{Thy splendour half divine, since all were less esteem'd:}

\begin{hangparas}{1.25em}{1}
For thy most worthy deeds, Maurus, accept the lays

Wherewith we celebrate thy praise.\\

Born of a noble stock, great honour was his due,

But palaces he spurn'd, and from the world withdrew;

Delights he trampl'd down, estates and robes unpric'd,

To undergo the yoke of Christ.\\

The holy Abbot's grace, before his eyes display'd,

By deeds of equal worth he eagerly portray'd;

The pattern of the life monastic shone in truth

From every action of the youth.\\

Sternly, with sackcloth rough, self-mastery he wrought,

And by the curb of law, unbroken silence sought;

The ever-watchful nights in fervent prayer he spent;

Whole days of fasting underwent\\

Right speedily he flew to do the father's hest,

Dry-shod, the waters deep with fearless feet he press'd;

And safely he return'd with Placidus, set free

Like Peter walking on the sea.\\

To thee, O Trinity, high praise and honour be,

Whose countenance desir'd the heaven-dwellers see:

Grant that the Holy Rule may be our pathway plain

The prize of Maurus to attain.  Amen.\\
\end{hangparas}

	℣. The Lord loved him and adorned him.

	℟. He clothed him with a robe of glory.
	
	\antiphon{Mag.}{O most blessed of men! {\dag} who, rejecting this world, bore the yoke of Holy Rule from tender years so lovingly; and being made obedient even unto death, he denied himself, that he might wholly cling to Christ his Master, alleluia.}
	
	\switchcolumn

%MANUAL:
\vspace{-1\baselineskip}

\mattins

\psalmAntiphon{Upheld by wings of obedience, {\dag} he walked upon the waters; and borne by the Spirit of God, he was saved from sinking in the flood.}

\mattinsInvitatoryHymn\label{MaurusInvitatory}
\hymnLett{I}{n} childhood Placidus was by his father giv'n;

\hymnSecondLine{Thus offer'd, he himself did freely yield to God.}

\begin{hangparas}{1.25em}{1}
Since first he came, he ever shone with wondrous grace,

A pattern to all zealous souls.\\

His Abbot sent him forth to fill an earthen crock,

He dips it in the lake, unwary, slips and falls,

A wave then carries him a bow-shot from the shore,

Out in the deep drowning is nigh.\\

But Maurus, what is this, that hast'ning to the lake,

Thou runnest o'er the waves as though thou wast on land?

Wont to obey, `tis thus thy holy father's voice

Doth lead thee to a miracle.\\

Straightway the lake restores Placidus safe again,

But whose the merit? Did his Nursian father draw

Him from the swirling depths, or was it Maurus' act?

The child resolves their questioning.\\

O Holy Trinity, through prayers of Placidus,

Grant to thy monks that by the narrow path of Rule

They may at length attain unto the courts of heav'n,

And mingle with celestial choirs. Amen.\\
\end{hangparas}

\mattinsOfficeHymn

\begin{rubric}
	Common of a Confessor not a Bishop (p. \pageref{CommonConfessorNotBishop}), the Versicle \& Antiphon of I Evensong.
\end{rubric}

\secondEvensong

\psalmAntiphon{He was chosen {\dag} by the Lord to be an example to the cloistered, and a chief observer of the Holy Rule.}

\secondEvensongHymn

\begin{rubric}
	Office Hymn of I Evensong.
\end{rubric}

℣. The Lord guided the righteous in right paths.

℟. And shewed him the kingdom of God.

\antiphon{Mag.}{To-day holy Maurus, {\dag} lying upon a goat-skin, died happily before the altar; to-day the first-begotten disciple of blessed Benedict, through the guiding of the Holy Rule, came up to Christ, rising untroubled, accompanied by choirs of Angels; today the obedient man, telling his victories, was worthy to be crowned by the Lord, alleluia.}

\fussy
\end{paracol}

\collect
\lett{O}{God,} who for a pattern of obedience didst cause blessed Maurus to walk dry-shod upon the waters: grant that we may both follow perfectly the example of his virtues, and also be worthy to share in his reward. Through.


\subbymemorial{16 January}{Pope St. Marcellus I of Rome}{B.M.}{MartyrRed}
\feastday{{Pope St. Marcellus I Rome}}
\feastdate{16}{January}

\memorialCollect
\lett{O}{Lord,} we beseech thee favourably to hear the prayers of thy people: that, as we do rejoice in the passion of blessed Marcellus thy Martyr and Bishop, so we may be succoured by his merits. Through.


\festival{Double}{17 January}{St. Anthony}{Ab.}{FestiveWhite}
\feastday{{St. Anthony}}
\feastdate{17}{January}

\begin{rubric}
	Common of a Confessor not a Bishop (p. \pageref{AbbotCollect}).
\end{rubric}


%MANUAL ADJUSTMENT:
\clearpage
\festival{Greater Double}{18 January}{Chair of St. Peter at Rome}{A.}{FestiveWhite}
\feastday{{Chair St. Peter Rome}}
\feastdate{18}{January}

\begin{rubric}
	As in the Feast of the Chair of St. Peter at Antioch (p. \pageref{CathedraAntioch}).
\end{rubric}

%\collect
\lett{O}{God,} who didst bestow upon thy blessed Apostle Peter the keys of the kingdom of heaven, and didst appoint unto him the high priesthood of binding and loosing: vouchsafe; that by the help of his intercession we may be delivered from the bonds of our iniquities. Who livest and reignest.

\lett{O}{God,} who by the preaching of the blessed Apostle Paul didst teach the multitude of the Gentiles: grant to us, we beseech thee; that we who celebrate his commemoration may know him to be our advocate with thee. (Through.)

\begin{rubric}
    Commemoration \blackrubric{of St. Prisca of Rome} (p. \pageref{PriscaCollect}).
\end{rubric}


\subbymemorial{18 January}{St. Prisca of Rome}{V.M.}{MartyrRed}
\feastday{{St. Prisca Rome}}
\feastdate{18}{January}

%\begin{rubric}
%	The Common of Virgins (p. \pageref{CommonVirgin}).
%\end{rubric}

\memorialCollect\label{PriscaCollect}
\lett{G}{rant,} we beseech thee, Almighty God: that we who celebrate the heavenly birthday of blessed Prisca, thy Virgin and Martyress; may both rejoice in her yearly solemnity, and profit by the example of so great a faith. Through.


\subbymemorial{19 January}{Sts. Marius, Martha, Audifax, \& Abachum}{M.}{MartyrRed}
\feastday{{Sts. Marius \&c.}}
\feastdate{19}{January}

\memorialCollect
\lett{G}{raciously} hear thy people, O Lord, who call upon thee with the advocacy of thy Saints: and grant us both to rejoice in peace in this temporal life; and to find succour unto life everlasting. Through.

\begin{rubric}
    Commemoration \blackrubric{of St. Mark of Ephesus} (p. \pageref{EphesianCollect}).
\end{rubric}


%PROVISIONAL:
\subbymemorial{19 January}{St. Mark of Ephesus (m.t.v.)}{B.C.}{FestiveWhite}
\feastday{{St. Mark Ephesus}}
\feastdate{19}{January}

%\begin{rubric}
%	The Common of a Confessor Bishop (p. \pageref{CommonConfessorBishop}), second Collect.
%\end{rubric}
%
%\par\noindent
%\textit{Yet to be approved.}
%
%\vspace{-1ex}
%
%\begin{tcolorbox}[
%  colback=white,           % background color
%  colframe=black,          % frame color
%  fonttitle=\itshape,     % optional font for title
%  title style={left=2mm},  % optional title left alignment
%  attach title to upper={},
%  enhanced,
%  boxed title style={size=small,colframe=black!50!white,colback=white},
%]
%\collect
\memorialCollect\label{EphesianCollect}
\lett{O}{Lord} Jesu Christ, the only Shepherd and Bishop of our souls; we beseech thee to keep us in thy truth, preserve us by thy grace, and govern us according to thy loving-kindness. That as we imitate the good example of blessed Mark thy Confessor and Bishop, we may grow in love for thy Word and service of thy Church. Who with.
%\end{tcolorbox}


%MANUAL ADJUSTMENT:
\clearpage
\festival{Double}{20 January}{Sts. Bishop Fabian \& Sebastian}{M.}{MartyrRed}
\feastday{{Sts. Fabian \& Sebastian}}
\feastdate{20}{January}

\begin{rubric}
	Common of Many Martyrs (p. \pageref{CommonMartyrs}).
\end{rubric}

%\collect
\lett{A}{lmighty} God, mercifully look upon our infirmities: that whereas we are oppressed by the burden of our sins, the glorious intercession of thy blessed Martyrs Fabian and Sebastian may be our succour and defence. Through.


\festival{Greater Double}{21 January}{St. Agnes of Rome}{V.M.}{MartyrRed}
\feastday{{St. Agnes Rome}}
\feastdate{21}{January}

\begin{rubric}
	Common of Virgins (p. \pageref{CommonVirgin}).
\end{rubric}

\begin{paracol}{2}
	\sloppy
	
\firstEvensong
\psalmAntiphon{When Agnes entered {\dag} the place of shame, she found the Angel of the Lord ready beforehand.}\\

\firstEvensongAntiphon{Blessed Agnes, {\dag} in the midst of the flames, stretched out her hands and prayed: I call on thee, O Father transcendent, august, and dread; for by thy holy Son's protection I have escaped the threats of an impious tyrant, and passed unscathed through the foulness of fleshly pollution: and behold, I come to thee, whom I have loved, whom I have sought, whom I have alway desired.}

\switchcolumn

\secondEvensong
\psalmAntiphon{Rejoice with me {\dag} and be exceeding glad, because among all these I have taken a throne of glory.}\\

℣. Full of grace are thy lips.

℟. Because God hath blessed thee for ever.

\antiphon{Mag.}{While the blessed Agnes {\dag} was standing in the midst of the flames, she stretched out her hands, and prayed unto God: Almighty Lord, worthy of all adoration, fear, and worship: I bless thy holy Name, and I glorify thee for ever and ever.}

\fussy
\end{paracol}

\mattins
\psalmAntiphon{With his ring {\dag} hath my Lord Jesus Christ betrothed me, and as his bride hath he crowned me.}

\mattinsAntiphon{Lo, that which I desired, {\dag} now I see; that for which I hoped, I now possess; I am united in heaven unto him, whom on earth I loved with a perfect devotion.}

\collect
\lett{A}{lmighty} and everlasting God, who dost choose the weak things of the world to confound those things that are strong: mercifully grant; that we who keep the feast of blessed Agnes thy Virgin and Martyress may feel the succour of her intercession in thy sight. Through.


\subbymemorial{22 January}{Sts. Vincent \& Anastasius}{M.}{MartyrRed}
\feastday{{St. Vincent \& Anastasius}}
\feastdate{22}{January}

%New propers needed for addition of Anastasius

\begin{rubric}
	Common of Many Martyrs (p. \pageref{CommonMartyrsNotBishopsCollect}).
\end{rubric}


\subbymemorial{23 January}{St. Emerentiana}{V.M.}{MartyrRed}
\feastday{{St. Emerentiana}}
\feastdate{23}{January}

\begin{rubric}
	Common of Virgins (p. \pageref{VirginMartyressCollect}), second Collect.
\end{rubric}


\festival{Double}{24 January}{St. Timothy}{B.M.}{MartyrRed}
\feastday{{St. Timothy}}
\feastdate{24}{January}

\begin{rubric}
	Common of a Martyr (p. \pageref{CommonMartyrBishopCollect}).
\end{rubric}


\festival{Greater Double}{25 January}{Conversion of St. Paul}{A.}{FestiveWhite}
\feastday{Conversion of St. Paul}
\feastdate{25}{January}

\begin{paracol}{2}[]
\sloppy

\firstEvensong

\psalmAntiphon{I have planted, {\dag} Apollos watered, but God gave the increase, alleluia.}

\firstEvensongHymn\label{PaulEvensong}

\hymnLett{O}{by} thy doctrine, Paul, thou sage illustrious,

\hymnSecondLine{Guide us in virtue, raise our spirits heavenwards;}

\begin{hangparas}{1.25em}{1}

Till perfect knowledge stream on us abundantly,

And that which only is in part be done away.\\

2. Glory eternal to the blessed Trinity,

With laud and honour, virtue and supremacy,

Trinal yet Onely, reigning in his majesty

Both now and ever, through the ages infinite. Amen.\\
\end{hangparas}

℣. Thou art a chosen vessel, holy Apostle Paul.

℟. A preacher of the truth throughout all the world.

\properantiphon{Mag.}{Go forth, Ananias, {\dag} and enquire for Saul, for behold, he prayeth: for he is a chosen vessel unto me, to bear my Name before the Gentiles, and Kings, and the children of Israel.}

\switchcolumn

\mattins

\psalmAntiphon{Most gladly will I glory {\dag} in my infirmities, that the power of Christ may rest upon me.}

\begin{rubric}
	Invitatory Hymn of I Evensong.
\end{rubric}

\begin{rubric}
	Office Hymn of I Evensong from the Common of the Apostles (p. \pageref{ApostlesEvensong}).
\end{rubric}

℣. Thou art a chosen vessel, holy Apostle Paul.

℟. A preacher of the truth throughout all the world.

\properantiphon{Ben.}{Ye which have followed me {\dag} shall sit upon twelve thrones, judging the twelve tribes of Israel, saith the Lord.}

\secondEvensong

\psalmAntiphon{The grace of God {\dag} which was bestowed upon me was not in vain: but his grace abideth with me always.}

\begin{rubric}
	The Office Hymn \& Versicle of I Evensong.
\end{rubric}

\properantiphon{Mag.}{O holy Apostle Paul, {\dag} thou preacher of the truth and Doctor of the Gentiles, intercede for us unto God, who hath chosen thee.}

\fussy
\end{paracol}

\collect
\lett{O}{God,} who, through the preaching of the blessed Apostle Saint Paul, hast caused the light of the Gospel to shine throughout the world; Grant, we beseech thee, that we, having his wonderful conversion in remembrance, may show forth our thankfulness unto thee for the same, by following the holy doctrine which he taught. Through.

\lett{O}{God,} who didst bestow upon thy blessed Apostle Peter the keys of the kingdom of heaven, and didst appoint unto him the high priesthood of binding and loosing: vouchsafe; that by the help of his intercession we may be delivered from the bonds of our iniquities. Who livest.


\subbymemorial{26 January}{St. Polycarp of Smyrna}{B.M.}{MartyrRed}
\feastday{{St. Polycarp Smyrna}}
\feastdate{26}{January}

\memorialCollect
\lett{O}{God,} who makest us glad with the yearly solemnity of blessed Polycarp, thy Martyr and Bishop: mercifully grant; that, as we now cebrate his birthday, so we may likewise rejoice in his protection. Through.


%MANUAL ADJUSTMENT:
\clearpage
\festival{Double (m.t.v.)}{27 January}{St. John Chrysostom}{B.C.D.}{FestiveWhite}
\feastday{{St. John Chrysostom}}
\feastdate{27}{January}

\begin{rubric}
	Common of a Confessor Bishop (p. \pageref{CommonConfessorBishop}).
\end{rubric}

%\collect
\lett{M}{ultiply,} we beseech thee, O Lord, thy Church with thy heavenly grace: even as thou didst vouchsafe to enlighten her with the glorious merits and doctrine of blessed John Chrysostom, thy Confessor and Bishop. Through.


%MANUAL ADJUSTMENT:
\vspace{-0.25\baselineskip}

\festival{Double}{28 January}{St. Cyril of Alexandria}{B.C.D.}{FestiveWhite}
\feastday{{St. Cyril Alexandria}}
\feastdate{28}{January}

\begin{rubric}
	Common of a Confessor Bishop (p. \pageref{CommonConfessorBishop}).
\end{rubric}

%\collect
\lett{O}{God,} who didst make blessed Cyril, thy Confessor and Bishop, an invincible defender of the divine Motherhood of the most blessed Virgin Mary: grant, by his intercession; that we, who believe her to be indeed the Mother of God, may be saved through her maternal protection. Through the same.

\begin{rubric}
    Commemoration \blackrubric{of St. Agnes} (p. \pageref{AgnesCollectII}).
\end{rubric}


%MANUAL ADJUSTMENT:
\vspace{-0.25\baselineskip}

\subbymemorial{28 January}{St. Agnes of Rome}{V.M.}{MartyrRed}
\feastday{{St. Agnes of Rome}}
\feastdate{28}{January}

\memorialCollect\label{AgnesCollectII}
\lett{O}{God,} who makest us glad with the yearly solemnity of blessed Agnes, thy Virgin and Martyress: grant, we beseech thee; that as we venerate her in our service, so we may follow the example of her godly conversation. Through.


%MANUAL ADJUSTMENT:
\vspace{-0.25\baselineskip}

\subbymemorial{30 January}{St. Martina}{V.M.}{MartyrRed}
\feastday{{St. Martina}}
\feastdate{30}{January}

\begin{rubric}
	Common of Virgins (p.\pageref{VirginMartyressCollect}).
\end{rubric}


%MANUAL ADJUSTMENT:
\vspace{-0.25\baselineskip}

\festival{Double}{1 February}{St. Ignatius of Antioch}{B.M.}{MartyrRed}
\feastday{{St. Ignatius Antioch}}
\feastdate{1}{February}

\begin{rubric}
	Common of a Martyr (p. \pageref{CommonMartyr}).
\end{rubric}

%\collect
\lett{A}{lmighty} God, mercifully look upon our infirmities; that whereas we are oppressed by the burden of our sins, the glorious intercession of blessed Ignatius thy Martyr and Bishop may be our succour and defence. Through.

\begin{rubric}
    Commemoration \blackrubric{of St. Bridget of Ireland} (p. \pageref{BridgetCollect}).
\end{rubric}


%MANUAL ADJUSTMENT:
\clearpage
\subbymemorial{1 February}{St. Bridget of Ireland}{V.}{FestiveWhite}
\feastday{{St. Bridget Ireland}}
\feastdate{1}{February}

%Sarum Propers:
%\begin{rubric}
%	The Common of Virgins (p. \pageref{CommonVirginOnlyII}.
%\end{rubric}

\memorialCollect\label{BridgetCollect}
\lett{W}{e} beseech thee, O Lord, that the prayers of Saint Bridget, thy Virgin, being pleasing unto thee, may help us: and never cease to entreat thy loving-kindness towards us. Through.


\festivalLL{Second Class Double}{2 February}{Purification of the Blessed Virgin Mary}
\feastday{Candlemas}
\feastdate{2}{February}

\begin{rubric}
	Common of the Blessed Virgin Mary (p. \pageref{CommonBVM}).
\end{rubric}

\openingSentence{It was revealed unto Simeon by the Holy Ghost, that he should not see death, before he had seen the Lord's Christ.}

\begin{paracol}{2}
\sloppy

\firstEvensong
\psalmAntiphon{When thou was born {\dag} ineffably of a Virgin, then were the Scriptures fulfilled: thou camest down like the rain into a fleece of wool, to bring salvation unto mankind: we praise thee, O our God.}\\

℣. It was revealed unto Simeon by the Holy Ghost. 

℟. That he should not see death before he had seen the Lord's Christ.

\firstEvensongAntiphon{The ancient {\dag} carried the Infant, but the Infant governed the ancient: he whom a Virgin bare, and after bearing, remained virgin, the same was worshipped by her who bare him.}

\switchcolumn

\mattins
\psalmAntiphon{Simeon, {\dag} a just and devout man, awaited the consolation of Israel; and the Holy Ghost was upon him.}

\mattinsAntiphon{And when the parents {\dag} brought in the Child Jesus, then Simeon took him up in his arms, and blessed God, saying: Lord, now lettest thou thy servant depart in peace.}

\secondEvensong
\psalmAntiphon{A light {\dag} to lighten the Gentiles, and the glory of thy people Israel.}

\secondEvensongAntiphon{To-day {\dag} the blessed Virgin Mary presented the Child Jesus in the temple; and Simeon, filled with the Holy Spirit, received him into his arms, and blessed God for ever.}

\fussy	
\end{paracol}

\collect
\lett{A}{lmighty} and everliving God, we humbly beseech thy Majesty, that, as thy only-begotten Son was this day presented in the temple in substance of our flesh, so we may be presented unto thee with pure and clean hearts. By the same thy Son Jesus Christ our Lord. Who liveth.


%MANUAL ADJUSTMENT:
\clearpage
\subbymemorial{3 February}{St. Blaise}{B.M.}{MartyrRed}
\feastday{{St. Blaise}}
\feastdate{3}{February}

\begin{rubric}
	Common of a Martyr (p. \pageref{CommonMartyrBishopCollect}), second Collect.
\end{rubric}


\subbymemorial{4 February}{St. Joseph of Aleppo}{M.}{MartyrRed}
\feastday{{St. Joseph Aleppo}}
\feastdate{4}{February}

\begin{rubric}
	Common of a Martyr (p. \pageref{CommonMartyrNotBishopCollect}).
\end{rubric}

\begin{rubric}
	Commemoration \blackrubric{of the New Martyrs of Russia,} from the Common of Many Martyrs (p. \pageref{CommonMartyrsNotBishopsCollect}).
\end{rubric}


\subbymemorial{4 February}{The New Martyrs of Russia}{M.}{MartyrRed}
\feastday{{New Russian Martyrs}}
\feastdate{4}{February}

\begin{rubric}
	Common of Many Martyrs (p. \pageref{CommonMartyrsNotBishopsCollect}).
\end{rubric}


\festival{Greater Double}{5 February}{St. Agatha}{V.M.}{MartyrRed}
\feastday{{St. Agatha}}
\feastdate{5}{February}

\begin{rubric}
	Common of Virgins (p. \pageref{CommonVirgin}).
\end{rubric}

%\begin{paracol}{2}
%\sloppy

\firstEvensong
\psalmAntiphon{Who art thou that comest unto me {\dag} to heal my wounds? I am the Apostle of Christ; be not doubtful of me, my daughter.}

\antiphon{Mag.}{The blessed Agatha, {\dag} standing in the midst of the prison, with outstretched hands entreated the Lord: O Lord Jesus Christ, my gracious Master, I give thanks unto thee, who hast enabled me to overcome the torments of the executioner: bid me now, O Lord, joyfully to enter into thine unfading glory.}

%\switchcolumn

\mattins
\psalmAntiphon{I give thee thanks, O Lord {\dag} Jesus Christ: because thou art mindful of me, and hast sent thine Apostle to heal my wounds.}

\antiphon{Ben.}{The multitude {\dag} of the heathen, fleeing to the tomb of the virgin, took thence her veil to defend them from the fire: that the Lord might shew himself a deliverer from the burning, for the merits of Agatha his blessed Martyr.}

%\fussy	
%\end{paracol}

\secondEvensong
\psalmAntiphon{On him who hath deigned {\dag} to heal me of all my wounds, and to restore my breast to my bosom; on him, the living God, do I call.}\\

℣. Full of grace are thy lips.

℟. Because God hath blessed thee for ever.

\antiphon{Mag.}{The blessed Agatha, {\dag} standing in the midst of the prison, with outstretched hands entreated the Lord: O Lord Jesus Christ, my gracious Master, I give thanks unto thee, who hast enabled me to overcome the torments of the executioner: bid me now, O Lord, joyfully to enter into thine unfading glory.}

\collect
\lett{O}{God,} who among the manifold works of thy power hast bestowed even upon the weakness of women the victory of martyrdom: mercifully grant; that we, who celebrate the birthday of blessed Agatha, thy Virgin and Martyress, may by her example be drawn nearer unto thee. Through.


\subbymemorial{6 February}{St. Dorothea}{V.M.}{MartyrRed}
\feastday{{St. Dorothea}}
\feastdate{6}{February}

%\begin{rubric}
%	The Common of Virgins (p. \pageref{CommonVirgin}.
%\end{rubric}

\memorialCollect
\lett{W}{e} beseech thee, O Lord, that as blessed Dorothea, thy Virgin and Martyress, was ever found pleasing unto thee, both by the merit of her chastity, and by her confession of thy power: so she may implore for us thy pardon. Through.


\subbymemorial{6 February}{St. Photius}{B.C.}{FestiveWhite}
\feastday{{St. Photius}}
\feastdate{6}{February}

\begin{rubric}
	Common of a Confessor Bishop (p. \pageref{ConfessorBishopCollect}).
\end{rubric}


\festival{Double (m.t.v.)}{7 February}{St. Romuald}{Ab.}{FestiveWhite}
\feastday{{St. Romuald}}
\feastdate{7}{February}

\begin{rubric}
	Common of a Confessor not a Bishop (p. \pageref{AbbotCollect}).
\end{rubric}


\subbymemorial{9 February}{St. Apollonia of Alexandria}{V.M.}{MartyrRed}
\feastday{{\small St. Apollonia Alexandria}}
\feastdate{9}{February}

%\begin{rubric}
%	The Common of Virgins (p. \pageref{CommonVirgin}.
%\end{rubric}

\memorialCollect
\lett{O}{God} who among the manifold works of thy power hast bestowed even upon the weakness of women the victory of martyrdom mercifully grant; that we, who celebrate the birthday of blessed Apollonia thy Virgin and Martyress, may by her example be drawn nearer unto thee. Through.


%MANUAL ADJUSTMENT:
\clearpage
\festival{Greater Double}{10 February}{St. Scholastica}{V.}{FestiveWhite}
\feastday{{St. Scholastica}}
\feastdate{10}{February}

\begin{paracol}{2}[]
\sloppy

%MANUAL ADJUSTMENT:
\vspace{-1\baselineskip}

\firstEvensong
\psalmAntiphon{Lo, I entreated thee, {\dag} and thou wouldest not hear me: I entreated my Lord, and he heard me.}

\firstEvensongHymn\label{ScholasticaEvensong}

\hymnLett{B}{lessed} bride of Christ the Bridegroom, Dove of virgins, holy maid,

\hymnSecondLine{Starry hosts proclaim thy praises, Praise the crown to virtue paid,}

\begin{hangparas}{1.25em}{1}

While our joyful hearts and voices Join the anthem unafraid.\\

2. Wise wert thou to spurn the riches And the crowns that heaven miss,

Wise to choose thy brother's teachings, Change thy way of life for his:

In the savour of thine ointments Didst thou enter heav'nly bliss.\\

3. Oh, the might of love unbounded! Victory of all most meet!

At thy tears the heavens opened, Poured their floods about thy feet,

While his soul to thine made answer All night long in converse sweet.\\

4. Crash of tempest, roll of thunder, Lightning flash from pole to pole,

But the storm of love is stronger, Brighter flashes in the soul,

While the peace of Christ the Bridegroom Holds thee still and keeps thee whole.\\

5. Now a gleaming cloud in heaven, Now a sun with golden rays,

Never darkling, never fading, Shining still through all our days,

Fill the hearts of all the faithful With the joy of heav'nly lays.\\

6. Glory to the Father sing we, Glory to the only Son;

To the Paraclete in glory, Equal tribute be begun,

At whose pleasure made and govern'd All the ages' course is run. Amen.\\
\end{hangparas}

℣. Who is this that flieth as a cloud? 

℟. And as a dove to her windows?

\antiphon{Mag.}{Let all the multitude {\dag} of the faithful exult in the glory of the gracious virgin Scholastica: and chiefly let the company of virgins be joyful, celebrating her Solemnity; for she besought the Lord, pouring forth her tears, and of him received greater power, because her love was greater.}

\secondEvensong
\psalmAntiphon{When holy Benedict {\dag} had withdrawn himself in his cell for three days, he lifted up his eyes, and saw the soul of his sister, departing from the body in the likeness of a dove, enter the secret places of the heavens.}\\

%MANUAL ADJUSTMENT:
\vspace{-1\baselineskip}

\begin{rubric}
	Office Hymn \& Versicle of I Evensong.
\end{rubric}

\antiphon{Mag.}{To-day {\dag} the holy virgin Scholastica, in the likeness of a dove, went forth with all gladness to the heavenly places: to-day she was found worthy to enjoy for ever the bliss of celestial life beside her brother.}

\collect
\lett{O}{God,} who didst reveal in a vision the soul of blessed Scholastica thy Virgin entering heaven in the likeness of a dove, that thou mightest shew the way of the undefiled: grant us by the aid of her merits and prayers so innocently to live, that we may worthily attain unto joys eternal. Through.

\switchcolumn

%MANUAL ADJUSTMENT:
\vspace{-1\baselineskip}

\mattins
\psalmAntiphon{Now speak we {\dag} even until morning of heavenly things, in holy converse regarding the life of the spirit.}

\mattinsInvitatoryHymn\label{ScholasticaInvitatory}

\hymnLett{N}{ow} we our voices raise in sweet angelic hymn,

\hymnSecondLine{While earthly hills and vales in silent night grow dim:}

\begin{hangparas}{1.25em}{1}

Scholastica calls forth a heav'nly melody

From souls of virgin purity.\\

2. She, sprung from Nursian race, her stem nobility,

As bride to Virgin Lamb more noble yet shall be;

The fragrance of the Spouse she ever doth impart,

The beauty of his wounded Heart.\\

3. At their fraternal meal besought she zealously:

For greater nourishment hunger'd her charity.

And so with still more joy her brother truth outpour'd

While she more fervently implor'd.\\

4. O peaceful hours of night, O time of quiet rest,

When heav'nly feasting so inebriates the breast,

While words of loving zeal from one another flow;

Thy will, O Jesus, they would know.\\

5. O very Rest of heart, O wondrous Trinity

Whose glance doth satiate with heav'nly radiancy,

May talk of thee be sweet, more sweet thee to attain,

Sweetest eternally to gain. Amen.
\end{hangparas}

\mattinsOfficeHymn\label{ScholasticaMattins}

\hymnLett{T}{he} shades of night are fled away,

\hymnSecondLine{The long desired day is born;}

\begin{hangparas}{1.25em}{1}

The changing wheel of time has brought

Scholastica her wedding morn.\\

2. The winter's weariness is past,

Tempests no longer doth she see;

The fields of heaven bloom with stars,

The blossoms of eternity.\\

3. He calls her forth, the Source of love,

He gives her wings that she may rise,

And to the kisses of his mouth,

A shining dove she swiftly flies.\\

4. O royal child, how fair thy course!

O maid, how beautiful thy way!

Thy brother marks thy upward flight,

Then turns to God his thanks to pay.\\

5. Lock'd in the Bridegroom's close embrace,

She takes the crown to virtue owed,

Sunk in a glorious stream of bliss,

Inebriated with her God.\\

6. Thou lily of the valley, Christ,

To thee our homage meet we pay:

To Father and to Paraclete,

While endless ages roll away. Amen.\\
\end{hangparas}

℣. Behold, thou art fair, my love.  

℟. Behold, thou art fair; thou hast dove's eyes.

\antiphon{Ben.}{O how illustrious {\dag} are the merits of blessed Scholastica! O how great the power of her tears! through which the renowned virgin, out of sunny clearness, drew down from the air a mighty flood of rain.}

\end{paracol}


\subbymemorial{11 February}{Pope St. Gregory II}{B.C.}{FestiveWhite}
\feastday{{Pope St. Gregory II}}
\feastdate{11}{February}

\begin{rubric}
	Common of a Confessor Bishop (p. \pageref{ConfessorBishopCollect}).
\end{rubric}


\subbymemorial{14 February}{St. Valentine}{M.}{MartyrRed}
\feastday{{St. Valentine}}
\feastdate{14}{February}

%\begin{rubric}
%	The Common of a Martyr (p. \pageref{CommonMartyr}).
%\end{rubric}

\memorialCollect
\lett{G}{rant,} we beseech thee, Almighty God: that we who observe the birthday of blessed Valentine thy Martyr may by his intercession be delivered from all evils that beset us. Through.


\subbymemorial{15 February}{Sts. Faustinus and Jovita}{M.}{MartyrRed}
\feastday{{Sts. Faustinus \& Jovita}}
\feastdate{15}{February}

%\begin{rubric}
%	The Common of Many Martyrs (p. \pageref{CommonMartyrs}).
%\end{rubric}

\memorialCollect
\lett{O}{God,} who makest us glad with the yearly solemnity of thy holy Martyrs Faustinus and Jovita: mercifully grant; that as we rejoice in their merits, so we may be enkindled by their example. Through.


\subbymemorial{18 February}{St. Simeon of Jerusalem}{B.M.}{MartyrRed}
\feastday{{St. Simeon Jerusalem}}
\feastdate{18}{February}

\begin{rubric}
	Common of a Martyr (p. \pageref{CommonMartyrBishopCollect}).
\end{rubric}


%MANUAL ADJUSTMENT:
\clearpage
\festival{Second Class Double}{22 February}{Chair of St. Peter at Antioch}{A.}{FestiveWhite}\label{CathedraAntioch}
\feastday{Chair St. Peter Antioch}
\feastdate{22}{February}


\begin{paracol}{2}
\sloppy

\firstEvensong\label{PeterAntiochEvensong}
\psalmAntiphon{Behold a great priest, {\dag} who in his days pleased God, and was found righteous.}

\firstEvensongHymn
\selectlanguage{ecclesiasticallatin}

%I fear I did a bad job translating, so I'm just supplying the Latin here.
\hymnLett{Q}{uodc\smash{ú}mque} vinclis super terram strínxeris,

\hymnSecondLine{Erit in astris religátum fórtiter:}
\begin{hangparas}{1.25em}{1}
Et quod resólvis in terris arbítrio,

Erit solútum super cæli rádium:

In fine mundi judex eris sǽculi.\\

2. Gloria Patriper imménsa sǽcula;

Sit tibi, Nate, decus et impérium;

Honor, potéstas, Sanctóque Spirítui:

Sit Trinitáti salus individua,

Per infiníta sæculórum sǽcula. Amen.\\
\end{hangparas}

\selectlanguage{british}
    
    ℣. Thou art Peter.

	℟. And upon this rock I will build my Church.
	
\firstEvensongAntiphon{Thou art the shepherd of the sheep, {\dag} Prince of the Apostles, unto thee were given the keys of the kingdom of heaven.}

\switchcolumn

\mattins
\psalmAntiphon{The Lord, {\dag} therefore, assured him by an oath that he would bless the nations in his seed.}

\begin{rubric}
	Invitatory Hymn of I Evensong.
\end{rubric}

\mattinsOfficeHymn\label{PeterAntiochMattins}

\hymnLett{P}{eter,} good shepherd, may thy ceaseless orisons,

\hymnSecondLine{For us prevailing, break the bands of wickedness:}
\begin{hangparas}{1.25em}{1}
For thou of old time didst receive authority

The gates to open, or to close, of Paradise.\\

2. Glory eternal to the blessed Trinity,

With laud and honour, virtue and supremacy,

Trinal yet Onely, reigning in his majesty

Both now and ever, through the ages infinite. Amen.\\
\end{hangparas}

    ℣. Thou art Peter.

	℟. And upon this rock I will build my Church.
	
	\properantiphon{Ben.}{Whatsoever {\dag} thou shalt bind on earth shall be bound in heaven: and whatsoever thou shalt loose on earth shall be loosed in heaven, saith the Lord unto Simon Peter.}

\fussy	
\end{paracol}

\secondEvensong
\psalmAntiphon{Good and faithful servant, {\dag} enter thou into the joy of thy Lord.}

\begin{rubric}
	Office Hymn \& Versicle of I Evensong.
\end{rubric}

\properantiphon{Mag.}{Whilst he was chief Bishop, he feared nothing on earth, but ascended gloriously to the heavenly kingdoms.}


\collect
\lett{O}{God,} who didst bestow upon thy blessed Apostle Peter the keys of the kingdom of heaven, and didst appoint unto him the high priesthood of binding and loosing: vouchsafe; that by the help of his intercession we may be delivered from the bonds of our iniquities. Who livest.

\lett{O}{God,} who by the preaching of the blessed Apostle Paul didst teach the multitude of the Gentiles: grant to us, we beseech thee; that we who celebrate his commemoration may know him to be our advocate with thee. (Through.)

\begin{rubric}
    If this Feast fall on a Saturday---not in a Leap Year---Commemoration is made \blackrubric{of the Vigil of St. Matthias,} from the Common of Vigils of the Apostles (p. \pageref{CommonVigilApostles}).
\end{rubric}

\begin{rubric}
    In Lent, Commemoration of the Feria.
\end{rubric}


\festival{Vigil}{23/24 February}{Vigil of St. Matthias}{A.}{LentenViolet}\label{MatthiasVigil}
\feastday{{St. Matthias Vigil}}
\feastdate{23/24}{February}

\begin{rubric}
	Common of Vigils of the Apostles (p. \pageref{CommonVigilApostles}).
\end{rubric}


\festival{Second Class Double}{24/25 February}{St. Matthias}{A.}{FestiveWhite}
\feastday{St. Matthias}
\feastdate{24/25}{February}

\begin{rubric}
	Common of Apostles (p. \pageref{CommonApostles}).
\end{rubric}

\lett{O}{God,} who didst join blessed Matthias to the college of thine Apostles: grant, we beseech thee; that through his mediation we may ever perceive thy tender mercy towards us. Through.
\begin{rubric}
    In Lent, Commemoration of the Feria.
\end{rubric}


%MANUAL ADJUSTMENT:
\clearpage
\subbymemorial{25/26 February}{St. Walburga of Heidenheim}{V.}{FestiveWhite}
\feastday{{St. Walburga Heidenheim}}
\feastdate{25/26}{February}

%\begin{rubric}
%	The Common of Virgins (p. \pageref{CommonVirgin}).
%\end{rubric}

\memorialCollect
\lett{O}{God,} who among the manifold gifts of thy grace dost also work great wonders even in the weaker sex: mercifully grant that we may feel the help of the intercession of blessed Walburga, thy Virgin, whose example of chastity doth enlighten us, and whose glory in miracles doth gladden us. Through.


\subbymemorial{26/27 February}{Pope St. Alexander of Alexandria}{B.C.}{FestiveWhite}
\feastday{{Pope St. Alexander}}
\feastdate{26/27}{February}

\begin{rubric}
	Common of a Confessor Bishop (p. \pageref{ConfessorBishopCollect}).
\end{rubric}


\festival{Greater Double}{27/28 February}{St. Raphael of Brooklyn}{B.C.}{FestiveWhite}
\feastday{{St. Raphael Brooklyn}}
\feastdate{27/28}{February}

\begin{rubric}
	In Lent, Commemoration only.
\end{rubric}

\begin{rubric}
	Common of a Confessor Bishop (p. \pageref{CommonConfessorBishop}).
\end{rubric}

%MANUAL ADJUSTMENT:
\lett{A}{lmighty}\textls[-15]{ God, who hast instructed thy Church with the heavenly doctrine of thy Bishop and Confessor Saint Raphael of Brooklyn; and hast provided in his life an example of the Good Shepherd who collects the lost sheep; grant us so to follow in his godly footsteps, that, ever helped by his intercession, we with him may attain to thy heavenly kingdom. Through the same.}


\subbymemorial{1 March}{St. David of Wales}{B.C.}{FestiveWhite}
\feastday{{St. David Wales}}
\feastdate{1}{March}

%\begin{rubric}
%	The Common of a Confessor Bishop (p. \pageref{CommonConfessorBishop}).
%\end{rubric}

\memorialCollect
\lett{G}{rant} to us, Almighty God: that the loving intercession of blessed David, thy Confessor and Bishop, may protect us; that while we celebrate his festival we may imitate his steadfastness in the defence of the Catholic faith. Through.


\subbymemorial{2 March}{St. Chad}{B.C.}{FestiveWhite}
\feastday{{St. Chad}}
\feastdate{2}{March}

%\begin{rubric}
%	The Common of a Confessor Bishop (p. \pageref{CommonConfessorBishop}).
%\end{rubric}

\memorialCollect
\lett{A}{lmighty} and everlasting God, who on this day dost gladden us by the festival of blessed Chad, thy Confessor and Bishop: we humbly beseech thy mercy; that we, who devoutly observe and venerate his festival, may by his loving advocacy obtain the reward of everlasting life. Through.


%MANUAL ADJUSTMENT:
\clearpage
\subbymemorial{4 March}{Pope St. Lucius I of Rome}{B.M.}{MartyrRed}
\feastday{{Pope St. Lucius I}}
\feastdate{4}{March}

%\begin{rubric}
%	The Common of a Martyr (p. \pageref{CommonMartyr}).
%\end{rubric}

\memorialCollect
\lett{O}{God,} who makest us glad with the yearly solemnity of blessed Lucius thy Martyr and Bishop: mercifully grant; that, as we now celebrate his birthday, so we may likewise rejoice in his protection. Through.


\festival{Double}{6 March}{Sts. Perpetua and Felicitas}{M.}{MartyrRed}
\feastday{{Sts. Perpetua \& Felicitas}}
\feastdate{6}{March}

\begin{rubric}
	In Lent, Commemoration only.
\end{rubric}

\begin{rubric}
	Common of Holy Matrons (p. \pageref{CommonHolyMatrons}).
\end{rubric}

\lett{G}{rant,} we beseech thee, O Lord our God, that we may at all times so devoutly honour the triumphs of thy holy Martyrs Perpetua and Felicitas: that, although we cannot worthily shew forth their praises, yet we may continually honour them with lowly service. Through.


\subbymemorial{9 March}{St. Gregory of Nyssa}{B.C.D.}{FestiveWhite}
\feastday{{St. Gregory Nyssa}}
\feastdate{9}{March}

\begin{rubric}
	Common of a Confessor Bishop (p. \pageref{ConfessorBishopDoctorCollect}).
\end{rubric}


\subbymemorial{10 March}{The Forty Holy Martyrs}{M.}{MartyrRed}
\feastday{{40 Holy Martyrs}}
\feastdate{10}{March}

\memorialCollect
\lett{G}{rant} we beseech thee, Almighty God: that, like as we have known thy glorious Martyrs to be constant in their confession, so we may perceive their loving intercession for us with thee. Through.


\festival{Second Class Double}{12 March}{Pope St. Gregory the Great}{B.C.D.}{FestiveWhite}
\feastday{Pope St. Gregory}
\feastdate{12}{March}

\begin{paracol}{2}
\sloppy

%MANUAL ADJUSTMENT:
\vspace{-1\baselineskip}

\firstEvensong\label{GregoryEvensong}
\psalmAntiphon{The blessed Gregory, {\dag} lifted up to the throne of Peter, shewed forth in deed his name of `the Watchful.'}

\firstEvensongHymn

\hymnLett{A}{postle} of the English, thou

\hymnSecondLine{Art comrade of the Angels now:}
\begin{hangparas}{1.25em}{1}
O help the faithful, Gregory,

Who heardest then the peoples' cry.\\

2. Abundant wealth was naught to thee

Who, seeking princely penury,

Didst spurn the glory of the earth

To follow Jesus in his dearth.\\

3. Like some poor needy shipwreck'd soul,

His messenger besought a dole;

Twin gifts thou gavest; nothing loth

To offer self and silver both.\\

4. First in his Church from that time forth

Christ set thee, in thy proven worth;

So didst thou rise to Peter's throne,

Whose life was pattern for thine own.\\

5. O Priest and Leader of the flock,

The Church's glory, light, and rock,

Instructed by thy wise command,

Let not the sheep in peril stand.\\

6. Praise to the unbegotten One,

And glory to his only Son;

And thine, O Spirit, Breath of God,

Be equal majesty and laud. Amen.\\
\end{hangparas}

    ℣. The Lord loved him and adorned him.

	℟. He clothed him with a robe of glory.

\firstEvensongAntiphon{O Teacher right excellent, {\dag} O light of Holy Church, O blessed Gregory, lover of the divine law: intercede for us unto the Son of God.}

\switchcolumn

%MANUAL ADJUSTMENT:
\vspace{-1\baselineskip}

\mattins\label{GregoryMattins}
\psalmAntiphon{While he revealed {\dag} the mysteries of Holy Writ, there appeared a dove, whiter than snow.}

\begin{rubric}
	Invitatory Hymn of I Evensong.
\end{rubric}

\mattinsOfficeHymn

\hymnLett{T}{hy} lips, O Gregory, distil

\hymnSecondLine{The honey that thy heart doth fill;}
\begin{hangparas}{1.25em}{1}

In spiced sweetness floweth thence

The savour of thine eloquence.\\

2. The hidden mysteries that lie

In Holy Scripture, to thine eye

Are plain; the Truth himself draws near

To make their subtle secrets clear.\\

3. At once the Apostolic throne

And majesty become thine own:

O free us from sin's binding chain

That heav'nly mansions we may gain.\\

4. O Priest and Leader of the flock,

The Church's glory, light, and rock,

Instructed by thy wise command,

Let not the sheep in peril stand.\\

5. Praise to the unbegotten One,

And glory to his only Son;

And thine, O Spirit, Breath of God,

Be equal majesty and laud. Amen.\\
\end{hangparas}

    ℣. The Lord guided the righteous in right paths.

	℟. And shewed him the kingdom of God.

\mattinsAntiphon{Gregory, {\dag} when he looked upon the youthful Angles, said: They have the countenance of Angels, and such as these should be of the fellowship of the Angels in heaven.}

\secondEvensong

\begin{rubric}
	Office Hymn \& Antiphon of I Evensong. Versicle of Mattins.
\end{rubric}

\fussy	
\end{paracol}


\collect
\lett{O}{God,} who on the soul of thy servant Gregory, hast bestowed the rewards of everlasting bliss: mercifully grant; that we, who are oppressed by the burden of our sins, may by his prayers before thee be relieved. Through.


\festival{Double}{17 March}{St. Patrick}{B.C.}{FestiveWhite}
\feastday{{St. Patrick}}
\feastdate{17}{March}

\begin{rubric}
	In Lent, Commemoration only.
\end{rubric}

\begin{rubric}
	Common of a Confessor Bishop (p. \pageref{CommonConfessorBishop}).
\end{rubric}

\lett{O}{God,} who for the preaching of thy glory unto the Gentiles wast pleased to send forth blessed Patrick, thy Confessor and Bishop: grant by his merits and intercession; that we may through thy mercy be enabled to accomplish those things which thou commandest us to do. Through.

\begin{rubric}
    Commemoration \blackrubric{of St. Joseph of Arimathea} (p. \pageref{ArimatheanCollect}).
\end{rubric}


\subbymemorial{17 March}{St. Joseph of Arimathea}{C.}{FestiveWhite}
\feastday{{St. Joseph Arimathea}}
\feastdate{17}{March}

%\begin{rubric}
%	The Common of a Confessor not a Bishop (p. \pageref{CommonConfessorNotBishop}).
%\end{rubric}

\memorialCollect\label{ArimatheanCollect}
\lett{O}{God,} who didst give such grace unto thy servant Joseph that he boldly craved the body of Jesus, and with great reverence laid him in the rock-hewn sepulchre: grant, we beseech thee, that we likewise may be so emboldened for thee, as to do works meet for thy Kingdom. Through the same.


\festival{Double}{18 March}{St. Cyril of Jerusalem}{B.C.D.}{FestiveWhite}
\feastday{{St. Cyril Jerusalem}}
\feastdate{18}{March}

\begin{rubric}
	In Lent, Commemoration only.
\end{rubric}

\begin{rubric}
	Common of a Confessor Bishop (p. \pageref{CommonConfessorBishop}).
\end{rubric}

\label{CyrilCollect}
\lett{G}{rant} to us, we beseech thee, Almighty God, at the intercession of the blessed Bishop Cyril: so to know thee, the only true God, and Jesus Christ whom thou hast sent; that we may be found worthy to be numbered for evermore among the sheep who hear his voice. Through the same.

\begin{rubric}
	Commemoration \blackrubric{of King St. Edward} (p. \pageref{EdwardCollect}).
\end{rubric}


%MANUAL ADJUSTMENT:
\clearpage
\subbymemorial{18 March}{King St. Edward}{M.}{MartyrRed}
\feastday{{King St. Edward}}
\feastdate{18}{March}

%\begin{rubric}
%	The Common of a Martyr (p. \pageref{CommonMartyr}).
%\end{rubric}

\memorialCollect\label{EdwardCollect}
\lett{O}{God,} the triumphant ruler of an everlasting kingdom, mercifully behold this thy family who celebrate the memory of blessed Edward, thy King and Martyr; and, by his merits and intercession, vouchsafe; that they, who glory in his triumph, may also attain unto his rewards. Through.

%\begin{rubric}
%	Commemoration \blackrubric{of St. Cyril} (p. \pageref{CyrilCollect}).
%\end{rubric}


\festival{Second Class Double}{19 March}{St. Joseph, Spouse of the Blessed Virgin Mary}{C.}{FestiveWhite}
\feastday{St. Joseph}
\feastdate{19}{March}

\begin{paracol}{2}
	


\firstEvensong
\psalmAntiphon{And Jacob {\dag} begat Joseph, the husband of Mary, of whom was born Jesus, who is called Christ.}

\firstEvensongHymn\label{JosephEvensong}

\hymnLett{O}{Joseph,} heav'nly hosts thy worthiness proclaim,

\hymnSecondLine{And Christendom conspires to celebrate thy fame,}
\begin{hangparas}{1.25em}{1}
	
Thou who in purest bonds wert to the Virgin bound;

How glorious is thy name renown'd.\\

2. Thou, when thou didst behold thy Spouse about to bear,

Wert sore oppress'd with doubt, wert fill'd with wondering care;

At length the Angel's word thy anxious heart reliev'd:

She by the Spirit hath conceiv'd.\\

3. Thou with thy new-born Lord didst seek far Egypt's land,

As wand'ring pilgrims, ye fled o'er the desert sand;

That Lord, when lost, by thee is in the temple found,

While tears are shed, and joys abound.\\

4. Not till death's hour is past do other men obtain

The meed of holiness, and glorious rest attain;

Thou, like to Angels made, in life completely blest,

Dost clasp thy God unto thy breast.\\

5. O Holy Trinity, thy suppliant servants spare;

Grant us to ride to heaven, for Joseph's sake and prayer,

And so our grateful hearts to thee shall ever raise

Exulting canticles of praise. Amen.\\
\end{hangparas}

    ℣. He made him lord of his house.

	℟. And ruler of all his substance.
	
	\firstEvensongAntiphon{Then Joseph, {\dag} being raised from sleep, did as the Angel of the Lord had bidden him, and took unto him his wife.}

\secondEvensong
\psalmAntiphon{Jesus went down with them, {\dag} and came to Nazareth, and was subject unto them.}

\begin{rubric}
	Office Hymn of I Evensong.
\end{rubric}

    ℣. Riches and plenteousness shall be in his house.

	℟. And his righteousness endureth for ever.
	
	\secondEvensongAntiphon{Behold a faithful {\dag} and wise servant, whom his Lord hath made ruler over his household.}
	
\collect
\lett{W}{e} beseech thee, O Lord, that we may be aided through the merits of the Spouse of thy most holy Mother: that those things, which by our own power we cannot obtain, may through his intercession be granted unto us. Who livest.

\begin{rubric}
    In Lent, Commemoration of the Feria.
\end{rubric}
	
	\switchcolumn
	
\mattins
\psalmAntiphon{When Mary {\dag} the mother of Jesus was espoused to Joseph, before they came together, she was found with child of the Holy Ghost.}
	
	\mattinsInvitatoryHymn\label{JosephInvitatory}
	
\hymnLett{J}{oseph,} the praise and glory of the heavens,

\hymnSecondLine{Sure pledge of life, and safety of the wide world,}
\begin{hangparas}{1.25em}{1}

As in our joy we sing to thee, in kindness

List to our praises.\\

2. Thou by the world’s Creator wert appointed

Spouse of the Virgin: thee he willed to honour

Naming thee Father of the Word, and guardian

Of our salvation.\\

3. Thou thy Redeemer, lying in a stable,

Whom long ago foretold the choir of prophets,

Sawest rejoicing, and thy God adoredst

Humble in childhood.\\

4. God, King of Kings, and Governor of the ages,

He at whose word the pow'rs of hell do tremble,

He whom the adoring heav'ns ever worship

Called thee protector.\\

5. Praise to the Triune Godhead everlasting,

Who with such honour mightily hath blest thee;

O may he grant us at thy blest petition

Joys everlasting. Amen.
\end{hangparas}
	

\mattinsOfficeHymn\label{JosephMattins}
	
\hymnLett{H}{e,} whom the faithful joyously do honour,

\hymnSecondLine{Singing his praises with devout affection,}
\begin{hangparas}{1.25em}{1}

Won on this feast day, in eternal glory

Life everlasting.\\

2. Blest beyond others, and exceeding blissful,

For, when the moment of his death was nearing,

Jesus and Mary at his side were standing,

Soothing his spirit.\\

3. Death doth he conquer, laying down his burden,

Calmly he slumbers, rest he gains eternal;

Lo, round his forehead, bright with rays of splendour,

Shineth a garland.\\

4. Then, as he reigneth, earnestly beseech we

That he may utter fervent intercessions,

Praying that pardon and the peace of heaven

May be our portion.\\

5. Glory we give thee, hymn of praise and blessing,

One in Three Persons, who above art reigning,

God, who hast honour'd with thy crown for ever

This thy true servant. Amen.\\
\end{hangparas}

    ℣. The mouth of the righteous is exercised in wisdom.

	℟. And his tongue will be talking of judgement.

\mattinsAntiphon{Jesus himself {\dag} began to be about thirty years of age, being, as was supposed, the son of Joseph.}


\end{paracol}



\festival{Double}{20 March}{St. Cuthbert}{B.C.}{FestiveWhite}
\feastday{{St. Cuthbert}}
\feastdate{20}{March}

\begin{rubric}
	In Lent, Commemoration only.
\end{rubric}

\begin{rubric}
	Common of a Confessor Bishop (p. \pageref{CommonConfessorBishop}).
\end{rubric}

%\collect
\lett{O}{God,} who dost make thy Saints glorious by the inestimable gift of thy grace: grant, we beseech thee; that at the intercession of blessed Cuthbert, thy Confessor and Bishop, we may be found worthy to attain to the perfection of all virtue. Through.


%MANUAL ADJUSTMENT:
\clearpage
\festival{Second Class Double}{21 March}{St. Benedict}{Ab.}{FestiveWhite}
\feastday{St. Benedict}
\feastdate{21}{March}

\begin{paracol}{2}
	\sloppy
	
%MANUAL ADJUSTMENT:
\vspace{-1\baselineskip}
	
\firstEvensong\label{BenedictEvensong}
\psalmAntiphon{There was a man of venerable life, {\dag} Benedict by grace and by name, who, from his younger years, carried always the mind of an old man; overpassing by virtue, he denied himself all vain delights.}

\firstEvensongHymn

\hymnLett{S}{hout,} all ye people! Let your measur'd praises

\hymnSecondLine{Ring through the churches solemnly and sweetly:}
\begin{hangparas}{1.25em}{1}
	
On this his feast day, Benedict ascended

Heaven's high summit.\\

2. He, when his youthful joyous years were blooming,

Yet in his boyhood, left his native dwelling,

Seeking concealment hid within a cavern

Lonely and silent.\\

3. There amid nettles, rigid thorns, and briars

Won he the battle over youth's enticement,

Nurse of pollution: then he wrote a Holy

Rule of blest living.\\

4. Thy brazen image, infamous Apollo,

Soon hath he smitten; burnt the grove of Venus;

Then to the Baptist, on the sacred mountain

Stablish'd a chapel.\\

5. Now doth he witness happily in heaven

Seraphim, leading throngs of shining Angels,

While he refreshes faithful hearts who need him

With living waters.\\

6. Praise to the Father, to the Sole-begotten,

And to thee, alway with the Twain co-equal,

Fostering Spirit; one and only Godhead

Through all the ages. Amen.\\
\end{hangparas}

    ℣. The Lord loved him and adorned him.

	℟. He clothed him with a robe of glory.
	
	\firstEvensongAntiphon{Let all the multitude {\dag} of the faithful exult in the glory of the gracious Father Benedict: and chiefly let the company of monks be joyful, celebrating his festival on earth, in whose goodly fellowship the Saints rejoice in heaven.}

\secondEvensong

\begin{rubric}
	Office Hymn of I Evensong. Versicle of Mattins.
\end{rubric}

\secondEvensongAntiphon{To-day {\dag} holy Benedict, while his disciples beheld it, ascending by way of the East on a straight pathway to heaven: to-day with hands uplifted, he died between the words of his supplication: to-day he was received into glory by the Angels.}

\collect
\lett{A}{lmighty} and everlasting God, who on this day didst release thy most blessed Confessor Benedict from the bondage of the flesh and take him up to heaven, grant, we beseech thee, to us thy servants who celebrate this feast pardon of all our sins, that we who with joyful hearts take pleasure in his renown may for his sake, and at his intercession, have fellowship with him. Through.
\begin{rubric}
    In Lent, Commemoration of the Feria.
\end{rubric}

\switchcolumn
	
%MANUAL ADJUSTMENT:
\vspace{-1\baselineskip}
	
\mattins
\psalmAntiphon{The blessed man Benedict {\dag} desired rather the miseries of the world than perpetual praises; to be wearied with labour for God’s sake, than to be exalted with transitory favour.}
	
	\mattinsInvitatoryHymn\label{BenedictInvitatory}
	
\hymnLett{W}{hate\smash{'}er} in former days befell the prophets,

\hymnSecondLine{Whate'er renowned in the law eternal:}
\begin{hangparas}{1.25em}{1}

In him, the greatest monk, whose feast we honour,

All is contained.\\

2. Virtue of Moses issued forth in mercy;

Wondrous the offspring unto Abram given;

Righteous and stern, the ruling of his parent

Bride gained for Isaac.\\

3. Benedict laden loftily with virtues:

Higher were those our Patriarch had gather'd,

Merits of Isaac, Abraham and Moses

One breast containeth.\\

4. Those by the world’s blows overcome he lifted,

Meekness establish'd in the place of anger,

Peace where it was not, rest he made to well up

From midst of terror.\\

5. Praise to the Father, to the Sole-begotten,

And to thee, alway with the Twain co-equal,

Fostering Spirit; one and only Godhead

Through all the ages. Amen.
\end{hangparas}

	\mattinsOfficeHymn\label{BenedictMattins}

\hymnLett{G}{em} of the highest, diadem immortal,

\hymnSecondLine{Cherish'd for ever after sacred struggle,}
\begin{hangparas}{1.25em}{1}

Thou, among many, first in lofty merit,

Benedict, shinest.\\

2. Thee in thy boyhood holy eld adorned;

Surely of nothing hath desire despoil'd thee;

Earth's blossom wither'd, to thy heart uplifted

Unto the highest.\\

3. Fleeing in secret, thou didst leave thy native

Country and parents: then, a forest-dweller,

Into subjection unto Christ, by torment

Broughtest thy body.\\

3. Not long securely couldest thou be hidden,

Good deeds and wonders brought thee to men's knowledge;

Through the world widely spread the happy rumour

Speedily flying.\\

4. Praise to the Father, to the Sole-begotten,

And to thee, alway with the Twain co-equal,

Fostering Spirit; one and only Godhead

Through all the ages. Amen.\\
\end{hangparas}

    ℣. The Lord guided the righteous in right paths.

	℟. And shewed him the kingdom of God.
	
	\mattinsAntiphon{Benedict, {\dag} thou father and guide of monks, thou most holy Confessor of the Lord, intercede for us all and for our salvation.}


	\fussy
\end{paracol}



\festivalAO{Greater Double}{24 March}{St. Gabriel the Archangel}
\feastday{St. Gabriel Archangel}
\feastdate{24}{March}

\begin{rubric}
	Commemoration only is made, except where he be the patron.
\end{rubric}

%\begin{rubric}
%	I Evensong: Psalms 9, 11, \& 15. Daniel 8:15-19. Revelation 8:1-5.\par
%	Mattins: Psalms 30, 34, \& 77. Daniel 9:20-26. Luke 1:5-17\par
%	II Evensong: Psalms 97, 99, \& 103. Isaiah 7:10-14. Luke 1:26-38.
%\end{rubric}
\par\noindent
%AOB:
\textit{Opening Sentence.} I saw another angel come down from heaven, having great power; and the earth was lightened with his glory.

\begin{paracol}{2}[]
\sloppy

\firstEvensong
\psalmAntiphon{When Zacharias {\dag} went into the temple, there appeared unto him the Archangel Gabriel, standing on the right side of the altar of incense.}

\firstEvensongHymn

\hymnLett{C}{hrist,} the fair glory of the holy Angels,

\hymnSecondLine{Thou who hast made us, thou who o'er us rulest,}
\begin{hangparas}{1.25em}{1}

Grant of thy mercy unto us thy servants

Steps up to heaven.\\

2. Send thy Archangel, Gabriel, the mighty,

Herald of heaven; may he from us mortals

Spurn the old serpent, watching o'er the temples

Where thou art worshipp'd.\\

3. May the blest Mother of our God and Saviour,

May the assembly of the Saints in glory,

May the celestial companies of Angels

Ever assist us.\\

4. This he vouchsafe us, God for ever blessed

Father eternal, Son, and Holy Spirit,

Whose is the glory which through all creation.

Ever resoundeth. Amen.\\
\end{hangparas}

    ℣. An Angel stood at the altar of the temple.

	℟. Having in his hand a golden censer.
	
	\firstEvensongAntiphon{The Angel Gabriel {\dag} came in unto Mary and said, Hail, thou that art full of grace, the Lord is with thee; blessed art thou among women.}
	
	\switchcolumn
	
\mattins
\psalmAntiphon{I am the Angel Gabriel, {\dag} that stands in the presence of God; and am sent to speak unto thee.}

\begin{rubric}
	Invitatory Hymn of I Evensong.
\end{rubric}

\mattinsOfficeHymn

\hymnLett{O}{Christ,} Redeemer of us all,

\hymnSecondLine{Protect thy servants when they call,}
\begin{hangparas}{1.25em}{1}

And hear with reconciling care

The blessed Virgin's holy prayer.\\

Be ever present, Angel high

Whose name `God's might' doth signify:

To all the weak new strength impart,

And solace to the sad of heart.\\

And ye, O ever blissful throng

Of heav'nly Spirits, guardians strong,

Our past and present ills dispel,

From future peril shield us well.\\

From lands wherein thy faithful dwell

Drive far away the infidel;

So we to Christ due hymns of praise

Henceforth with eager hearts may raise.\\

To thee, O Father, born of none,

And thee, O sole-begotten Son,

One with the Holy Paraclete,

Be glory ever, as is meet. Amen.\\
\end{hangparas}

    ℣. An Angel stood at the altar of the temple.

	℟. Having in his hand a golden censer.
	
	\mattinsAntiphon{The Angel Gabriel {\dag} descended to Zacharias, and said unto him: Thy wife shall bear thee a son, and thou shalt call his name John; and many shall rejoice at his birth; for he shall go before the face of the Lord, to prepare his ways.}

\fussy
\end{paracol}

%MANUAL ADJUSTMENT:
\clearpage
\secondEvensong

\begin{rubric}
	Office Hymn of I Evensong.
\end{rubric}

    ℣. In the presence of the Angels I will sing praise unto thee, O my God.

	℟. I will worship toward thy holy temple, and praise thy Name.
	
	\properantiphon{Mag.}{The Archangel Gabriel {\dag} said unto Mary: With God nothing shall be impossible. And Mary said, Behold the handmaid of the Lord: be it unto me according to thy word. And the Angel departed from her.}

%MANUAL ADJUSTMENT:
\vspace{-0.5\baselineskip}

\collect
\lett{O}{God,} who from the company of Angels didst choose the Archangel Gabriel to proclaim the mystery of thine Incarnation: mercifully grant; that we who celebrate his festival (commemoration) on earth, may perceive his advocacy in heaven. Who livest.


\festivalLL{First Class Double}{25 March}{Annunciation of the Blessed Virgin Mary}
\feastday{Annunciation B.V.M.}
\feastdate{25}{March}

\begin{rubric}
	Common of the Blessed Virgin Mary (p. \pageref{CommonBVM}).
\end{rubric}

\openingSentence{Send, O Lord, the Lamb, the Ruler of the land, from the rock of the desert unto the mount of the daughter of Sion.}

%MANUAL ADJUSTMENT:
\vspace{-1.5\baselineskip}

\firstEvensong
\psalmAntiphon{Hail Mary, {\dag} thou that art full of grace, the Lord is with thee: blessed art thou among women.}\\

    ℣. Hail Mary, full of grace.

	℟. The Lord is with thee.
	
	\properantiphon{Mag.}{The Holy Ghost {\dag} shall come upon thee, Mary: and the power of the Highest shall overshadow thee.}
	
%MANUAL ADJUSTMENT:
\vspace{-0.5\baselineskip}

\mattins
\psalmAntiphon{The Lord shall give unto him {\dag} the throne of his father David, and he shall reign for ever.}\\
	
	℣. Hail Mary, full of grace.

	℟. The Lord is with thee.

	\properantiphon{Ben.}{How shall this be, {\dag} O Angel of God, seeing I know not a man? Hearken, O Virgin Mary: The Holy Ghost shall come upon thee, and the power of the Highest shall overshadow thee.}
	
\secondEvensong
\psalmAntiphon{Behold the handmaid of the Lord; {\dag} be it unto me according unto thy word.}\\
	
	℣. Hail Mary, full of grace.

	℟. The Lord is with thee.

	\properantiphon{Mag.}{The Angel Gabriel {\dag} spake unto Mary, saying: Hail, thou that art full of grace, the Lord is with thee; blessed art thou among women.}

\collect
\lett{O}{God,} who wast pleased that thy Word should take flesh of the womb of the Blessed Virgin Mary at the message of an Angel: grant to thy humble servants; that we who believe her to be truly the Mother of God may be aided by her intercession in thy sight. Through the same.
\begin{rubric}
    Commemoration of the Feria.% in non-conventual Masses.
\end{rubric}


\festival{Double (m.t.v.)}{27 March}{St. John Damascene}{C.D.}{FestiveWhite}
\feastday{{St. John Damascene}}
\feastdate{27}{March}

\begin{rubric}
	Commemoration only.
\end{rubric}

%\begin{rubric}
%	The Common of a Confessor not a Bishop (p. \pageref{CommonConfessorNotBishop}).
%\end{rubric}

%\collect
\lett{A}{lmighty} and everlasting God, who, for the defence of the veneration of sacred images, didst endue blessed John with heavenly doctrine and wondrous strength of spirit: grant unto us by his intercession and example; that we may imitate the virtues and perceive the advocacy of those whose images we venerate. Through.


\subbymemorial{30 March}{St. John Climacus}{Ab.}{FestiveWhite}
\feastday{{St. John Climacus}}
\feastdate{30}{March}

\begin{rubric}
	Common of a Confessor not a Bishop (p. \pageref{AbbotCollect}).
\end{rubric}


\subbymemorial{31 March}{St. Innocent of Alaska}{B.C.}{FestiveWhite}
\feastday{{St. Innocent Alaska}}
\feastdate{31}{March}

\begin{rubric}
	Common of a Confessor Bishop (p. \pageref{CommonConfessorBishop}), second Collect.
\end{rubric}


%MANUAL ADJUSTMENT:
\clearpage
\festival{Double}{4 April}{St. Isidore of Seville}{B.C.D.}{FestiveWhite}
\feastday{{St. Isidore Seville}}
\feastdate{4}{April}

\begin{rubric}
	Commemoration only.
\end{rubric}

\begin{rubric}
	Common of a Confessor Bishop (p. \pageref{ConfessorBishopDoctorCollect}).
\end{rubric}


\festival{First Class Double with a Common Octave}{7 April}{St. Tikhon of Moscow}{B.C.}{FestiveWhite}
\feastday{{St. Tikhon Moscow}}
\feastdate{7}{April}

\begin{rubric}
	Common of a Confessor Bishop (p. \pageref{CommonConfessorBishop}).
\end{rubric}

\begin{rubric}
    In Lent, Commemoration of the Feria.
\end{rubric}

\lett{O}{Lord} God of the nations, who for the defence of the Orthodox Faith didst raise up and strengthen blessed Tikhon, thy Bishop and Confessor, with great courage and self-denial: mercifully grant; that by his example and intercession, the hearts of the erring may be turned to the wisdom of the just, and the faithful may ever persevere in confession of thy true Religion. Through.


\festival{Double (m.t.v.)}{11 April}{Pope St. Leo the Great}{B.C.D.}{FestiveWhite}
\feastday{{Pope St. Leo the Great}}
\feastdate{11}{April}

\begin{rubric}
	Commemoration only.
\end{rubric}

%\begin{rubric}
%	The Common of a Confessor Bishop (p. \pageref{CommonConfessorBishop}).
%\end{rubric}

%\collect
\lett{W}{e} beseech thee, O Lord, graciously to hear the prayers which we offer unto thee on the solemnity of blessed Leo, thy Confessor and Bishop: that, like as he was found worthy to do thee faithful service, so by his merits and intercession we may be absolved from all our sins. Through.


\subbymemorial{13 April}{St. Hermengild}{M.}{MartyrRed}
\feastday{{St. Hermengild}}
\feastdate{13}{April}

%\begin{rubric}
%	In Lent, the Common of a Martyr (p. \pageref{CommonMartyr}). And Commemoration of the Feria.
%\end{rubric}

%\begin{rubric}
%	In Eastertide, the Common of a Martyr in Eastertide (p. \pageref{CommonMartyrEaster}).
%\end{rubric}

\memorialCollect
\lett{O}{God,} who didst teach blessed Hermengild, thy Martyr, to lay down an earthly for a heavenly kingdom: grant to us, we beseech thee; by his example to despise things temporal, and to seek after things eternal. Through.


%MANUAL ADJUSTMENT:
\clearpage
\subbymemorial{14 April}{St. Justin}{M.}{MartyrRed}
\feastday{{St. Justin}}
\feastdate{14}{April}

\memorialCollect
\lett{O}{God,} who through the foolishness of the Cross didst wondrously teach blessed Justin Martyr the excellent knowledge of Jesus Christ: grant to us through his intercession; that we, driving away the errors that beset us, may attain unto steadfastness of faith. Through the same.

\begin{rubric}
	Commemoration \blackrubric{of Sts. Tiburtius, Valerian, \& Maximus} (p. \pageref{TiburtiusCollect}).
\end{rubric}


\subbymemorial{14 April}{Sts. Tiburtius, Valerian, and Maximus}{M.}{MartyrRed}
\feastday{{Sts. Tiburtius \&c.}}
\feastdate{14}{April}

%\begin{rubric}
%	Outside of Eastertide, the Common of Many Martyrs (p. \pageref{CommonMartyrs}).
%\end{rubric}

%\begin{rubric}
%	In Eastertide, the propers are from the Common of Many Martyrs in Eastertide (p. \pageref{CommonMartyrsEaster}).
%\end{rubric}

\memorialCollect\label{TiburtiusCollect}
\lett{G}{rant,} we beseech thee, Almighty God: that we, who celebrate the festival of thy holy Martyrs, Tiburtius, Valerian, and Maximus; may also imitate their virtues. Through.


\subbymemorial{17 April}{St. Anicetus}{B.M.}{MartyrRed}
\feastday{{St. Anicetus}}
\feastdate{17}{April}

\begin{rubric}
	Outside of Eastertide, the Common of a Martyr (p. \pageref{CommonMartyr}), the second Collect.
\end{rubric}

\begin{rubric}
	In Eastertide, the Common of a Martyr in Eastertide (p. \pageref{CommonMartyrEaster}), the second Collect.
\end{rubric}


\subbymemorial{22 April}{Sts. Soter and Caius}{B.M.}{MartyrRed}
\feastday{{Sts. Soter \& Caius}}
\feastdate{22}{April}

\begin{rubric}
	Outside of Eastertide, the Common of Many Martyrs (p. \pageref{CommonMartyrs}).
\end{rubric}

\begin{rubric}
	In Eastertide, the Common of Many Martyrs in Eastertide (p. \pageref{CommonMartyrsEaster}).
\end{rubric}


\festival{First Class Double with a Common Octave}{23 April}{St. George}{M.}{MartyrRed}
\feastday{St. George}
\feastdate{23}{April}

\begin{rubric}
	Outside of Eastertide, the Common of a Martyr (p. \pageref{CommonMartyr}).
\end{rubric}

\begin{rubric}
	In Eastertide, the Common of a Martyr in Eastertide (p. \pageref{CommonMartyrEaster}).
\end{rubric}

%\collect
\lett{O}{God,} who causest us to rejoice in the good deeds and intercession of Saint George, thy Martyr, mercifully grant that, by the gift of thy grace, we may obtain the benefits we ask of him. Through.
\begin{rubric}
    In Lent, Commemoration of the Feria.% in non-conventual Masses.
\end{rubric}


%MANUAL ADJUSTMENT:
\clearpage
\festival{Second Class Double}{25 April}{St. Mark}{E.}{FestiveWhite}
\feastday{St. Mark}
\feastdate{25}{April}

\begin{rubric}
	Outside of Eastertide, the Common of Apostles (p. \pageref{CommonApostles}).
\end{rubric}

\begin{rubric}
	In Eastertide, the Common of Apostles in Eastertide (p. \pageref{CommonApostlesEaster}).
\end{rubric}

%\collect
\lett{O}{God,} who hast exalted blessed Mark thine Evangelist by the grace of the preaching of the Gospel: grant, we beseech thee; that we may ever both profit by his learning, and be defended by his prayer. Through.
\begin{rubric}
    In Lent, Commemoration of the Feria.% in non-conventual Masses.
\end{rubric}


\subbymemorial{26 April}{Sts. Cletus and Marcellinus}{M.}{MartyrRed}
\feastday{{Sts. Cletus \& Marcellinus}}
\feastdate{26}{April}

%\begin{rubric}
%	Outside of Eastertide, the Common of Many Martyrs (p. \pageref{CommonMartyrs}).
%\end{rubric}
%
%\begin{rubric}
%	In Eastertide, the Common of Many Martyrs in Eastertide (p. \pageref{CommonMartyrsEaster}).
%\end{rubric}

\memorialCollect
\lett{O}{Lord,} let the precious confession of thy blessed Martyrs and Bishops Cletus and Marcellinus be our defence: and let their loving intercession be a continual defence. Through.

\lettspace


\subbymemorial{28 April}{St. Vitalis}{M.}{MartyrRed}
\feastday{{St. Vitalis}}
\feastdate{28}{April}

\begin{rubric}
	Outside of Eastertide, the Common of a Martyr (p. \pageref{CommonMartyrNotBishopCollect}).
\end{rubric}

\begin{rubric}
	In Eastertide, the Common of a Martyr in Eastertide (p. \pageref{CommonMartyrNotBishopEasterCollect}).
\end{rubric}


\festival{Second Class Double}{1 May}{Sts. Philip \& James}{A.}{FestiveWhite}
\feastday{Sts. Philip \& James}
\feastdate{1}{May}

\begin{rubric}
	Common of Apostles in Eastertide (p. \pageref{CommonApostlesEaster}).
\end{rubric}

\openingSentence{The Lord was known of them in breaking of bread. Alleluia.}

%MANUAL ADJUSTMENT:
\vspace{-1\baselineskip}

\firstEvensong
\psalmAntiphon{Lord, {\dag} shew us the Father, and it sufficeth us, alleluia.}\\

\properantiphon{Mag.}{Let not your heart {\dag} be troubled, neither let it be afraid; ye believe in God, believe also in me: in my Father's house are many mansions, alleluia, alleluia.}

\mattins
\psalmAntiphon{Have I been {\dag} so long with you, and yet hast thou not known me, Philip? he that hath seen me hath seen my Father also, alleluia.}\\

\properantiphon{Ben.}{I am the way, {\dag} the truth, and the life: no man cometh unto the Father, but by me, alleluia.}

\secondEvensong
\psalmAntiphon{If ye love me, {\dag} keep my commandments, alleluia, alleluia, alleluia.}\\

	℣. Right dear in the sight of the Lord, alleluia.

	℟. Is the death of his Saints, alleluia.

\properantiphon{Mag.}{If ye abide in me, {\dag} and my words abide in you, ye shall ask what ye will, and it shall be done unto you, alleluia, alleluia, alleluia.}

\collect
\lett{O}{Almighty} God, whom truly to know is everlasting life; Grant us perfectly to know thy Son Jesus Christ to be the way, the truth, and the life; that, following the steps of thy holy Apostles, Saint Philip and Saint James, we may stedfastly walk in the way that leadeth to eternal life through the same thy Son Jesus Christ our Lord. Who liveth.


\festival{Double}{2 May}{Pope St. Athanasius of Alexandria}{B.C.D.}{FestiveWhite}
\feastday{{St. Athanasius}}
\feastdate{2}{May}

\begin{rubric}
	Common of a Confessor Bishop (p. \pageref{CommonConfessorBishop}).
\end{rubric}

%\collect
\lett{W}{e} beseech thee, O Lord, graciously to hear the prayers which we offer unto thee on the solemnity of blessed Athanasius, thy Confessor and Bishop: that, like as he was found worthy to do thee faithful service, so by his merits and intercession we may be absolved from all our sins. Through.


%MANUAL ADJUSTMENT:
\clearpage
\festivalLL{Second Class Double}{3 May}{Invention of the Holy Cross}\label{InventionCross}
\feastday{Invention of the Cross}
\feastdate{3}{May}

\openingSentence{Tell it out among the heathen that the Lord hath reigned from the tree. Alleluia.}

\begin{paracol}{2}
\sloppy

%MANUAL ADJUSTMENT:
\vspace{-2\baselineskip}

\firstEvensong
\psalmAntiphon{O mighty {\dag} work of mercy! death then died, when life died upon the Tree, alleluia.}

\firstEvensongHymn
\hymnLett{T}{he} royal banners forward go;

\hymnSecondLine{The Cross shines forth in mystic glow;}

\begin{hangparas}{1.25em}{1}

Where he in flesh, our flesh who made,

Our sentence bore, our ransom paid.\\

2. Where deep for us the spear was dy'd,

Life's torrent rushing from his side,

To wash us in that precious flood

Where mingled Water flowed and Blood.\\

3. Fulfill'd is all that David told

In true prophetic song of old:

Amidst the nations, God, saith he,

Hath reign'd and triumph'd from the Tree.\\

4. O Tree of beauty! Tree of light!

O Tree with royal purple dight!

Elect on whose triumphal breast

Those holy limbs should find their rest:\\

5. On whose dear arms, so widely flung,

The weight of this world’s ransom hung,

The price of humankind to pay,

And spoil the spoiler of his prey.
\end{hangparas}

\begin{rubric}
	During the following stanza all kneel, and the last stanza is never changed.
\end{rubric}

\begin{hangparas}{1.25em}{1}

6. O Cross, our one reliance, hail!

In this our Paschal joy avail

To give fresh merit to the saint

And pardon to the penitent.\\

7. To thee, eternal Three in One,

Let homage meet by all be done:

Whom by the Cross thou dost restore,

Preserve and govern evermore. Amen.\\
\end{hangparas}

℣. This sign of the Cross shall be in heaven, alleluia.

℟. When the Lord shall come to judgement, alleluia.

\firstEvensongAntiphon{O Cross, {\dag} surpassing all the stars in splendour, world renowned, exceeding dear unto the hearts of men, holier than all things: thou only wert counted worthy to uphold the world’s random. Sweet the wood, sweet the iron, bearing so sweet a burden: bring aid to this congregation, who are here assembled to celebrate thy praises, alleluia, alleluia.}


\secondEvensong
\psalmAntiphon{But as for us, {\dag} it behoveth us to glory in the Cross of our Lord Jesus Christ, alleluia.}

\begin{rubric}
	Office Hymn of I Evensong.
\end{rubric}

℣. This sign of the Cross shall be in heaven, alleluia.

℟. When the Lord shall come to judgement, alleluia.

\secondEvensongAntiphon{He the holy Cross endured, {\dag} Burst the gates of hell in twain: Begirt with might and majesty, On Easter morn he rose again, alleluia.}


\collect
\lett{O}{God,} who in the wondrous Finding of the Cross of salvation didst renew the miracles of thy Passion: vouchsafe; that by the ransom of the tree of life we may attain thy succour unto life eternal. Who livest.
\begin{rubric}
	 Commemoration \blackrubric{of Sts. Alexander, Eventius, Theodulus, \& Juvenalis} (p. \pageref{AlexanderEtcCollect}).
\end{rubric}


\switchcolumn

%MANUAL ADJUSTMENT:
\vspace{-2\baselineskip}

\mattins
\psalmAntiphon{Behold the Cross of the Lord; {\dag} flee away, ye adversaries: the Lion of Judah, the Root of David, hath prevailed, alleluia.}

\mattinsInvitatoryHymn
\hymnLett{S}{ing,} my tongue, the glorious battle

\hymnSecondLine{Sing the last, the dread affray;}

\begin{hangparas}{1.25em}{1}

O'er the Cross, the victor's trophy,

Sound the high triumphal lay:

Tell how Christ, the world's Redeemer,

As a victim won the day.\\

2. God, our Maker, sorely grieving

That the first-made Adam fell,

When he ate the fruit of sorrow,

Whose reward was death and hell,

Noted then this Wood, the ruin

Of the ancient wood to quell.\\

3. For the work of our salvation

Needs would have his order so,

And the multiform deceiver's

Art by art would overthrow,

And from thence would bring the med'cine

Whence the insult of the foe.\\

4. Wherefore, when the sacred fulness

Of th'appointed time was come,

This world's Maker left his Father,

Sent the heav'nly mansion from,

And proceeded, God Incarnate,

Of the Virgin's holy womb.\\

5. Weeps the Infant in the manger

That in Bethl'hem's stable stands;

And his limbs the Virgin Mother

Doth compose in swaddling bands,

Meetly thus in linen folding

Of her God the feet and hands.\\

6. To the Trinity be glory

Everlasting, as is meet;

Equal to the Father, equal

To the Son, and Paraclete:

Trinal Unity, whose praises

All created things repeat. Amen.
\end{hangparas}

\mattinsOfficeHymn
\hymnLett{T}{hirty} years among us dwelling,

\hymnSecondLine{His appointed time fulfill'd,}

\begin{hangparas}{1.25em}{1}
Born for this, he meets his Passion,

For that this he freely will'd:

On the Cross the Lamb is lifted,

Where his life-blood shall be spill'd.\\

2. He endured the nails, the spitting,

Vinegar, and spear, and reed:

From that holy Body broken

Blood and Water forth proceed:

Earth, and stars, and sky, and ocean,

By that flood from stain are freed.\\

3. Faithful Cross, above all other

One and only noble Tree;

None in foliage, none in blossom,

None in fruit thy peer may be:

Sweetest wood, and sweetest iron!

Sweetest weight is hung on thee.\\

4. Bend thy boughs, O Tree of glory,

Thy relaxing sinews bend:

For awhile the ancient rigour

That thy birth bestow'd, suspend:

And the King of heav'nly beauty

On thy bosom gently tend.\\

5. Thou alone wast counted worthy

This world's ransom to sustain;

That a shipwreck'd race for ever

Might a port of refuge gain:

With the sacred Blood anointed

Of the Lamb for sinners slain.\\

6. Glory be to God, and honour

In the highest, as is meet,

To the Son, and to the Father,

And th' eternal Paraclete,

Whose is boundless praise and power

Through the ages infinite. Amen.\\
\end{hangparas}

℣. We adore thee, O Christ, and we bless thee, alleluia.

℟. Because by thy Cross thou hast redeemed the world, alleluia.

\mattinsAntiphon{Thou alone {\dag} excellest in stature all the cedars of Lebanon: for on thee the Life of the world was hanged, on thee was Christ victorious, and death over death did for ever triumph, alleluia.}

\fussy	
\end{paracol}


%MANUAL ADJUSTMENT:
\vspace{-0.25\baselineskip}

\subbymemorial{3 May}{Sts. Pope Alexander, Eventius, Theodulus, and Bishop Juvenalis}{M.}{MartyrRed}
\feastday{{Sts. Alexander \&c.}}
\feastdate{3}{May}

%\begin{rubric}
%	The Common of Many Martyrs in Eastertide (p. \pageref{CommonMartyrsNotBishopsEasterCollect}).
%\end{rubric}

\memorialCollect\label{AlexanderEtcCollect}
\lett{G}{rant,} we beseech thee, Almighty God: that we, who devoutly celebrate the birthday of thy Saints Alexander, Eventius, Theodulus, and Juvenal; may, by their intercession, be delivered from all evils that beset us. Through.


%MANUAL ADJUSTMENT:
\vspace{-0.5\baselineskip}

\subbymemorial{4 May}{St. Monica}{W.}{FestiveWhite}
\feastday{{St. Monica}}
\feastdate{4}{May}

%\begin{rubric}
%	The Common of Holy Matrons (p. \pageref{CommonHolyMatrons}).
%\end{rubric}

\memorialCollect
\lett{O}{God,} the comforter of them that mourn, and the salvation of them that hope in thee, who didst mercifully receive the loving tears of blessed Monica for the conversion of Augustine her son: grant to us by the intercession of them both; that we may bewail our sins, and obtain the pardon of thy grace. Through.


%MANUAL ADJUSTMENT:
\vspace{-0.25\baselineskip}

\festival{Greater Double}{6 May}{St. John before the Latin Gate}{A.}{FestiveWhite}
\feastday{St. John Latin Gate}
\feastdate{6}{May}

\begin{rubric}
	Common of Apostles in Easter (p. \pageref{CommonApostlesEaster}).
\end{rubric}

\openingSentence{Right worthy of honour is blessed John the Apostle, who leaned on the Lord's bosom at the last supper.}

\antiphon{I \& II Evensong Mag.}{John the Apostle, {\dag} being cast into a caldron of boiling oil, by virtue of protecting grace came forth unharmed, alleluia.}

\collect
\lett{O}{God,} who seest that we are beset by evils on every side: grant, we beseech thee; that the glorious intercession of blessed John, thine Apostle and Evangelist, may protect us. Through.


%MANUAL ADJUSTMENT:
\clearpage
\subbymemorial{7 May}{St. Alexis Toth}{C.}{FestiveWhite}
\feastday{{St. Alexis Toth}}
\feastdate{7}{May}

%CHECK:
\begin{rubric}
	Common of a Confessor not a Bishop (p. \pageref{ConfessorNotBishopCollect}).
\end{rubric}


\festivalAO{Greater Double}{8 May}{The Apparition of St. Michael the Archangel}
\feastday{{St. Michael Apparition}}
\feastdate{8}{May}

\begin{rubric}
	Hymns, Versicles, \& Antiphons of \blackrubric{Michaelmas} (p. \pageref{Michaelmas}).
\end{rubric}

%\collect
\lett{O}{everlasting} God, who hast ordained and constituted the services of Angels and men in a wonderful order; Mercifully grant that, as thy holy Angels always do thee service in heaven, so, by thy appointment, they may succour and defend us on earth. Through.
\begin{rubric}
	Commemoration \blackrubric{of Pope St. Boniface IV of Rome,} from the Common of a Confessor Bishop (p. \pageref{ConfessorBishopCollect}).
\end{rubric}

%8 May is the latest day of Easter


\subbymemorial{8 May}{Pope St. Boniface IV of Rome}{B.C.}{FestiveWhite}
\feastday{{Pope St. Boniface IV}}
\feastdate{8}{May}

%CHECK:
\begin{rubric}
	Common of a Confessor Bishop (p. \pageref{ConfessorBishopCollect}).
\end{rubric}


\festival{Double}{9 May}{St. Gregory Nazianzen}{B.C.D.}{FestiveWhite}
\feastday{{St. Gregory Nazianzen}}
\feastdate{9}{May}

\begin{rubric}
	Common of a Confessor Bishop (p. \pageref{CommonConfessorBishop}).
\end{rubric}


\subbymemorial{10 May}{Sts. Gordian \& Epimachus}{M.}{MartyrRed}
\feastday{{Sts. Gordian \& Epimachus}}
\feastdate{10}{May}

%\begin{rubric}
%	The Common of Many Martyrs in Eastertide (p. \pageref{CommonMartyrsEaster}).
%\end{rubric}

\memorialCollect
\lett{G}{rant,} we beseech thee, Almighty God: that we, who celebrate the festival of thy blessed Martyrs Gordian and Epimachus, may be aided by their intercession with thee. Through.


\subbymemorial{12 May}{Sts. Nereus, Achilles, Domitilla, and Pancras}{M.}{MartyrRed}
\feastday{{Sts. Nereus \&c.}}
\feastdate{12}{May}

\memorialCollect
\lett{W}{e} beseech thee, O Lord, that the blessed solemnity of thy Martyrs, Nereus, Achilles, Domitilla, and Pancras, may ever protect us: and render us worthy of thy service. Through.


%MANUAL ADJUSTMENT:
\clearpage
\subbymemorial{14 May}{St. Boniface of Tarsus}{M.}{MartyrRed}
\feastday{{St. Boniface Tarsus}}
\feastdate{14}{May}

%\begin{rubric}
%	The Common of a Martyr in Eastertide (p. \pageref{CommonMartyrEaster}).
%\end{rubric}

\memorialCollect
\lett{G}{rant,} we beseech thee, Almighty God: that we, who celebrate the festival of blessed Boniface thy Martyr, may be aided by his intercession with thee. Through.


\festival{Double}{18 May}{St. Venantius}{M.}{MartyrRed}
\feastday{{St. Venantius}}
\feastdate{18}{May}

\begin{rubric}
	Common of a Martyr in Eastertide (p. \pageref{CommonMartyrEaster}).
\end{rubric}

%\collect
\lett{O}{God,} who hast hallowed this day by the triumph of blessed Venantius, thy Martyr: graciously hear the prayers of thy people, and grant; that we who venerate his merits may imitate the constancy of his faith. Through.



\subbymemorial{19 May}{St. Pudentiana}{V.}{FestiveWhite}
\feastday{{St. Pudentiana}}
\feastdate{19}{May}

\begin{rubric}
	Common of Virgins (p. \pageref{VirginOnlyCollect}).
\end{rubric}


\subbymemorial{22 May}{St. Romanus}{Ab.}{FestiveWhite}
\feastday{{St. Romanus}}
\feastdate{22}{May}

\begin{rubric}
	Common of a Confessor not a Bishop (p. \pageref{AbbotCollect}).
\end{rubric}


\subbymemorial{24 May}{St. Vincent of Lerins}{C.}{FestiveWhite}
\feastday{{St. Vincent Lerins}}
\feastdate{24}{May}

\begin{rubric}
	Common of a Confessor not a Bishop (p. \pageref{ConfessorNotBishopCollect}).
\end{rubric}


\subbymemorial{25 May}{Pope St. Urban I of Rome}{B.M.}{MartyrRed}
\feastday{{Pope St. Urban I}}
\feastdate{25}{May}

\begin{rubric}
	Common of a Martyr (p. \pageref{CommonMartyrBishopEasterCollect}).
\end{rubric}


\festival{Double}{26 May}{{St. Augustine of Canterbury}}{B.C.}{FestiveWhite}
\feastday{{\small St. Augustine Canterbury}}
\feastdate{26}{May}

\begin{rubric}
	Common of a Confessor Bishop (p. \pageref{CommonConfessorBishop}).
\end{rubric}

%\collect
\lett{O}{God,} who didst give the blessed Bishop Augustine to be the first Teacher of the English people: grant us, we beseech thee; that we who proclaim his merits on earth may perceive his intercession in heaven. Through.

\begin{rubric}
	Commemoration \blackrubric{of St. Eleutherius} (p. \pageref{EleutheriusCollect}).
\end{rubric}


%MANUAL ADJUSTMENT:
\clearpage
\subbymemorial{26 May}{Pope St. Eleutherius of Rome}{B.M.}{MartyrRed}
\feastday{{Pope St. Eleutherius}}
\feastdate{26}{May}

%\begin{rubric}
%	The Common of a Martyr in Eastertide (p. \pageref{CommonMartyrEaster}).
%\end{rubric}

\memorialCollect\label{EleutheriusCollect}
\lett{A}{lmighty} God, mercifully look upon our infirmities: that, whereas we are oppressed by the burden of our sins, the glorious intercession of blessed Eleutherius, thy Martyr and Bishop, may be our succour and defence. Through.


\festival{Double (m.t.v.)}{27 May}{St. Bede the Venerable}{C.D.}{FestiveWhite}
\feastday{{St. Bede Venerable}}
\feastdate{27}{May}

\begin{rubric}
	Common of a Confessor not a Bishop (p. \pageref{CommonConfessorNotBishop}).
\end{rubric}

%\collect
\lett{O}{God,} who dost illumine thy Church with the learning of blessed Bede, thy Confessor and Doctor: mercifully grant unto thy servants; that they may ever be enlightened by his wisdom, and succoured by his merits. Through.

\begin{rubric}
	Commemoration \blackrubric{of Pope St. John I of Rome,} from the Second Collect of the Common of a Martyr in Eastertide (p. \pageref{CommonMartyrBishopEasterCollect}).
\end{rubric}


\subbymemorial{27 May}{Pope St. John I of Rome}{B.M.}{MartyrRed}
\feastday{{Pope St. John I}}
\feastdate{27}{May}

\begin{rubric}
	Common of a Martyr in Eastertide (p. \pageref{CommonMartyrBishopEasterCollect}), the second Collect.
\end{rubric}


\subbymemorial{30 May}{Pope St. Felix I of Rome}{B.M.}{MartyrRed}
\feastday{{Pope St. Felix I}}
\feastdate{30}{May}

\begin{rubric}
	In Eastertide, the Common of a Martyr in Eastertide (p. \pageref{CommonMartyrBishopEasterCollect}).
\end{rubric}

\begin{rubric}
	Outside of Eastertide, the Common of a Martyr (p. \pageref{CommonMartyrBishopCollect}).
\end{rubric}


\subbymemorial{31 May}{St. Petronilla}{V.}{FestiveWhite}
\feastday{{St. Petronilla}}
\feastdate{31}{May}

\begin{rubric}
	Common of Virgins (p. \pageref{VirginOnlyCollect}), the second Collect.
\end{rubric}


\subbymemorial{2 June}{Sts. Marcellinus, Peter, \& Erasmus}{M.}{MartyrRed}
\feastday{{Sts. Marcellinus \&c.}}
\feastdate{2}{June}

\memorialCollect
\lett{O}{God,} who makest us glad with the yearly solemnity of thy blessed Martyrs, Marcellinus, Peter, and Erasmus: grant, we beseech thee; that, as we rejoice in their merits, so we may be enkindled by their example. Through.


%MANUAL ADJUSTMENT:
\clearpage
\festival{Double}{5 June}{St. Boniface}{B.M.}{MartyrRed}
\feastday{{St. Boniface}}
\feastdate{5}{June}

\begin{rubric}
	In Eastertide, the Common of a Martyr in Eastertide (p. \pageref{CommonMartyrEaster}).
\end{rubric}

\begin{rubric}
	Outside of Eastertide, the Common of a Martyr (p. \pageref{CommonMartyr}).
\end{rubric}

%\collect
\lett{O}{God,} who by the zeal of blessed Boniface, thy Martyr and Bishop, didst vouchsafe to call a multitude of peoples to the knowledge of thy name: mercifully grant; that we who celebrate his festival may likewise perceive his advocacy. Through.


\festival{Double}{9 June}{St. Columba of Iona}{Ab.}{FestiveWhite}
\feastday{{St. Columba Iona}}
\feastdate{9}{June}


\begin{rubric}
	Common of a Confessor not a Bishop (p. \pageref{CommonConfessorNotBishop}), Commemoration \blackrubric{of Sts. Primus \& Felician} (p. \pageref{PrimusFelician}).
\end{rubric}


\subbymemorial{9 June}{Sts. Primus \& Felician}{M.}{MartyrRed}\label{PrimusFelician}
\feastday{{Sts. Primus \& Felician}}
\feastdate{9}{June}

\memorialCollect
\lett{O}{Lord,} we beseech thee, make us ever to rejoice in the festival of thy holy Martyrs Primus and Felician: that by their prayers we may perceive the gifts of thy protection. Through.


\subbymemorial{10 June}{Queen St. Margaret of Scotland}{Ma.}{FestiveWhite}
\feastday{{St. Margaret Scotland}}
\feastdate{10}{June}

\memorialCollect
\lett{O}{God,} who didst render the blessed Queen Margaret wondrous by reason of her eminent charity towards the poor: grant; that, by her intercession and example, thy charity may continually increase in our hearts. Through.


\festival{Greater Double}{11 June}{St. Barnabas}{A.}{FestiveWhite}
\feastday{St. Barnabas}
\feastdate{11}{June}

\begin{rubric}
	Outside of Eastertide, the Common of Apostles (p. \pageref{CommonApostles}).
\end{rubric}

\begin{rubric}
	In Eastertide, the Common of Apostles in Eastertide (p. \pageref{CommonApostlesEaster}).
\end{rubric}

%\collect
\lett{W}{e} beseech thee, O Lord, let the prayers of thy Apostle, Saint Barnabas, commend thy Church to thee, and let him continue to intercede for her whom by his doctrine and death he doth glorify. Through.


%MANUAL ADJUSTMENT:
\clearpage
\subbymemorial{12 June}{Sts. Basilides, Cyrinus, Nabor, and Nazarius}{M.}{MartyrRed}
\feastday{{Sts. Basilides \&c.}}
\feastdate{12}{June}

\memorialCollect
\lett{W}{e} beseech thee, O Lord, that the observance of the birthday of thy holy Martyrs Basilides, Cyrinus, Nabor, and Nazarius may shine brightly upon us: that the everlasting excellence which they have obtained may increase by the fruits of our devotion. Through.


\festival{Greater Double (m.t.v.)}{14 June}{St. Basil the Great}{B.C.D.}{FestiveWhite}
\feastday{{St. Basil the Great}}
\feastdate{14}{June}

\begin{rubric}
	Common of a Confessor Bishop (p. \pageref{CommonConfessorBishop}).
\end{rubric}

%\collect
\lett{W}{e} beseech thee, O Lord, graciously to hear the prayers which we offer unto thee on the solemnity of blessed Basil, thy Confessor and Bishop: that, like as he was found worthy to do thee faithful service, so by his merits and intercession we may be absolved from all our sins. Through.


\subbymemorial{15 June}{Sts. Vitus, Modestus, and Crescentia}{M.}{MartyrRed}
\feastday{{Sts. Vitus \&c.}}
\feastdate{15}{June}

\memorialCollect
\lett{G}{rant,} O Lord, we beseech thee, that, at the intercession of thy holy Martyrs, Vitus, Modestus, and Crescentia, thy Church may learn not to be highminded, but to grow in such lowliness as is acceptable unto thee: that, despising things evil, she may with bounteous charity perform those things that be right. Through.


\festival{Double}{18 June}{St. Ephrem the Syrian}{C.D.}{FestiveWhite}
\feastday{{St. Ephrem Syrian}}
\feastdate{18}{June}

\begin{rubric}
	Common of a Confessor not a Bishop (p. \pageref{CommonConfessorNotBishop}).
\end{rubric}

%\collect
\lett{O}{God,} who wast pleased to adorn thy Church with the wondrous learning and glorious merits of the life of blessed Ephram, thy Confessor and Doctor: we humbly entreat thee; that at his intercession thou wouldest defend her with thy continual power against the snares of error and wickedness. Through.

\begin{rubric}
	Commemoration \blackrubric{of Sts. Marcus \& Marcellianus} (p. \pageref{MarcusMarcellianusCollect}).
\end{rubric}


\subbymemorial{18 June}{Sts. Marcus and Marcellianus}{M.}{MartyrRed}
\feastday{{Sts. Marcus \&c.}}
\feastdate{18}{June}

\memorialCollect\label{MarcusMarcellianusCollect}
\lett{G}{rant,} we beseech thee, Almighty God: that we who devoutly celebrate the birthday of thy holy Martyrs Marcus and Marcellianus; may, by their intercession, be delivered from all evils that beset us. Through.


%MANUAL ADJUSTMENT:
\clearpage
\subbymemorial{19 June}{Sts. Gervase and Protase}{M.}{MartyrRed}
\feastday{{Sts. Gervase \& Protase}}
\feastdate{19}{June}

\memorialCollect
\lett{O}{God,} who makest us glad with the yearly solemnity of thy holy Martyrs Gervase and Protase: mercifully grant; that as we do rejoice in their merits, so we may be enkindled by their example. Through.


\subbymemorial{20 June}{Pope St. Silverius of Rome}{B.M.}{MartyrRed}
\feastday{{St. Silverius}}
\feastdate{20}{June}

\begin{rubric}
	Outside of Eastertide, the Common of a Martyr (p. \pageref{CommonMartyrBishopCollect}).
\end{rubric}

\begin{rubric}
	In Eastertide, the Common of a Martyr in Eastertide (p. \pageref{CommonMartyrBishopEasterCollect}).
\end{rubric}


\festival{Double}{22 June}{St. Alban}{M.}{MartyrRed}
\feastday{{St. Alban}}
\feastdate{22}{June}

\begin{rubric}
	Outside of Eastertide, the Common of a Martyr (p. \pageref{CommonMartyr}).
\end{rubric}

\begin{rubric}
	In Eastertide, the Common of a Martyr in Eastertide (p. \pageref{CommonMartyrEaster}).
\end{rubric}

%\collect
\lett{O}{God,} who hast hallowed this day by the martyrdom of blessed Alban: grant, we beseech thee; that as year by year we rejoice to pay him honour, so we may be defended by his continual help. Through.

\begin{rubric}
	Commemoration \blackrubric{of St. Paulinus} (p. \pageref{PaulinusCollect}).
\end{rubric}


\subbymemorial{22 June}{St. Paulinus}{B.C.}{FestiveWhite}
\feastday{{St. Paulinus}}
\feastdate{22}{June}

\memorialCollect\label{PaulinusCollect}
\lett{O}{God,} who to them that forsake all things in this world for thee hast promised a hundredfold in the world to come, and life eternal: mercifully grant; that, following in the footsteps of thy holy Bishop, Paulinus, we may be enabled to despise earthly things, and to seek only after things heavenly: Who livest.

\begin{rubric}
	If to-day be Saturday, Commemoration \blackrubric{of the Vigil of St. John Baptist} (p. \pageref{JohnBaptistVigilCollect}).
\end{rubric}


\festivalAO{Vigil}{23 June}{Vigil of St. John Baptist}
\feastday{{St. John Baptist Vigil}}
\feastdate{23}{June}

\vigilCollect\label{JohnBaptistVigilCollect}
\lett{G}{rant,} we beseech thee, Almighty God: that thy family may walk in the way of salvation; and, following the teachings of blessed John the Forerunner, may attain in safety unto him whom he fore-told, Jesus Christ thy Son, our Lord: Who liveth.


%MANUAL ADJUSTMENT:
\clearpage
\subbymemorial{23 June}{Queen St. Etheldreda}{V.}{FestiveWhite}
\feastday{{St. Etheldreda}}
\feastdate{23}{June}

\memorialCollect
\lett{O}{God,} who makest us glad with the yearly solemnity of blessed Etheldreda, thy Virgin: mercifully grant; that as we are enlightened by the example of her purity, so we may be succoured by her prayers. Through.


\festivalAO{First Class Double with a Common Octave}{24 June}{St. John Baptist}\label{StJohnBaptist}
\feastday{St. John Baptist}
\feastdate{24}{June}

\openingSentence{There was a man sent from God, whose name was John.}

\begin{paracol}{2}
\sloppy

\firstEvensong
\psalmAntiphon{He shall go before him {\dag} in the spirit and power of Elias, to make ready a people prepared for the Lord.}

\firstEvensongHymn
\hymnLett{O}{for} thy spirit, holy John, to chasten

\hymnSecondLine{Lips sin-polluted, fetter'd tongues to loosen;}

\begin{hangparas}{1.25em}{1}
So by thy children might thy deeds of wonder

Meetly be chanted.\\

2. Lo! a swift herald, from the skies descending,

Bears to thy father promise of thy greatness;

How he shall name thee, what thy future story,

Duly revealing.\\

3. Scarcely believing message so transcendent,

Him for a season pow'r of speech forsaketh,

Till, at thy wondrous birth, again returneth

Voice to the voiceless.\\

4. Thou, in thy mother’s womb all darkly cradled,

Knewest thy Monarch, biding in his chamber;

Whence the two parents, through their children’s merits,

Mysteries utter'd.\\

5. Praise to the Father, to the Sole-begotten,

And to thee, alway with the Twain co-equal,

Fostering Spirit; one and only Godhead

Through all the ages. Amen.\\	
\end{hangparas}

℣. There was a man sent from God.

℟. Whose name was John.

\firstEvensongAntiphon{When Zacharias {\dag} went into the temple, there appeared unto him the Angel Gabriel, standing on the right side of the altar of incense.}

\secondEvensong
\psalmAntiphon{Thou, child {\dag} shalt be called the prophet of the Highest; for thou shalt go before the face of the Lord, to prepare his ways.}

\begin{rubric}
	Office Hymn of I Evensong.
\end{rubric}

℣. This child shall be great in the sight of the Lord.

℟. For the hand of the Lord is with him.

\secondEvensongAntiphon{The child {\dag} that is born unto us is more than a prophet; for this is he of whom the Saviour saith: Among them that are born of women, there hath not risen a greater than John the Baptist.}

\collect
\label{JohnBaptistCollect}
\lett{A}{lmighty}\textls[-25]{ God, by whose providence thy servant John Baptist was wonderfully born,} and sent to prepare the way of thy Son our Saviour by preaching repentance; Make us so to follow his doctrine and holy life, that we may truly repent according to his preaching; and after his example constantly speak the truth, boldly rebuke vice, and patiently suffer for the truth's sake. Through the same.

\switchcolumn

\mattins
\psalmAntiphon{Elisabeth, {\dag} the wife of Zacharias, brought forth a mighty son, John the Baptist, the Lord's Forerunner.}

\mattinsInvitatoryHymn
\hymnLett{T}{hou} in thy childhood, to the desert caverns

\hymnSecondLine{Fleddest for refuge from the cities' turmoil,}

\begin{hangparas}{1.25em}{1}
Where the world's slander might not dim thy lustre,

Lonely abiding.\\

2. Camel's hair raiment cloth'd thy saintly members;

Leathern the girdle which thy loins encircled;

Locusts and honey, with the fountain-water,

Daily sustain'd thee.\\

3. Oft in past ages, seers with hearts expectant

Sang the far-distant advent of the Day-Star;

Thine was the glory, as the world’s Redeemer

First to proclaim him.\\

4. Far as the wide world reacheth, born of woman,

Holier was there none than John the Baptist;

Meetly in water laving him who cleanseth

Man from pollution.\\

5. Praise to the Father, to the Sole-begotten,

And to thee, alway with the Twain co-equal,

Fostering Spirit; one and only Godhead

Through all the ages. Amen.
\end{hangparas}

\mattinsOfficeHymn
\hymnLett{O}{more} than blessed, merit high attaining,

\hymnSecondLine{Pure as the snow-drift, innocent of evil,}

\begin{hangparas}{1.25em}{1}
Child of the desert, mightiest of martyrs,

Greatest of prophets.\\

2. Thirty-fold increase some with glory crowneth;

Sixty-fold fruitage prize for others winneth;

Hundred-fold measure, thrice repeated, decks thee,

Blest one, for guerdon.\\

3. O may the virtue of thine intercession,

All stony hardness from our hearts expelling,

Smooth the rough places, and the crooked straighten

Here in the desert.\\

4. Thus may our gracious Maker and Redeemer,

Seeking a station for his hallow'd footsteps,

Find, when he cometh, temples undefiled,

Meet to receive him.\\

5. Now as the Angels celebrate thy praises,

Godhead essential, Trinity co-equal;

Spare thy redeem'd ones, as they bow before thee,

Pardon imploring. Amen.\\
\end{hangparas}

℣. This child shall be great in the sight of the Lord.

℟. For the hand of the Lord is with him.

\mattinsAntiphon{The mouth of Zacharias {\dag} was opened, and he prophesied, saying: Blessed be the God of Israel.}


\fussy
\end{paracol}


%MANUAL ADJUSTMENT:
\clearpage
\festival{Double}{26 June}{Sts. John and Paul}{M.}{MartyrRed}
\feastday{{Sts. John \& Paul}}
\feastdate{26}{June}

\begin{rubric}
	Common of Many Martyrs (p. \pageref{CommonMartyrs}).
\end{rubric}

%\collect
\lett{W}{e} beseech thee, Almighty God: that, as their faith and passion did cause blessed John and Paul to be brothers indeed; so this day's festival may bestow upon us a twofold gladness in their glory. Through.

\begin{rubric}
	Commemoration \blackrubric{of the Octave of St. John Baptist} (p. \pageref{JohnBaptistCollect}).
\end{rubric}


\festival{Vigil}{28 June}{Vigil of St. Peter}{A.}{LentenViolet}
\feastday{{Sts. Peter \& Paul Vigil}}
\feastdate{28}{June}

%\collect
\lett{G}{rant,} we beseech thee, Almighty God: that as thou hast stablished us on the rock of the apostolic confession; so thou wouldest not suffer us to be troubled by any adversities. Through.

\begin{rubric}
	Commemoration \blackrubric{of the Octave of St. John Baptist} (p. \pageref{JohnBaptistCollect}).
\end{rubric}

\begin{rubric}
	Commemoration \blackrubric{of Pope St. Leo II,} from the Common of a Confessor Bishop (p. \pageref{ConfessorBishopCollect}).
\end{rubric}


\subbymemorial{28 June}{Pope St. Leo II of Rome}{B.C.}{FestiveWhite}
\feastday{{Pope St. Leo II}}
\feastdate{28}{June}

\begin{rubric}
	Common of a Confessor Bishop (p. \pageref{ConfessorBishopCollect}).
\end{rubric}


%MANUAL ADJUSTMENT:
\clearpage
%\festival{First Class Double with a Common Octave}{29 June}{Sts. Peter \& Paul}

\subsection{{\color{FestiveWhite}29 June. St. Peter,} {\color{RubricRed}A.}}\label{SaintPeter}

\centerline{\small{(Sts. Peter \& Paul)}}

\begin{inhead}
	First Class Double with a Common Octave
\end{inhead}

\feastday{St. Peter}
\feastdate{29}{June}

\openingSentence{Thou art Peter, and upon this rock I will build my church.}

\begin{paracol}{2}
\sloppy

\firstEvensong
\psalmAntiphon{The Angel said unto Peter, {\dag} Cast thy garment about thee, and follow me.}

\firstEvensongHymn
\hymnLett{W}{ith} golden splendour, and with roseate loveliness,

\hymnSecondLine{Thou didst illumine, Light of Light, the universe;}

\begin{hangparas}{1.25em}{1}
The heav'ns adorning with a glorious martyrdom

This day, which bringeth pardon to the penitent.\\

2. Celestial Warder! earth’s Instructor eloquent!

The world's dread judges, lights mankind enlightening,

By cross triumphant, by the sword victorious,

Now are ye laurell'd, Life’s immortal senators.\\

3. O Rome, thrice happy, that so great a martyrdom

Thy walls, empurpled with rich life-blood, halloweth!

Not thine the merit; by thy Princes' excellence

Thou shinest fairer than the world's magnificence.\\

4. Glory eternal to the Blessed Trinity,

With laud and honour, virtue and supremacy,

Trinal yet Onely, reigning in his majesty

Both now and ever, through the ages infinite. Amen.\\
\end{hangparas}

℣. Their sound is gone out into all lands.

℟. And their words into the ends of the world.

\firstEvensongAntiphon{Thou art the shepherd of the sheep, {\dag} O chief of the Apostles: unto thee were given the keys of the kingdom of heaven.}

\secondEvensong
\psalmAntiphon{Just as the glorious princes loved one another in their lives: {\dag} so also in their death were they not separated.}

\begin{rubric}
	Office Hymn of I Evensong.
\end{rubric}

℣. They declared the work of God.

℟. And wisely considered of his doing.

\secondEvensongAntiphon{To-day {\dag} Simon Peter ascended the gibbet of the cross, alleluia: to-day the Key-Bearer of the Kingdom joyfully departed to Christ: to-day Paul the Apostle, the light of the whole world, bowed down his head; and for the Name of Christ, received the crown of martyrdom, alleluia.}

\switchcolumn

\mattins
\psalmAntiphon{The Lord hath sent {\dag} his Angel, and hath delivered me out of the hand of Herod, alleluia.}

\mattinsInvitatoryHymn
\begin{rubric}
	Invitatory Hymn of the Common of the Apostles (\pageref{ApostlesInvitatory}).
\end{rubric}

\mattinsOfficeHymn
\hymnLett{P}{eter,} good shepherd, may thy ceaseless orisons,

\hymnSecondLine{For us prevailing, break the bands of wickedness:}

\begin{hangparas}{1.25em}{1}
For thou of old time didst receive authority

The gates to open, or to close, of Paradise.\\

2. O by thy doctrine, Paul, thou sage illustrious,

Guide us in virtue, raise our spirits heavenwards;

Till perfect knowledge stream on us abundantly,

And that which only is in part be done away.\\

3. Glory eternal to the blessed Trinity,

With laud and honour, virtue and supremacy,

Trinal yet Onely, reigning in his majesty

Both now and ever, through the ages infinite. Amen.\\
\end{hangparas}

℣. They declared the work of God.

℟. And wisely considered of his doing.

\mattinsAntiphon{Whatsoever {\dag} thou shalt bind on earth shall be bound in heaven: and whatsoever thou shalt loose on earth shall be loosed in heaven: saith the Lord unto Simon Peter.}

\collect
\lett{O}{Almighty} God, who by thy Son Jesus Christ didst give to thy Apostle Saint Peter many excellent gifts, and commandedst him earnestly to feed thy flock; Make, we beseech thee, all Bishops and Pastors diligently to preach thy holy Word, and the people obediently to follow the same, that they may receive the crown of everlasting glory. Through the same.

\lett{O}{God,} who hast hallowed this day by the martyrdom of thine Apostles Peter and Paul: grant unto thy Church in all things to follow the commandment of those; through whom she received the beginning of religion. Through.

\fussy	
\end{paracol}



%MANUAL ADJUSTMENT:
\clearpage
\festival{Greater Double}{30 June}{Commemoration of St. Paul}{A.}{FestiveWhite}
\feastday{\small St. Paul Commemoration}
\feastdate{30}{June}

\firstEvensong
\psalmAntiphon{I have planted, {\dag} Apollos watered, but God gave the increase, alleluia.}

\firstEvensongHymn

\begin{multicols}{2}
\sloppy

\hymnLett{O}{by} thy doctrine, Paul, thou sage illustrious,

\hymnSecondLine{Guide us in virtue, raise our spirits heavenwards;}

\begin{hangparas}{1.25em}{1}
Till perfect knowledge stream on us abundantly,

And that which only is in part be done away.

\newcolumn

2. Glory eternal to the blessed Trinity,

With laud and honour, virtue and supremacy,

Trinal yet Onely, reigning in his majesty

Both now and ever, through the ages infinite. Amen.\\

\end{hangparas}

\fussy	
\end{multicols}

℣. Thou art a chosen vessel, holy Apostle Paul.

℟. A preacher of the truth throughout all the world.

\firstEvensongAntiphon{O holy Apostle Paul, {\dag} thou preacher of the truth and Doctor of the Gentiles, intercede for us unto God, who hath chosen thee.}

\mattins
\psalmAntiphon{Most gladly will I glory {\dag} in mine infirmities, that the power of Christ may rest upon me.}

\begin{rubric}
	Invitatory Hymn of I Evensong.
\end{rubric}

\begin{rubric}
	Office Hymn of the Common of the Apostles (p. \pageref{CommonApostles}).
\end{rubric}

℣. Thou art a chosen vessel, holy Apostle Paul.

℟. A preacher of the truth throughout all the world.

\mattinsAntiphon{Ye which have followed me {\dag} shall sit upon twelve thrones, judging the twelve tribes of Israel, saith the Lord.}

\secondEvensong
\psalmAntiphon{The grace of God {\dag} which was bestowed upon me was not in vain: but his grace abideth with me always.}

\begin{rubric}
	Office Hymn, Versicle, \& Canticle Antiphon of I Evensong.
\end{rubric}

\collect
\lett{O}{God,} who by the preaching of the blessed Apostle Paul didst teach the multitude of the Gentiles: grant to us, we beseech thee; that we who celebrate his birthday (commemoration) may know him to be our advocate with thee. Through.

\begin{rubric}
	Commemoration \blackrubric{of St. Peter} is made before all other Commemorations, as followeth.
\end{rubric}

\lett{O}{God,} who didst bestow upon thy blessed Apostle Peter the keys of the kingdom of heaven, and didst  appoint unto him the high priesthood of binding and loosing: vouchsafe; that by the help of his intercession we may be delivered from the bonds of our iniquities. (Who livest.)

\begin{rubric}
	Commemoration \blackrubric{of the Octave of St. John Baptist} (p. \pageref{JohnBaptistCollect}).
\end{rubric}


\festivalLL{Second Class Double}{1 July}{The Most Precious Blood of Our Lord Jesus Christ}
\feastday{Precious Blood}
\feastdate{1}{July}

\openingSentence{We therefore pray thee, help thy servants, whom thou hast
redeemed with thy precious blood.}


\firstEvensong
\psalmAntiphon{He was clothed {\dag} with a vesture dipped in blood: and his Name is called The Word of God.}

\firstEvensongHymn

\begin{multicols}{2}
\sloppy

\hymnLett{W}{ith} glad and joyous strains, now let each street resound,

\hymnSecondLine{And let the laurel wreath each Christian brow entwine;}

\begin{hangparas}{1.25em}{1}
With torches waving bright, let old and young go forth,

And swell the train in solemn line.\\

2. Whilst we with bitter tears, with sighs and grief profound,

Wail o'er the saving Blood, pour'd forth upon the Tree,

Oh, deeply let us muse, and count the heavy price

Which Christ hath paid to make us free.\\

3. The primal man of old, who fell by serpent's guile,

Brought death and many woes upon his fallen race;

But our New Adam, Christ, new life unto us gave,

And brought to all ne'er-ending grace.\\

4. To heaven's highest height, the wailing cry went up

Of him who hung in pain, God's own eternal Son;

His saving, priceless Blood, his Father's wrath appeas'd,

And for his sons full pardon won.\\

5. Who-e'er in that pure Blood his guilty soul shall wash,

Shall from his stains be freed, be made as roses bright;

Shall vie with Angels pure, shall please his King and Lord,

And precious shine in his glad sight.\\

6. Oh, from the path of right ne'er let thy steps depart,

But haste thee to the goal in virtue's peaceful ways;

Thy God who reigns on high will e'er direct thy steps,

And crown thy deeds with blissful days.\\

7. Father of all things made, to us propitious be,

For whom thy own dear Son his saving Blood did spill;

O Holy Spirit, grant the souls by thee refresh'd

Eternal bliss may ever fill. Amen.\\
\end{hangparas}

℣. Thou hast redeemed us, O Lord, by thy Blood.

℟. And hast made us unto our God a kingdom.

\firstEvensongAntiphon{But ye are come {\dag} unto mount Sion and unto the city of the living God, the heavenly Jerusalem, and to Jesus, the mediator of the new covenant, and to the Blood of sprinkling, that speaketh better things than that of Abel.}

\fussy
\end{multicols}


\mattins
\psalmAntiphon{And they overcame the dragon {\dag} by the Blood of the Lamb and the word of their testimony.}

\begin{paracol}{2}
\sloppy

\mattinsInvitatoryHymn
\hymnLett{R}{ighteous} anger of our Maker

\hymnSecondLine{Pour'd his vengeance in the Flood,}

\begin{hangparas}{1.25em}{1}
And the crime-filled world submerged;

Noah sav'd in ark of wood.

Then he came with love stupendous,

Wash'd the world in Precious Blood.\\

2. So the happy earth and healthful,

Water'd by such saving show'rs,

Though once sown with thorns and trouble,

Now sprouts forth with tiny flow'rs:
All the bitter plants of wormwood

Change to flowing nectar's dow'rs.\\

3. Straightway ceas'd the serpent's venom,

All its poison laid aside;

And the wild beasts' raging ended

While their cruel fierceness died;

For the gentle Lamb and wounded

Won the vict'ry glorifi'd.\\

4. O supernal, wondrous Knowledge

High exalted, heights unknown!

O how sweet thy loving-kindness

To the weak and sinful shown,

That the King of kings, all Goodness,

Should for guilty slaves atone.\\

5. When we rouse the Judge's anger

By our sinfulness abhorr'd,

Then from wrath we are protected

By thy pleading Blood outpour'd,

Then from violent waves of evil

We to safety are restor'd.\\

6. Thee, O Christ, the world's Redeemer,

With our thankful praise we greet,

Ever dwelling with the Father,

And the Holy Paraclete,

Whose is boundless praise and power

Through the ages; as is meet. Amen.
\end{hangparas}

\switchcolumn

\mattinsOfficeHymn
\hymnLett{H}{ail,} holy wounds of Jesus, hail,

\hymnSecondLine{Sweet pledges of the saving Rood,}

\begin{hangparas}{1.25em}{1}
Whence flow the streams that never fail,

The purple streams of his dear Blood.\\

2. Brighter than brightest stars ye shew,

Than sweetest rose your scent more rare,

No Indian gem may match your glow,

No honey's taste with yours compare.\\

3. Portals ye are to that dear home

Wherein our wearied souls may hide,

Whereto no angry foe can come,

The Heart of Jesus crucifi'd.\\

4. What countless stripes our Jesus bore,

All naked left in Pilate’s hall!

From his torn flesh how red a show'r

Did round his sacred person fall!\\

5. His beauteous brow, oh, shame and grief,

By the sharp thorny crown is riv'n;

Through hands and feet, without relief,

The cruel nails are rudely driv'n.\\

6. But when for our poor sakes he died,

A willing Priest by love subdued,

The soldier’s lance transfix'd his side,

Forth flow'd the Water and the Blood.\\

7. In full atonement of our guilt,

Careless of self, the Saviour trod,

E’en till his Heart’s best Blood was spilt,

The wine-press of the wrath of God.\\

8. Come, bathe you in the healing flood,

All ye who mourn, by sin oppress'd,

Your only hope is Jesus’ Blood,

His sacred Heart your only rest.\\

9. All praise to him, th' Eternal Son,

At God’s right hand enthroned above,

Whose Blood our full redemption won,

Whose Spirit seals the gift of love. Amen.\\
\end{hangparas}

℣. Being justified by the Blood of Christ.

℟. We shall be saved from wrath through him.

\mattinsAntiphon{The Blood of the Lamb {\dag} shall be to you for a token, saith the Lord: and when I see the Blood, I will pass over you, and the plague shall not be upon you to destroy you.}

\fussy	
\end{paracol}

\secondEvensong
\psalmAntiphon{Blessed {\dag} are they which have washed their robes in the Blood of the Lamb.}

%\secondEvensongHymn
\begin{rubric}
	Office Hymn \& Versicle of I Evensong.
\end{rubric}

\secondEvensongAntiphon{And this day {\dag} shall be unto you for a memorial: and ye shall keep it a feast to the Lord throughout your generations by an ordinance for ever.}

\collect
\lett{A}{lmighty} and everlasting God, who didst appoint thine only-begotten Son to be the Redeemer of the world, and hast vouchsafed to be reconciled to us by his Blood: grant us, we beseech thee, so to venerate with solemn worship the price of our salvation, that by its power we may be defended from the evils of this present life on earth; and may rejoice in the everlasting fruit thereof in heaven. Through.
\begin{rubric}
    Commemoration \blackrubric{of the Octave Day of St. John Baptist} (p. \pageref{JohnBaptistCollect}).
\end{rubric}


%MANUAL ADJUSTMENT:
\clearpage
\festivalLL{Second Class Double}{2 July}{Visitation of the Blessed Virgin Mary}
\feastday{Visitation B.V.M.}
\feastdate{2}{July}

\begin{rubric}
	Common of the Blessed Virgin Mary (p. \pageref{CommonBVM}).
\end{rubric}

\firstEvensong
\psalmAntiphon{Mary arose, {\dag} and went into the hill country with haste, into the city of Judah.}\\

℣. Blessed art thou amongst women.

℟. And blessed is the fruit of thy womb.

\firstEvensongAntiphon{Blessed art thou, {\dag} O Mary, for thou hast believed: and there shall be a performance in thee of those things which were told thee from the Lord, alleluia.}

\mattins
\psalmAntiphon{When Elisabeth heard {\dag} the salutation of Mary, the babe leaped in her womb, and she was filled with the Holy Ghost, alleluia.}\\

℣. Blessed art thou amongst women.

℟. And blessed is the fruit of thy womb.

\mattinsAntiphon{When Elisabeth {\dag} heard the salutation of Mary, she spake out with a loud voice, and said: Whence is this to me, that the mother of my Lord should come to me? Alleluia.}

\secondEvensong
\psalmAntiphon{For lo, as soon {\dag} as the voice of my salutation sounded in mine ears, the babe leaped in my womb for joy, alleluia.}\\

℣. Blessed art thou amongst women.

℟. And blessed is the fruit of thy womb.

\secondEvensongAntiphon{All generations {\dag} shall call me blessed: for God hath regarded the lowliness of his handmaiden, alleluia.}

%MANUAL ADJUSTMENT:
\clearpage
\collect
\lett{W}{e} beseech thee, O Lord, to grant unto us thy servants the gift of heavenly grace: that as the child-bearing of the blessed Virgin was unto us the beginning of salvation; so the devout observance of her Visitation may bestow on us an increase of peace. Through.

\begin{rubric}
	 Commemoration \blackrubric{of St. John of San Francisco,} from the Common of a Confessor Bishop (p. \pageref{ConfessorBishopCollect}), the second Collect.
\end{rubric}

\begin{rubric}
	 Commemoration \blackrubric{of Sts. Processus and Martinian} (p. \pageref{ProcessusCollect}).
\end{rubric}


\subbymemorial{2 July}{St. John of San Francisco}{B.C.}{FestiveWhite}
\feastday{{St. John San Francisco}}
\feastdate{2}{July}

\begin{rubric}
	Common of a Confessor Bishop (p. \pageref{ConfessorBishopCollect}), the second Collect.
\end{rubric}


\subbymemorial{2 July}{Sts. Processus and Martinian}{M.}{MartyrRed}
\feastday{{Sts. Processus \&c.}}
\feastdate{2}{July}

\memorialCollect\label{ProcessusCollect}
\lett{O}{God,} who dost encompass and protect us by the glorious confession of thy holy Martyrs Processus and Martinian: grant us both to profit by their example, and to rejoice in their intercession. Through.


\festival{Double}{3 July}{St. Iren{\ae}us of Lyon}{B.M.}{MartyrRed}
\feastday{{St. Iren{\ae}us Lyon}}
\feastdate{3}{July}

\begin{rubric}
	Common of a Martyr (p. \pageref{CommonMartyr}).
\end{rubric}

%\collect
\lett{O}{God,} who gavest grace to blessed Iren{\ae}us thy Martyr and Bishop to overcome false doctrine by the truth of his teaching, and favourably to establish peace in thy Church: give thy people, we beseech thee, constancy in holy religion, and grant us thy peace all the days of our life. Through.


\festival{Greater Double}{6 July}{Octave Day of St. Peter}{A.}{FestiveWhite}
\feastday{{Sts. Peter \& Paul Octave}}
\feastdate{6}{July}

\begin{rubric}
	Hymns, Versicles, \& Antiphons of the Feast (p. \pageref{SaintPeter}).
\end{rubric}

%\collect
\lett{O}{God,} whose right hand upheld blessed Peter lest he should sink when walking on the waves, and thrice from the depths of the sea delivered his fellow-Apostle Paul when suffering shipwreck: graciously hear us and grant; that, through the merits of them both, we may obtain the glory of everlasting life: Who livest.


%Last Day of Eastertide Possible---

%MANUAL ADJUSTMENT:
\clearpage
\festival{Double (m.t.v.)}{7 July}{Sts. Cyril and Methodius}{B.C.}{FestiveWhite}
\feastday{{Sts. Cyril \& Methodius}}
\feastdate{7}{July}

\begin{rubric}
	Common of a Confessor Bishop (p. \pageref{CommonConfessorBishop}).
\end{rubric}

%\collect
\lett{A}{lmighty} and everlasting God, who by thy blessed Confessors and Bishops, Cyril and Methodius, didst suffer the peoples of Slavonia to come to the knowledge of thy name: vouchsafe; that we, who glory in their festival may be joined unto their fellowship. Through.


\subbymemorial{10 July}{Seven Holy Brothers, with Sts. Rufina and Secunda}{M.}{MartyrRed}
\feastday{{7 Holy Brothers}}
\feastdate{10}{July}

\memorialCollect
\lett{G}{rant,} we beseech thee, Almighty God: that, like as we have known thy glorious Martyrs to be constant in their confession, so we may perceive their loving intercession. Through.


\subbymemorial{10 July}{St. Joseph of Damascus and Companions}{M.}{MartyrRed}
\feastday{{Sts. Joseph Damascene \&c.}}
\feastdate{10}{July}

\begin{rubric}
	Common of Many Martyrs (p. \pageref{CommonMartyrsNotBishopsCollect}), the second Collect.
\end{rubric}


\subbymemorial{11 July}{Pope St. Pius I of Rome}{B.M.}{MartyrRed}
\feastday{{Pope St. Pius I}}
\feastdate{11}{July}

\begin{rubric}
	Common of a Martyr (p. \pageref{CommonMartyrBishopCollect}).
\end{rubric}


\subbymemorial{12 July}{Sts. Nabor and Felix}{M.}{MartyrRed}
\feastday{{Sts. Nabor \& Felix}}
\feastdate{12}{July}

\memorialCollect
\lett{G}{rant,} we beseech thee, O Lord: that as the birthday of thy holy Martyrs, Nabor and Felix, faileth not to return for our observance; so they may continually assist us by their prayers. Through.


\subbymemorial{13 July}{Pope St. Anacletus of Rome}{B.M.}{MartyrRed}
\feastday{{St. Anacletus}}
\feastdate{13}{July}

\begin{rubric}
	Common of a Martyr (p. \pageref{CommonMartyrBishopCollect}), the second Collect.
\end{rubric}


\festival{Double}{15 July}{King St. Vladimir of Kiev}{C.}{FestiveWhite}
\feastday{{St. Vladimir Kiev}}
\feastdate{15}{July}

\begin{rubric}
	Common of a Confessor not a Bishop (p. \pageref{CommonConfessorNotBishop}).
\end{rubric}


%MANUAL ADJUSTMENT:
\clearpage
\subbymemorial{16 July}{Sts. Nicholas and Habib Khasha}{M.}{MartyrRed}
\feastday{Sts. Nicholas \& Habib}
\feastdate{16}{July}

%%CHECK:
\begin{rubric}
	Common of Many Martyrs (p. \pageref{CommonMartyrsNotBishopsCollect}), the second Collect.
\end{rubric}


\subbymemorial{17 July}{St. Alexius}{C.}{FestiveWhite}
\feastday{{St. Alexius}}
\feastdate{17}{July}

\begin{rubric}
	Common of a Confessor not a Bishop (p. \pageref{ConfessorNotBishopCollect}).
\end{rubric}


\festival{Double (m.t.v.)}{18 July}{Translation of St. Raphael of Brooklyn}{B.C.}{FestiveWhite}
\feastday{{St. Raphael Brooklyn Translation}}
\feastdate{18}{July}

\begin{rubric}
	Common of a Confessor Bishop (p. \pageref{CommonConfessorBishop}).
\end{rubric}


\subbymemorial{18 July (m.t.v.)}{St. Sergius of Radonezh}{Ab.}{FestiveWhite}
\feastday{{St. Sergius Radonezh}}
\feastdate{18}{July}

\begin{rubric}
	Common of a Confessor not a Bishop (p. \pageref{AbbotCollect}).
\end{rubric}


\subbymemorial{18 July}{Sts. Symphorosa and Her Seven Sons}{M.}{MartyrRed}
\feastday{{Sts. Symphorosa \&c.}}
\feastdate{18}{July}

\memorialCollect
\lett{O}{God,} who vouchsafest unto us to keep the heavenly birthday of thy holy Martyrs Symphorosa and her sons: grant, we beseech thee; that we may rejoice in the perpetual felicity of their friendship. Through.


\subbymemorial{19 July (m.t.v.)}{St. Seraphim of Sarov}{C.}{FestiveWhite}
\feastday{{St. Seraphim Sarov}}
\feastdate{19}{July}

\begin{rubric}
	Common of a Confessor not a Bishop (p. \pageref{ConfessorNotBishopCollect}).
\end{rubric}


\festival{Double}{20 July}{St. Elias the Prophet}{C.}{FestiveWhite}
\feastday{{St. Elias Prophet}}
\feastdate{20}{July}

\openingSentence{Then stood up Elias the prophet as fire, and his word burned like a lamp, alleluia.}

\begin{paracol}{2}
\sloppy

\firstEvensong
\psalmAntiphon{With zeal have I been zealous {\dag} for the Lord God of hosts: for the children of Israel have forsaken thy covenant.}

\firstEvensongHymn
\hymnLett{T}{he} lofty peaks of Carmel

\hymnSecondLine{With tuneful praises ring,}

\hymnThirdLine{The anthems of Elias}

\begin{hangparas}{1.25em}{1}

`Tis our delight to sing.\\

%Order changed to Churches:
2. The glory of our Churches,

Our leader, prop, and stay,

From east to west his offspring

Increaseth day by day.\\

3. When sorely press'd with famine,

A raven serv'd him bread,

With meal and cruse unfailing,

The widow'd hearth was fed.\\

4. The boy from death delivered

Is to his home restor'd,

And light so much desir'd,

In radiant flood is pour'd.\\

5. Behold the Heaven closeth,

To open at his voice,

And copious welcome showers

The thirsty lands rejoice.\\

6. To Father, Son, and Spirit,

Be equal power and praise,

All glory and dominion

Henceforth for endless days. Amen.\\
\end{hangparas}

    ℣. By the word of the Lord he shut up the heaven.

	℟. And he brought down fire from heaven thrice.

\firstEvensongAntiphon{Behold, I {\dag} will send you Elias the prophet before the coming of the great and dreadful day of the Lord: And he shall turn the heart of the fathers to the children, and the heart of the children to their fathers}

\switchcolumn

\secondEvensong
\psalmAntiphon{I have not troubled Israel, {\dag} but thou and thy father’s house, in that ye have forsaken the commandments of the Lord.}

\secondEvensongHymn
\hymnLett{T}{hou} prop of our Churches, thou pride of our race,

\hymnSecondLine{Let thy praises resound far and near,}

\begin{hangparas}{1.25em}{1}
Let sea and the land and the air give them place,

Rehearsing in gladness thy glory and grace,

Till the earth and the heav'ns give ear.\\

2. O sun of the heavens, how lovely thy rays!

What power thy wonders unfold,

How fruitful in merits the length of thy days,

Commission'd by God in His manifold ways,

For noble endeavours of old!\\

3. To regions celestial, in power and might,

Triumphant thy chariot speeds;

Uplifted by angels to marvellous height,

While shining in splendour and dazzling with light,

Thou guidest the fiery steeds.\\

4. As witness to men of his Sonship divine,

With Jesus thy glory we view;

The Father hath call'd thee on Tabor to shine,

Companion to Moses, and with him to sign

A testament faithful and true.\\

5. Protect us we pray, ‘neath thy powerful shield,

Incline to our aid from above,

Let thy fost'ring guidance be ever reveal'd

To thy children of Carmel, whose bosoms are seal'd

With the strength of thy fath'rly love\\

6. All power, dominion, all glory and praise,

Be given to Father and Son,

To thee, Holy Spirit, for numberless days

Our homage eternal we equally raise

All glory to God, Three in One. Amen.\\
\end{hangparas}

    ℣. Blessed are they that saw thee.

	℟. And were honoured with thy friendship.
	
%Based on KJV with changes based off the Latin.
\secondEvensongAntiphon{And Elias took {\dag} his mantle, and wrapped it together, and smote the waters of the Jordan, and they were divided hither and thither, so that he and Eliseus went over on dry ground. As they still went on, behold, there appeared a chariot of fire, and horses of fire, and parted them both asunder; and Elias went up by a whirlwind into heaven, and Eliseus saw him no more.}

\fussy
\end{paracol}

\mattins
\psalmAntiphon{How long halt ye {\dag} between two opinions? if the Lord be God, follow him.}

\begin{paracol}{2}
\sloppy

\mattinsInvitatoryHymn

\hymnLett{C}{reator} of the universe, our grateful hearts rejoice,

\hymnSecondLine{Exalting thee in this, the mighty Thesbite of thy choice;}

\begin{hangparas}{1.25em}{1}
2. Whose zeal for thy great glory, enkindled as a flame,

Defied the impious Prophets, and slew them in thy name.\\

3. His holocaust ascendeth, for thy fire may never fail,

While scorn and deep derision mock the clamorous priests of Baal.\\

4. The baneful rage of Jezabel he flieth in his dread,

And sleeping `neath a juniper is by an angel fed.\\

5. Empowered by the vision, with new courage he hath trod

Unto the heights of Horeb, unto the Mount of God.\\

6. O virtue of this bread divine, imparted by the Lord!

Full forty days he fasteth in the strength it doth afford.\\

7. All glory be to thee, O Father, Son, and Paraclete,

O undivided Trinity, thy praise all hearts repeat. Amen.
\end{hangparas}

\switchcolumn

\mattinsOfficeHymn
\hymnLett{C}{ome,} blest companions, let our joy resounding

\hymnSecondLine{Extol to heav'n the leader of our line.}

\begin{hangparas}{1.25em}{1}
`Tis meet the mem'ry of his deeds abounding

Should waken ceaseless canticles divine.\\

2. He knows the gentle breathing of the Spirit

Cloth'd in the whistling murmur of the air,

By God's command the chastisements they merit

Proud Jezebel and Ahab justly share.\\

3. The caverns green of Carmel form his dwelling,

With leathern tunic is he rudely clad,

To impious Ahaziah his foretelling

Gives portent of a dissolution sad.\\

4. Twice at his pray'r the fire from Heav'n descending

Consumeth trembling soldiers in its flame,

The flowing wat'rs met with his mantle rending,

Dry shod he passeth safely through the same.\\

5. O Father, let thy help and thy protection

Be o'er thy children as they humbly plead,

Entreat the Spirit, by his sweet election,

To multiply his graces in their need.\\

6. O unbegotten Father, we adore thee,

O Son begotten, rev'rence be to thee,

O glorious Spirit, bow we low before thee,

Thou simple undivided Trinity. Amen.\\
\end{hangparas}

    ℣. Elias was covered with the whirlwind.

	℟. And his spirit was filled up in Eliseus.

\antiphon{Ben.}{Elias was a man {\dag} subject to like passions as we are, and he prayed earnestly that it might not rain: and it rained not on the earth by the space of three years and six months. And he prayed again, and the heaven gave rain, and the earth brought forth her fruit.}\par\noindent

\fussy	
\end{paracol}


\collect
\lett{G}{rant,} we beseech thee, Almighty God: that we who believe that thou didst marvellously lift up thy Prophet Elias in a fiery chariot, while yet in this life; may at his intercession, while still alive, be raised to spiritual heights and rejoice in the resurrection of the just. Through.

%%CHECK:
\begin{rubric}
	Commemoration \blackrubric{of St. Margaret of Antioch,} from the Common of Virgins (p. \pageref{VirginMartyressCollect}), the second Collect.
\end{rubric}


\subbymemorial{20 July}{St. Margaret of Antioch}{V.M.}{MartyrRed}
\feastday{{St. Margaret Antioch}}
\feastdate{20}{July}

%%CHECK:
\begin{rubric}
	Common of Virgins (p. \pageref{VirginMartyressCollect}), the second Collect.
\end{rubric}


\subbymemorial{21 July}{St. Praxedes}{V.}{FestiveWhite}
\feastday{{St. Praxedes}}
\feastdate{21}{July}

%%CHECK:
\begin{rubric}
	Common of Virgins (p. \pageref{VirginOnlyCollect}), the second Collect.
\end{rubric}


%MANUAL ADJUSTMENT:
\clearpage
%Sarum Psalm Antiphons:
\festivalAO{Greater Double}{22 July}{St. Mary Magdalene, Penitent}
\feastday{St. Mary Magdalene}
\feastdate{22}{July}

\begin{paracol}{2}
\sloppy


%MANUAL ADJUSTMENT:
\vspace{-1\baselineskip}


\firstEvensong
\psalmAntiphon{As Jesus was reclining in the house of Simon the Pharisee, {\dag} Mary Magdalene approached him, bringing a pound of precious ointment.}
%Recumbente * Jesu in domo pharisei Symonis, accedit ad eum Maria Magdalene afferens preciosi libram unguenti.

\firstEvensongHymn
\hymnLett{O}{Father} of celestial rays,

\hymnSecondLine{When thou on Magdalene dost gaze,}

\begin{hangparas}{1.25em}{1}
The flame of burning love appears,

Her icy heart dissolves in tears.\\

2. Wounded by love, she hastens o'er

The feet of Christ her tears to pour,

Anoints them, wipes them with her hair,

And prints adoring kisses there.\\

3. Fearless, the Cross she will not leave:

And grieving, to the Tomb doth cleave:

No ruthless soldiers cause her dread:

For from her love all fear hath fled.\\

4. O Christ, true Charity thou art;

Purge thou the foulness of our heart,

Fill every soul with grace and love,

And give us thy rewards above.\\

5. All laud to God the Father be;

All praise, eternal Son, to thee;

All glory, as is ever meet,

To God the Holy Paraclete. Amen.\\
\end{hangparas}

℣. Mary hath chosen that good part.

℟. Which shall not be taken away from her.

\firstEvensongAntiphon{A woman {\dag} in the city which was a sinner, when she knew that Jesus sat at meat in the house of Simon the leper, brought an alabaster box of ointment, and stood behind at the feet of Jesus, and began to wash his feet with tears, and did wipe them with the hairs of her head, and kissed his feet, and anointed them with the ointment.}

\secondEvensong
\psalmAntiphon{O Mary Magdalene, {\dag} continually and humbly beseeching the Lord Jesus, intercede for us!}
%Fifth Lauds Psalm Antiphon
%Intercede *supplicans assidue pro nobis Jesum Dominum Maria Magdalene.

\begin{rubric}
	Office Hymn of I Evensong.
\end{rubric}

℣. God hath chosen her and preferred her.

℟. He hath made her to dwell in his tabernacle.

\secondEvensongAntiphon{A woman {\dag} in the city which was a sinner, brought an alabaster box of ointment, and stood at the Lord’s feet, and began to wash his feet with tears, and did wipe them with the hairs of her head.}

\switchcolumn

%MANUAL ADJUSTMENT:
\vspace{-1\baselineskip}

\mattins
\psalmAntiphon{Mary, {\dag} concealing her pure heart, anointed the Lord, cleansing herself by the holy stream of baptism.}
%Second Lauds Psalm Antiphon:
%Pectore * sincero Dominum Maria recondens unxit purgantem se baptismi gurgite sancto.

\mattinsInvitatoryHymn
\hymnLett{T}{he} Lord's dear feet did Mary lave

\hymnSecondLine{With her own tears; her hair she gave}

\begin{hangparas}{1.25em}{1}
To wipe them dry; and then she pour'd

The precious ointment, and ador'd.\\

2. To God the Father, God the Son,

And God the Spirit, Three in One,

Praise, honour, and dominion be

From age to age eternally. Amen.
\end{hangparas}

\mattinsOfficeHymn
\hymnLett{T}{hou} only Son of God on high,

\hymnSecondLine{Regard us with a gracious eye,}

\hymnThirdLine{Who weeping Magdalene dost own}

\begin{hangparas}{1.25em}{1}
And call unto thy glorious throne.\\

2. Lo! in the royal coffers laid

Again the long lost coin display'd;

The noble gem of sparkling sheen,

From mire recover'd, glows serene.\\

3. Jesu, our refuge sure and sweet,

Thee, hope of penitents, we greet;

Forgive the hearts that fain would break,

For that repentant sinner's sake.\\

4. And may that Mother kind and meek

Think on our nature frail and weak,

And raise her prayer that we may gain

A passage safe o'er life's rough main.\\

5. To God alone be honour paid

For grace so manifold displayed:

Their guilt he pardons who repent,

And gives reward for punishment. Amen.\\
\end{hangparas}

℣. Her sins, which are many, are forgiven.

℟. For she loved much.

\mattinsAntiphon{Mary therefore {\dag} anointed the feet of Jesus, and wiped them with her hair, and the house was filled with the odour of the ointment.}

\collect
\lett{G}{rant} us, most merciful Father, like as Saint Mary Magdalene, by loving thy Only Begotten Son above all things, obtained forgiveness of her sins, so she may procure for us in thy compassionate presence everlasting blessedness. Through.

\fussy	
\end{paracol}



\subbymemorial{23 July}{St. John Cassian}{Ab.}{FestiveWhite}
\feastday{{St. John Cassian}}
\feastdate{23}{July}

\begin{rubric}
	Common of a Confessor not a Bishop (p. \pageref{AbbotCollect}).
\end{rubric}


\subbymemorial{23 July}{St. Apollinaris of Ravenna}{B.M.}{MartyrRed}
\feastday{{St. Apollinaris Ravenna}}
\feastdate{23}{July}

\memorialCollect
\lett{O}{God,} the rewarder of faithful souls, who hast hallowed this day by the martyrdom of blessed Apollinaris, thy Priest: grant, we beseech thee, unto us thy servants; that we who keep his solemn festival may by his prayers obtain thy pardon. Through.


%MANUAL ADJUSTMENT:
\clearpage
\subbymemorial{23 July}{St. Liborius}{B.C.}{FestiveWhite}
\feastday{{St. Liborius}}
\feastdate{23}{July}

\begin{rubric}
	Common of a Confessor Bishop (p. \pageref{ConfessorBishopCollect}).
\end{rubric}


\festival{Vigil}{24 July}{Vigil of St. James}{A.}{LentenViolet}
\feastday{{St. James Vigil}}
\feastdate{23}{July}

%%CHECK:
\begin{rubric}
	Common of Vigils of the Apostles (p. \pageref{CommonVigilApostles}), Commemoration \blackrubric{of St. Christina,} from the Common of Virgins (p. \pageref{CommonVirgin}), the second Collect.
\end{rubric}


\subbymemorial{24 July}{St. Christina}{V.M.}{MartyrRed}
\feastday{{St. Christina}}
\feastdate{24}{July}

%%CHECK:
\begin{rubric}
	Common of Virgins (p. \pageref{VirginMartyressCollect}), the second Collect.
\end{rubric}


\festival{Second Class Double}{25 July}{St. James}{A.}{FestiveWhite}
\feastday{St. James}
\feastdate{25}{July}

\begin{rubric}
	Common of the Apostles (p. \pageref{CommonApostles}).
\end{rubric}

%\collect
\lett{B}{e} thou, O Lord, the sanctifier and guardian of thy people, that under the protection of thy Apostle, James, they may please thee in their conversation, and serve thee in all quietness. Through.
\begin{rubric}
	 Commemoration \blackrubric{of St. Christopher,} from the Common of a Martyr (p. \pageref{CommonMartyrNotBishopCollect}).
\end{rubric}


\subbymemorial{25 July}{St. Christopher}{M.}{MartyrRed}
\feastday{{St. Christopher}}
\feastdate{25}{July}

\begin{rubric}
	Common of a Martyr (p. \pageref{CommonMartyrNotBishopCollect}).
\end{rubric}


%MANUAL ADJUSTMENT:
\clearpage
\festivalAO{Second Class Double}{26 July}{St. Anne, Mother of the Blessed Virgin Mary}
\feastday{St. Anne}
\feastdate{26}{July}

\begin{paracol}{2}
\sloppy

\firstEvensong
\psalmAntiphon{I will give them an everlasting name, saith the Lord. They shall obtain joy and gladness for ever.}
%Nomen sempiternum * dabo Sanctis meis, dicit Dominus, gaudium et laetitiam obtinebunt in aeternum.

\firstEvensongHymn
\hymnLett{M}{other} Anne, be joyful;

\hymnSecondLine{Sing, O mother holy,}

\hymnThirdLine{Since thou art the parent}

\begin{hangparas}{1.25em}{1}
Of God's Mother lowly.\\

2. Praise thy wondrous daughter;

Joachim, too, raises

To the Virgin Mary

His paternal praises.\\

3. For in her our planet

First hath benediction

Which in hapless Eva

Suffer'd malediction.\\

4. Therefore take the praises

Joyous hearts are paying;

And from all defilement

Cleanse us by thy praying.\\

5. Father, Son eternal,

Holy Ghost supernal,

With one praise we bless thee,

Three in One confess thee. Amen.\\
\end{hangparas}

℣. Be glad, O ye righteous and rejoice in the Lord.

℟. And be joyful, all ye that are true of heart.

\firstEvensongAntiphon{The kingdom of heaven {\dag} is likened unto a merchantman seeking goodly pearls; who, when he had found one pearl of great price, went and sold all that he had, and bought it.}


\switchcolumn

\secondEvensong
\psalmAntiphon{The people {\dag} will tell of the wisdom of the Saints, and the whole congregation will shew forth their praise.}
%Sapientiam * Sanctorum narrant populi, et laudes eorum pronuntiat omnis ecclesia.

\secondEvensongHymn
\hymnLett{L}{et} all the Saints in concert sing,

\hymnSecondLine{The mother of that Maid to laud}

\begin{hangparas}{1.25em}{1}
By whose unstained Childbearing

Salvation came to men from God.\\

2. She sought her faithful progeny

From God the Father, Lord of light;

And merited most worthily

The pride of virgins, Mary bright.\\

3. To Joachim, esteemed the peer

Of any in his goodliness,

Anne brought forth Mary, Mother dear

Of Jesus, King of righteousness.\\

4. Let Jesse's noble stock efface

From Mother Eve her olden taints:

Anne bears her child, the child of grace,

The fairest blossom of the Saints!\\

5. All honour, laud, and glory be,

O Jesu, Virgin-born, to thee;

All glory, as is ever meet,

To Father and to Paraclete. Amen.\\
\end{hangparas}

℣. Let the Saints be joyful with glory.

℟. Let them rejoice in their beds.

\secondEvensongAntiphon{She stretched out {\dag} her hand to the poor; yea, she reacheth forth her hands to the needy; she eateth not the bread of idleness.}

\fussy	
\end{paracol}

\mattins
\psalmAntiphon{O ye holy spirits and souls of the just, sing a hymn unto God, alleluia.}
%Sancti spíritus + et ánimæ justórum, hymnum dícite Deo, alleluja. 

\begin{rubric}
	Invitatory Hymn of I Evensong.
\end{rubric}

\mattinsOfficeHymn

\begin{multicols}{2}
\sloppy

\hymnLett{T}{he} morning star succeeds to night,

\hymnSecondLine{The dawn soon follows, growing white}

\begin{hangparas}{1.25em}{1}
To herald in the sunrise bright

That floods the waking world with light.\\

2. Christ is the sun of righteousness,

The dawn, the Mother full of grace;

Bright Anne precedes her, like the star,

To drive the shades of law afar.\\

3. Lo, Anne; the very fruitful root,

The tree of healing, whence a shoot

To richest blossoming did spring,

And brought us Christ: to whom we sing,\\

4. All honour, laud, and glory be,

O Jesu, Virgin-born, to thee;

All glory, as is ever meet,

To Father and to Paraclete. Amen.\\
\end{hangparas}

℣. Let the Saints be joyful with glory.

℟. Let them rejoice in their beds.

\mattinsAntiphon{Give her {\dag} of the fruit of her hands; and let her own works praise her in the gates.}

\fussy	
\end{multicols}

\collect
\lett{O}{God,} who on blessed Anne didst vouchsafe to bestow grace, that she might be made worthy to become the mother of her who bore thine only-begotten Son: mercifully grant; that we, who celebrate her festival (commemoration) may be aided by her intercession with thee. Through the same.


%MANUAL ADJUSTMENT:
\clearpage
\subbymemorial{27 July}{St. Pantaleon}{M.}{MartyrRed}
\feastday{{St. Pantaleon}}
\feastdate{27}{July}

%%CHECK:
\begin{rubric}
	Common of a Martyr (p. \pageref{CommonMartyrNotBishopCollect}), the second Collect.
\end{rubric}


\subbymemorial{28 July}{Sts. Nazarius, Celsus, Pope Victor, and Innocent}{M.}{MartyrRed}
\feastday{{Sts. Nazarius \&c.}}
\feastdate{28}{July}

%\begin{rubric}
%	The Common of Many Martyrs (p. \pageref{CommonMartyrsNotBishopsCollect}).
%\end{rubric}

\memorialCollect
\lett{O}{Lord,} let the blessed confession of thy Saints, Nazarius, Celsus, Victor, and Innocent, strengthen us: and obtain for our frailty the succour of thy goodness. Through.


%MANUAL ADJUSTMENT:
\vspace{-0.25\baselineskip}

\festival{Double}{29 July}{St. Martha of Bethany}{V.}{FestiveWhite}
\feastday{{St. Martha Bethany}}
\feastdate{29}{July}

\begin{rubric}
	Common of Virgins (p. \pageref{VirginOnlyCollect}), with a Commemoration \blackrubric{of Sts. Felix, etc} (p. \pageref{FelixEtc}).
\end{rubric}


%MANUAL ADJUSTMENT:
\vspace{-0.25\baselineskip}

\subbymemorial{29 July}{Sts. Pope Felix II, Simplicius, Faustinus, and Beatrice}{M.}{MartyrRed}\label{FelixEtc}
\feastday{{Sts. Felix II \&c.}}
\feastdate{29}{July}

\memorialCollect
\lett{G}{rant,} we beseech thee, O Lord: that as thy Christian people rejoice in the temporal solemnity of thy Martyrs, Felix, Simplicius, Faustinus, and Beatrice, so they may eternally enjoy the same; that those things which they devoutly celebrate, they may effectually obtain. Through.


%MANUAL ADJUSTMENT:
\vspace{-0.25\baselineskip}

\subbymemorial{30 July}{Sts. Abdon and Sennen}{M.}{MartyrRed}
\feastday{{Sts. Abdon \& Sennen}}
\feastdate{30}{July}

\memorialCollect
\lett{O}{God,} who on thy Saints Abdon and Sennen didst bestow the bounteous gift of grace to attain unto thy glory: grant unto thy servants the remission of their sins; that, by the intercession of the merits of thy Saints, they may be found worthy to be delivered from all adversities. Through.


%MANUAL ADJUSTMENT:
\vspace{-0.25\baselineskip}

\festival{Greater Double}{1 August}{St. Peter in Chains}{A.}{FestiveWhite}
\feastday{Lammas}
\feastdate{1}{August}

\openingSentence{Thou art Peter, and upon this rock I will build my church.}

\begin{paracol}{2}
\sloppy

%MANUAL ADJUSTMENT:
\vspace{-1\baselineskip}

\firstEvensong
\psalmAntiphon{Herod the king {\dag} proceeded further to take Peter also: and when he apprehended him, he put him in prison, intending after Easter to bring him forth to his people.}

\firstEvensongHymn
\hymnLett{P}{eter} the blessed, bidden by Christ's messenger,

\hymnSecondLine{Breaketh his fetters, liberated wondrously;}

\begin{hangparas}{1.25em}{1}
Shepherd and teacher of  the Church’s fellowship,

The flock's preserver and the sheepfold's guardian,

Lo, he repelleth all the wolves’ ferocity.\\

2. Glory, O Father, be to thee eternally;

Grace and dominion to the Son adorable;

Unto the Spirit, honour and authority:

Praise to the holy undivided Trinity,

Praise everlasting, through the ages infinite. Amen.\\
\end{hangparas}

℣. Thou art Peter.

℟. And upon this rock I will build my Church.

\firstEvensongAntiphon{Thou art the shepherd of the sheep, {\dag} O chief of the Apostles: unto thee were given the keys of the kingdom of heaven.}

\secondEvensong
\psalmAntiphon{The Lord hath sent {\dag} his Angel, and hath delivered me out of the hand of Herod, alleluia.}

\secondEvensongHymn
\begin{rubric}
	Office Hymn \& Versicle of I Evensong.
\end{rubric}

\secondEvensongAntiphon{Loosen these earthly fetters {\dag} at the command of God, O Peter; and let the kingdom of heaven be opened unto the blessed.}


\switchcolumn


%MANUAL ADJUSTMENT:
\vspace{-1\baselineskip}


\mattins
\psalmAntiphon{Peter therefore {\dag} was kept in prison; but prayer was made without ceasing of the Church unto God for him.}

\mattinsInvitatoryHymn
\hymnLett{W}{hate\smash{'}er} thou bindest here on earth in punishment,

\hymnSecondLine{Likewise remaineth firmly bound in paradise;}

\begin{hangparas}{1.25em}{1}
What thou dost loosen here by thine authority

Over the heavens ever hath its liberty:

When this world endeth, thou shalt judge eternally.\\

2. Glory, O Father, be to thee eternally;

Grace and dominion to the Son adorable;

Unto the Spirit, honour and authority;

Praise to the holy undivided Trinity

Praise everlasting through the ages infinite. Amen.
\end{hangparas}

\mattinsOfficeHymn
\begin{rubric}
	Office Hymn \blackrubric{of the Feast of the Chair of St. Peter at Antioch} (p. \pageref{CathedraAntioch}).
\end{rubric}

℣. Thou art Peter.

℟. And upon this rock I will build my Church.

\mattinsAntiphon{Whatsoever {\dag} thou shalt bind on earth shall be bound in heaven: and whatsoever thou shalt loose on earth shall be loosed in heaven: saith the Lord unto Simon Peter.}

\fussy	
\end{paracol}


\collect
\lett{O}{God,} who didst deliver blessed Peter the Apostle from his chains, and didst cause him to depart unhurt: loose, we beseech thee, the chains of our sins; and graciously keep from us all evils. Through.

\begin{rubric}
	 Commemoration \blackrubric{of St. Paul,} as followeth,
\end{rubric}
\lett{O}{God,} who by the preaching of the blessed Apostle Paul didst teach the multitude of the Gentiles: grant to us, we beseech thee; that we who celebrate his commemoration may know him to be our advocate with thee. (Through.)

\begin{rubric}
	 Commemoration \blackrubric{of the Holy Maccabees} (p. \pageref{HolyMaccabeesCollect}).
\end{rubric}


\subbymemorial{1 August}{Holy Maccabees}{M.}{MartyrRed}
\feastday{{Holy Maccabees}}
\feastdate{1}{August}

\memorialCollect\label{HolyMaccabeesCollect}
\lett{O}{Lord,} let the crown of the brethren, thy Martyrs, cause us to rejoice: that we may thereby be strengthened and increased in our faith; and comforted by their manifold intercession. Through.


\subbymemorial{2 August}{Pope St. Stephen of Rome}{B.M.}{MartyrRed}
\feastday{{Pope St. Stephen}}
\feastdate{2}{August}

%%CHECK:
\begin{rubric}
	Common of a Martyr (p. \pageref{CommonMartyrBishopCollect}), the second Collect.
\end{rubric}


\festival{Double}{3 August}{Invention of St. Stephen}{M.}{MartyrRed}
\feastday{{Invention of St. Stephen}}
\feastdate{3}{August}

%%ADD REFERENCE:
\begin{rubric}
	Hymns, Versicles, \& Antiphons \blackrubric{of the Feast of St. Stephen} (p. \pageref{SaintStephen}).
\end{rubric}

%\collect
\lett{G}{rant} us, we beseech thee, O Lord, so to imitate that which we revere: that we may learn to love even our enemies; forasmuch as we celebrate the Finding of him, who was able to pray even for his persecutors to our Lord Jesus Christ thy Son. Who liveth.


\festivalLL{Greater Double}{5 August}{Dedication of Our Lady of the Snows}
\feastday{{Our Lady of the Snows}}
\feastdate{5}{August}

\begin{rubric}
	Common of the Blessed Virgin Mary (p. \pageref{CommonBVM}), Commemoration \blackrubric{of King St. Oswald} (p. \pageref{Oswald}).
\end{rubric}


%MANUAL ADJUSTMENT:
\clearpage
\festival{Double}{5 August}{King St. Oswald}{M.}{MartyrRed}
\feastday{{St. Oswald}}\label{Oswald}
\feastdate{5}{August}

\begin{rubric}
	Common of a Martyr (p. \pageref{CommonMartyr}).
\end{rubric}

%\collect
\lett{A}{lmighty} and everlasting God, who by the martyrdom of the blessed King Oswald hast hallowed this day with holy joy and gladness: grant unto our hearts the increase of thy charity; that we, who honour his glorious battle for the faith, may imitate his constancy even unto death. Through.


\festivalLL{Second Class Double}{6 August}{Transfiguration of Our Lord Jesus Christ}
\feastday{Transfiguration}
\feastdate{6}{August}

\begin{paracol}{2}
\sloppy

%MANUAL ADJUSTMENT:
\vspace{-1\baselineskip}

\firstEvensong
\psalmAntiphon{Jesus taketh {\dag} Peter, James, and John his brother, and bringeth them up into an high mountain apart, and was transfigured before them.}

\firstEvensongHymn
\hymnLett{A}{ll} ye who seek for Jesus, raise

\hymnSecondLine{Your eyes above, and upward gaze:}

\begin{hangparas}{1.25em}{1}
There may ye see the wondrous sign

Of never ending glory shine.\\

2. Behold him in celestial rays

Who never knoweth end of days;

Exalted, infinite, sublime;

Older than heav'n or hell or time.\\

3. This is the Gentiles’ King and Lord;

The Prince by Judah's race adored,

Promis'd to Abraham of yore

And to his seed for evermore.\\

4. To him the prophets testify;

And that same witness, from on high,

The Father seals by his decree:

Hear and believe my Son, saith he.\\

5. All glory, Lord, to thee we pay,

Transfigur'd on the mount to-day;

All glory, as is ever meet,

To Father and to Paraclete. Amen.\\
\end{hangparas}

℣. Glorious didst thou appear in the sight of the Lord.

℟. Because the Lord hath clothed thee with majesty.

\firstEvensongAntiphon{Christ Jesus, {\dag} the brightness of the Father and the express image of his person, who upholdeth all things by the word of his power, while he was by himself purging away our sins, vouchsafed on this day to shew himself in glory upon an high mountain.}

\switchcolumn

%MANUAL ADJUSTMENT:
\vspace{-1\baselineskip}

\mattins
\psalmAntiphon{His face did shine {\dag} as the sun, and his raiment was white as the light, alleluia.}

\begin{rubric}
	Invitatory Hymn of I Evensong.
\end{rubric}

\mattinsOfficeHymn
\hymnLett{O}{love} of Jesus, sweet and dear,

\hymnSecondLine{When to the heart thou dost appear,}

\begin{hangparas}{1.25em}{1}
Away its clouds and darkness roll,

And sweetness overflows the soul.\\

2. How happy he who feels thy light,

Thou Sharer of the Father's might,

True Radiance of our native land,

Surpassing all we understand.\\

3. Thou Brightness of the Father’s throne,

Goodness that never can be known,

The fulness of thy love impart

By thy true Presence in the heart.\\

4. All glory, Lord, to thee we pay,

Transfigur'd on the mount to-day;

All glory, as is ever meet,

To Father and to Paraclete. Amen.\\
\end{hangparas}

℣. A crown of gold is upon his head.

℟. A visible sign of holiness, glory, and honour.

\mattinsAntiphon{And behold a voice {\dag} out of the cloud, which said, This is my beloved Son, in whom I am well pleased; hear ye him, alleluia.}

\fussy	
\end{paracol}

\secondEvensong
\psalmAntiphon{And behold, {\dag} there appeared unto them Moses and Elias, talking with Jesus.}

\begin{rubric}
	Office Hymn \& Versicle of I Evensong.
\end{rubric}

\secondEvensongAntiphon{And when the disciples heard it, {\dag} they fell on their face, and were sore afraid: and Jesus came and touched them, and said, Arise, and be not afraid, alleluia.}


\collect
\lett{O}{God,} who on this day didst reveal from heaven thine only-begotten Son, transfigured in a wonderful manner, to the fathers of both testaments, grant unto us, we beseech thee, that by actions acceptable unto thee we may attain unto the perpetual contemplation of his glory in whom thou hast testified that thou, the Father, wast well pleased. Through the same.

\begin{rubric}
	 Commemoration \blackrubric{of Sts. Sixtus II, Felicissimus, \& Agapitus,} from the Common of Many Martyrs (p. \pageref{CommonMartyrsNotBishopsCollect}), the second Collect.
\end{rubric}


%MANUAL ADJUSTMENT:
\clearpage
\subbymemorial{6 August}{Sts. Sixtus II, Felicissimus, and Agapitus}{M.}{MartyrRed}
\feastday{{Sts. Sixtus II \&c.}}
\feastdate{6}{August}

%%CHECK:
\begin{rubric}
	Common of Many Martyrs (p. \pageref{CommonMartyrsNotBishopsCollect}), the second Collect.
\end{rubric}


\festivalLL{Second Class Double}{7 August}{Most Holy Name of Jesus Christ}
\feastday{Holy Name}
\feastdate{7}{August}

\begin{rubric}
    The propers of the Christmastide Feast Day of the same name (p. \pageref{MostHolyName}).
\end{rubric}


\subbymemorial{7 August}{St. Donatus}{B.M.}{MartyrRed}
\feastday{{St. Donatus}}
\feastdate{7}{August}

\memorialCollect
\lett{O}{God,} the glory of thy priests: grant, we beseech thee; that we may perceive the succour of thy holy Martyr and Bishop Donatus, whose festival we celebrate. Through.


\subbymemorial{8 August}{Sts. Cyriacus, Largus, and Smaragdus}{M.}{MartyrRed}
\feastday{{Sts. Cyriacus \&c.}}
\feastdate{8}{August}

\memorialCollect
\lett{O}{God,} who makest us glad with the yearly solemnity of thy holy Martyrs Cyriacus, Largus, and Smaragdus: mercifully grant; that as we now celebrate their birthday, so we may imitate their constancy in suffering. Through.

\begin{rubric}
	If today be Saturday, Commemoration \blackrubric{of the Vigil of St. Lawrence} (p. \pageref{LawrenceVigil}).
\end{rubric}


\festival{Vigil}{9 August}{Vigil of St. Lawrence}{M.}{MartyrRed}\label{LawrenceVigil}
\feastday{{St. Lawrence Vigil}}
\feastdate{9}{August}

%\collect
\lett{A}{ssist} us, O Lord, in these our supplications: and at the intercession of thy blessed Martyr Lawrence, whose festival we prevent, graciously bestow upon us thy perpetual mercy. Through.

\begin{rubric}
	Commemoration \blackrubric{of St. Romanus} (p. \pageref{RomanusCollect}).
\end{rubric}


\subbymemorial{9 August}{St. Romanus}{M.}{MartyrRed}
\feastday{{St. Romanus}}
\feastdate{9}{August}

\memorialCollect\label{RomanusCollect}
\lett{G}{rant,} we beseech thee, Almighty God: that, at the intercession of blessed Romanus, thy Martyr, we may both be delivered from all adversities which may happen to the body, and from all evil thoughts which may assault and hurt the soul. Through.


%MANUAL ADJUSTMENT:
\clearpage
\festival{Second Class Double, Simple Octave}{10 August}{St. Lawrence}{M.}{MartyrRed}
\feastday{St. Lawrence}
\feastdate{10}{August}

\begin{rubric}
	Common of a Martyr (p. \pageref{CommonMartyr}).
\end{rubric}

\firstEvensong
\psalmAntiphon{Lawrence {\dag} hath wrought a pious work; who by the sign of the Cross enlightened the blind.}\\

℣. Thou hast crowned him with glory and honour, O Lord.

℟. And madest him to have dominion of the works of thy hands.

\firstEvensongAntiphon{Lawrence the Deacon {\dag} hath wrought a pious work; who by the sign of the Cross enlightened the blind: and distributed to the poor the Church's treasures.}

\mattins
\psalmAntiphon{My soul {\dag} cleaveth fast unto thee, forasmuch as my flesh hath been burned for thee, O my God.}\\

℣. He hath dispersed abroad and given to the poor.
℟. And his righteousness remaineth for ever.

\mattinsAntiphon{On the iron grate, {\dag} O God, I denied thee not; and when fire was kindled beneath me, O Christ, I confessed thee: thou hast proved my heart and hast visited me in the night season; thou hast tried me with fire, and hast found no wickedness in me.}

\secondEvensong
\psalmAntiphon{Blessed Lawrence {\dag} prayed, saying: I render thanks unto thee, O Lord, who hast counted me worthy to enter within thy gates.}\\

℣. Lawrence the Deacon hath wrought a pious work.

℟. Who by the sign of the Cross enlightened the blind.

\secondEvensongAntiphon{Blessed Lawrence, {\dag} when laid and burning on the iron grating, spake to the impious tyrant, saying: The feast is ready, turn and eat; but the Church's treasures, which thou claimest, have been garnered up in heaven, by the hands of the poor and needy.}

\collect
\lett{G}{rant} to us, we beseech thee, Almighty God: that we may quench the flames of our sins; as thou didst enable blessed Lawrence to overcome the fires of his torments. Through.


\subbymemorial{11 August}{Sts. Tiburtius and Susanna}{M.}{MartyrRed}
\feastday{{\small Sts. Tiburtius \& Susanna}}
\feastdate{11}{August}

\memorialCollect
\lett{O}{Lord,} let the protection of thy holy Martyrs Tiburtius and Susanna continually defend us: forasmuch as thou failest not to look with mercy on those to whom thou dost grant the succour of their assistance. Through.


\festival{Double}{13 August}{St. Maximus of Constantinople}{C.}{FestiveWhite}
\feastday{{St. Maximus Constantinople}}
\feastdate{13}{August}

%%CHECK:
\begin{rubric}
	Common of a Confessor not a Bishop (p. \pageref{CommonConfessorNotBishop}), the second Collect. Commemoration \blackrubric{of Sts. Hippolytus \& Cassian} (p. \pageref{HippolytusCassianCollect}).
\end{rubric}


\subbymemorial{13 August}{Sts. Hippolytus and Cassian}{M.}{MartyrRed}
\feastday{{\small Sts. Hippolytus \& Cassian}}
\feastdate{13}{August}

\memorialCollect\label{HippolytusCassianCollect}
\lett{G}{rant,} we beseech thee, Almighty God: that the venerable solemnity of thy blessed Martyrs, Hippolytus and Cassian, may increase our devotion and set forward our salvation. Through.

\begin{rubric}
	If today be Saturday, the anticipated Vigil of the Assumption of the Blessed Virgin Mary is kept, as is noted on the following day, but with Commemoration \blackrubric{of Sts. Hippolytus and Cassian,} instead of St. Eusebius.
\end{rubric}


\festivalLL{Vigil}{14 August}{Vigil of the Assumption\\of the Blessed Virgin Mary}
\feastday{{\small Assumption B.V.M. Vigil}}
\feastdate{14}{August}


%\collect
\lett{O}{God,} who didst vouchsafe to choose the virgin womb of blessed Mary wherein to make thy dwelling: grant, we beseech thee; that, being defended by her protection, we may by thee be enabled to attain with gladness to her festival. Who livest.

\begin{rubric}
	Commemoration \blackrubric{of St. Eusebius} (p. \pageref{EusebiusCollect}).
\end{rubric}


%MANUAL ADJUSTMENT:
\clearpage
\subbymemorial{14 August}{St. Eusebius}{C.}{FestiveWhite}
\feastday{{St. Eusebius}}
\feastdate{14}{August}

\memorialCollect\label{EusebiusCollect}
\lett{O}{God,} who makest us glad with the yearly solemnity of blessed Eusebius thy Confessor: mercifully grant; that we, who celebrate his birthday, may by his example, be drawn nearer unto thee. Through.


\festivalLL{First Class Double, Common Octave}{15 August}{Assumption of the Blessed Virgin Mary}
\feastday{Assumption B.V.M.}
\feastdate{15}{August}

\begin{rubric}
	Common of the Blessed Virgin Mary (p. \pageref{CommonBVM}).
\end{rubric}

\openingSentence{The holy Mother of God is exalted. Above the choirs of Angels to the heavenly kingdom.}

\firstEvensong
\psalmAntiphon{Mary is taken up {\dag} into heaven: the Angels triumph, exalting and blessing the Lord.}

\begin{multicols}{2}
\sloppy

\hymnLett{O}{with} what glorious lustre thou shinest.

\hymnSecondLine{Daughter of royalty, David's descendant!}

\begin{hangparas}{1.25em}{1}
Throned in majesty, Mary the Virgin,

Thou `mid the blessed ones sittest exalted.\\

2. Keeping thy virginal honour unspotted

E'en in thy motherhood, chastely thou gavest

Shrine for the Holy One, Lord of the Angels;

Thus in humanity God was incarnate;\\

3. Whom the whole universe lowly adoreth.

Duly on bended knee tendering homage:

We on thy festival pray him to grant us

Light and felicity, darkness dispelling.\\

4. This, of thy clemency, Father of glory,

Grant through thine only Son, who, with the Spirit,

Evermore one with thee liveth and reigneth

In the bright firmament, ordering all things. Amen.\\
\end{hangparas}

℣. The holy Mother of God is exalted.

℟. Above choirs of Angels to the heavenly kingdom.

\firstEvensongAntiphon{O most prudent Virgin, {\dag} whither goest thou, shining resplendent like the glowing dawn? Daughter of Sion, thou art all comely and beautiful, fair as the moon, clear as the sun.}

\fussy	
\end{multicols}


\mattins
\psalmAntiphon{The Virgin Mary {\dag} is taken up into heavenly chambers, wherein the King of kings sitteth upon a starry throne.}\\

℣. The holy Mother of God is exalted.

℟. Above choirs of Angels to the heavenly kingdom.

\mattinsAntiphon{Who is she {\dag} that riseth up as the morning, fair as the moon, clear as the sun, and terrible as an army with banners?}

\secondEvensong
\psalmAntiphon{Thou art fair {\dag} and comely, O daughter of Jerusalem: thou art terrible as an army with banners.}\\

\secondEvensongAntiphon{On this day {\dag} the Virgin Mary went up into heaven: rejoice, because she reigneth with Christ for ever and ever.}


\collect
\lett{W}{e} beseech thee, O Lord, let us be continually aided by the sacred feast of this day, whereon the holy Mother of God underwent death in this world, and yet could not be holden by the chains of death: who did bring forth thy Son our Lord. Who liveth.


\festivalAO{Second Class Double}{16 August}{St. Joachim, Father of the Blessed Virgin Mary}
\feastday{St. Joachim}
\feastdate{16}{August}

\begin{rubric}
	Common of a Confessor not a Bishop (p. \pageref{CommonConfessorNotBishop}).
\end{rubric}

\begin{rubric}
	\textsc{Note,} The following Versicle \& Antiphon is said in I Evensong \& Mattins.
\end{rubric}

℣. His seed shall be mighty upon earth.

℟. The generation of the faithful shall be blessed.

\antiphon{Mag. \& Ben.}{Let us praise a man famous {\dag} in his generation, with whom the Lord established the blessing of all nations, and the covenant, and made it rest upon his head.}

%MANUAL ADJUSTMENT:
\vspace{-1\baselineskip}

\collect
\lett{O}{God,} who from amongst all thy Saints didst choose blessed Joachim to be the father of the Mother of thy Son: grant, we beseech thee; that we who venerate his festival may also continually perceive his advocacy. Through the same.


%MANUAL ADJUSTMENT:
\clearpage
\festival{Simple}{17 August}{Octave Day of St. Lawrence}{M.}{MartyrRed}
\feastday{{St. Lawrence Octave Day}}
\feastdate{17}{August}

%\collect
\lett{S}{tir} up, O Lord, in thy Church that Spirit whom the blessed Levite Lawrence served: that we, being filled with the same, may study to love that which he loved, and to perform in deed that which he taught. Through . . . in the unity of the same Holy Spirit.


\subbymemorial{18 August}{Empress St. Helen}{Ma.}{FestiveWhite}
\feastday{{St. Helen}}
\feastdate{18}{August}

\memorialCollect
\lett{O}{Lord} Jesu Christ, who didst reveal unto blessed Helen the place where thy Cross lay hid, that through her thou mightest enrich thy Church with this precious treasure: grant unto us at her intercession; that by the ransom of the life-giving tree we may attain unto the rewards of everlasting life. Who livest.


\subbymemorial{18 August}{St. Agapitus}{M.}{MartyrRed}
\feastday{{St. Agapitus}}
\feastdate{18}{August}

\memorialCollect
\lett{O}{Lord,} let thy Church trust with gladness in the advocacy of thy blessed Martyr Agapitus: that by his glorious prayers she may continue in devotion, and abide in safety. Through.


\festivalLL{Greater Double}{22 August}{Octave Day of the Assumption of the Blessed Virgin Mary}
\feastday{{Assumption Octave Day}}
\feastdate{22}{August}

\begin{rubric}
	As on the Feast, Commemoration \blackrubric{of Sts. Timothy, Hippolytus, \& Symphorian} (p. \pageref{TimothyHippolytusSymphorianCollect}).
\end{rubric}
\begin{rubric}
	If today be Saturday, \blackrubric{the Vigil of St. Bartholomew} is kept, as is noted on the following day.
\end{rubric}


\subbymemorial{22 August}{Sts. Bishop Timothy, Hippolytus, and Symphorian}{M.}{MartyrRed}
\feastday{{Sts. Timothy \&c.}}
\feastdate{22}{August}

\memorialCollect\label{TimothyHippolytusSymphorianCollect}
\lett{W}{e} beseech thee, O Lord, graciously to impart unto us thy help: and, at the intercession of thy blessed Martyrs Timothy, Hippolytus, and Symphorian, stretch forth upon us the right hand of thy mercy. Through.


%MANUAL ADJUSTMENT:
\clearpage
\festival{Vigil}{23 August}{Vigil of St. Bartholomew}{A.}{LentenViolet}
\feastday{{St. Bartholomew Vigil}}
\feastdate{23}{August}

\begin{rubric}
	Common of the Vigil of Apostles (p. \pageref{CommonVigilApostles}).
\end{rubric}


\festival{Second Class Double}{24 August}{St. Bartholomew}{A.}{FestiveWhite}
\feastday{St. Bartholomew}
\feastdate{24}{August}

\begin{rubric}
	Common of Apostles (p. \pageref{CommonApostles}).
\end{rubric}

\firstEvensongAntiphon{They will deliver you up {\dag} to the councils, and they will scourge you in their synagogues; and ye shall be brought before governors and kings for my sake, for a testimony against them and the Gentiles.}

\mattinsAntiphon{Ye which have forsaken all, {\dag} and followed me, shall receive an hundredfold, and shall inherit everlasting life.}

\secondEvensongAntiphon{Be ye valiant in warfare, {\dag} and contend ye with the old serpent: and ye shall receive an eternal kingdom, alleluia.}


\collect
\lett{O}{Almighty} and Everlasting God, who didst give to thine Apostle Bartholomew grace truly to believe and to preach thy Word; Grant, we beseech thee, unto thy Church to love that Word which he believed, and both to preach and receive the same. Through.


\subbymemorial{26 August}{Pope St. Zephyrinus of Rome}{B.M.}{MartyrRed}
\feastday{{St. Zephyrinus}}
\feastdate{26}{August}

\memorialCollect
\lett{G}{rant,} we beseech thee, Almighty God: that we who rejoice in the merits of blessed Zephyrinus, thy Martyr and Bishop, may be instructed by his example. Through.


%MANUAL ADJUSTMENT:
\clearpage
\festival{Greater Double}{28 August}{St. Augustine of Hippo}{B.C.D.}{FestiveWhite}
\feastday{{St. Augustine Hippo}}
\feastdate{28}{August}

\begin{rubric}
	Common of a Confessor Bishop (p. \pageref{CommonConfessorBishop}).
\end{rubric}

%\collect
\lett{A}{ssist} us, Almighty God, in these our supplications: that as thou dost suffer us to put our trust and confidence in thy mercy, so, at the intercession of blessed Augustine thy Confessor and Bishop, thou wouldest graciously vouchsafe unto us the wonted effects of thy compassion. Through.

\begin{rubric}
    Commemoration \blackrubric{of St. Hermes} (p. \pageref{HermesCollect}).
\end{rubric}


\subbymemorial{28 August}{St. Hermes}{M.}{MartyrRed}
\feastday{{St. Hermes}}
\feastdate{28}{August}

\memorialCollect\label{HermesCollect}
\lett{O}{God,} who didst strengthen blessed Hermes thy Martyr with the virtue of constancy in his passion: Grant unto us by his example; to despise for love of thee the prosperity of the world, and to fear none of its adversities. Through.


\festivalAO{Greater Double}{29 August}{Decollation of St. John Baptist}
\feastday{Decollation}
\feastdate{29}{August}

\begin{rubric}
	Hymns \& Versicles of the Common of Apostles (p. \pageref{CommonApostles}).
\end{rubric}

\firstEvensong
\psalmAntiphon{Herod himself had sent forth {\dag} and laid hold upon John, and bound him in prison, for Herodias' sake.}

\firstEvensongAntiphon{Herod sent {\dag} an executioner, and commanded him to behead John in the prison: and when his disciples heard of it, they came, and took up his body, and laid it in a tomb.}

\mattins
\psalmAntiphon{John reproved {\dag} Herod for the sake of Herodias, whom he took for wife from his brother Philip.}

\mattinsAntiphon{Herod sent {\dag} an executioner, and commanded him to behead John in the prison: and when his disciples heard of it, they came, and took up his body, and laid it in a tomb.}

\secondEvensong
\psalmAntiphon{Give me in a charger {\dag} the head of John the Baptist: and the king was sorry for his oath’s sake.}

\secondEvensongAntiphon{The king, the unbeliever, {\dag} sent his detestable servants, and commanded them to behead John the Baptist.}

\collect
\lett{W}{e} beseech thee, O Lord: that the venerable festival of thy Forerunner and Martyr, Saint John Baptist, may effectually bestow upon us thy succour unto our salvation. Who livest.
\begin{rubric}
	 Commemoration \blackrubric{of St. Sabina,} from the Common of Holy Matrons (p. \pageref{CommonHolyMatronsCollect}).
\end{rubric}


\subbymemorial{29 August}{St. Sabina}{M.}{MartyrRed}
\feastday{{St. Sabina}}
\feastdate{29}{August}

\begin{rubric}
	Common of Holy Matrons (p. \pageref{CommonHolyMatronsCollect}).
\end{rubric}


\subbymemorial{30 August}{Sts. Felix and Adauctus}{M.}{MartyrRed}
\feastday{{Sts. Felix \& Adauctus}}
\feastdate{30}{August}

\memorialCollect
\lett{O}{Lord,} we humbly entreat thy Majesty: that, like as thou dost continually gladden us with the commemoration of thy Saints; so thou wouldest evermore defend us with their supplication. Through.


\subbymemorial{31 August}{St. Aidan of Lindisfarne}{B.C.}{FestiveWhite}
\feastday{{St. Aidan Lindisfarne}}
\feastdate{31}{August}

\begin{rubric}
	Common of a Confessor Bishop (p. \pageref{ConfessorBishopCollect}).
\end{rubric}


\subbymemorial{1 September}{St. Giles}{Ab.}{FestiveWhite}
\feastday{{St. Giles}}
\feastdate{1}{September}

\begin{rubric}
	Common of a Confessor not a Bishop (p. \pageref{AbbotCollect}).
\end{rubric}


\subbymemorial{1 September}{Twelve Holy Brothers}{M.}{MartyrRed}
\feastday{{12 Holy Brothers}}
\feastdate{1}{September}

\memorialCollect
\lett{O}{Lord,} let the crown of the brethren thy Martyrs, cause us to rejoice: that we may thereby be strengthened and increased in our faith; and comforted by their manifold intercession. Through.


%MANUAL ADJUSTMENT:
\clearpage
\subbymemorial{2 September}{King St. Stephen of Hungary (m.t.v.)}{C.}{FestiveWhite}
\feastday{{St. Stephen Hungary}}
\feastdate{2}{September}

\memorialCollect
\lett{G}{rant,} we beseech thee, Almighty God unto thy Church: that as thy blessed Confessor Stephen, while he reigned on earth, did spread abroad her faith, so she may be found worthy to have him for her glorious defender in the heavens. Through.


\festival{Double}{4 September}{St. Gorazde of Prague}{B.M.}{MartyrRed}
\feastday{{St. Gorazde Prague}}
\feastdate{4}{September}

%%CHECK:
\begin{rubric}
	Common of a Martyr (p. \pageref{CommonMartyr}), the second Collect.
\end{rubric}


\festivalLL{Second Class Double, Simple Octave}{8 September}{Nativity of the Blessed Virgin Mary}
\feastday{Nativity B.V.M.}
\feastdate{8}{September}

\begin{rubric}
	Common of the Blessed Virgin Mary (p. \pageref{CommonBVM}).
\end{rubric}

\firstEvensong
\psalmAntiphon{The Nativity {\dag} of the glorious Virgin Mary, of the seed of Abraham, sprung from the tribe of Juda, of the noble stem of David.}\\

℣. To-day is the Nativity of the holy Virgin Mary.

℟. Whose glorious life illumineth all the churches.

\firstEvensongAntiphon{Let us celebrate {\dag} the worshipful Nativity of the blessed and glorious Virgin Mary: for she hath obtained the dignity of Motherhood, and yet lost not her maiden purity.}

\mattins
\psalmAntiphon{Sprung from a royal race, {\dag} lo, Mary shineth resplendent: by whose prayers we ask devoutly that we may be aided both in mind and spirit.}\\

℣. To-day is the Nativity of the holy Virgin Mary.

℟. Whose glorious life illumineth all the churches.

\mattinsAntiphon{To-day let us celebrate {\dag} with due solemnity the Nativity of God’s Mother, the ever Virgin Mary: from whom the Son of the Highest proceeded, alleluia.}

\secondEvensong
\psalmAntiphon{With all our heart and soul, {\dag} let us sing glory unto Christ on this holy solemnity of Mary, the exalted Mother of God.}\\

℣. To-day is the Nativity of the holy Virgin Mary.

℟. Whose glorious life illumineth all the churches.

\secondEvensongAntiphon{Thy Nativity, {\dag} O Virgin Mother of God, hath proclaimed joyful tidings unto all the world: for out of thee hath arisen the Sun of righteousness, even Christ our God: who, taking away the curse, hath bestowed a blessing; and confounding death, hath given unto us life everlasting.}

\collect
\lett{W}{e} beseech thee, O Lord, to grant unto us thy servants the gift of heavenly grace: that as the childbearing of the blessed Virgin was unto us the beginning of salvation; so the devout observance of her Nativity may bestow an increase of peace. Through.


\subbymemorial{8 September}{St. Hadrian}{M.}{MartyrRed}
\feastday{{St. Hadrian}}
\feastdate{8}{September}

\begin{rubric}
	Common of a Martyr (p. \pageref{CommonMartyrNotBishopCollect}).
\end{rubric}


\subbymemorial{9 September}{St. Gorgonius}{M.}{MartyrRed}
\feastday{{St. Gorgonius}}
\feastdate{9}{September}

\memorialCollect
\lett{L}{et} thy Saint Gorgonius, O Lord, gladden us by his intercession; and make us to rejoice in this holy solemnity. Through.

\lettspace


\subbymemorial{11 September}{Sts. Protus and Hyacinth}{M.}{MartyrRed}
\feastday{{Sts. Protus \& Hyacinth}}
\feastdate{11}{September}

\memorialCollect
\lett{O}{Lord,} let the meritorious confession of thy blessed Martyrs, Protus and Hyacinth, comfort us; and may their loving intercession ever defend us. Through.


%MANUAL ADJUSTMENT:
\clearpage
\festivalLL{Greater Double}{12 September}{Most Holy Name of Mary}
\feastday{{Holy Name of Mary}}
\feastdate{12}{September}

\begin{rubric}
	Common of the Blessed Virgin Mary (p. \pageref{CommonBVM}).
\end{rubric}

\begin{rubric}
	\textsc{Note,} The following Antiphon is said during I Evensong.
\end{rubric}

\firstEvensongAntiphon{O holy Mary, {\dag} help thou the suffering, strengthen the faint-hearted, comfort the sorrowful; pray for the people, entreat for the clergy, intercede for all womankind vowed unto God: may all acknowledge the help of thy prayer, who celebrate the commemoration of thine holy Name.}

\collect
\lett{G}{rant,} we beseech thee, Almighty God: that thy faithful people, who rejoice in the Name and protection of the most Holy Virgin Mary; may by her loving intercession be delivered from all evils upon earth, and be found worthy to attain unto everlasting joys in heaven. Through.


\festivalLL{Greater Double}{14 September}{Exaltation of the Holy Cross}
\feastday{Roodmas}
\feastdate{14}{September}

\begin{rubric}
	Hymns, Versicles, \& Antiphons \blackrubric{of the Feast of the Invention of the Holy Cross} (p. \pageref{InventionCross}), omitting the \blackrubric{Alleluia.}
\end{rubric}

\firstEvensong
\psalmAntiphon{O mighty {\dag} work of mercy! death then died, when life died upon the Tree.}

\mattins
\psalmAntiphon{Behold the Cross of the Lord; {\dag} flee away, ye adversaries: the Lion of Judah, the Root of David, hath prevailed, alleluia.}

\secondEvensong
\psalmAntiphon{But as for us, {\dag} it behoveth us to glory in the Cross of our Lord Jesus Christ.}

\secondEvensongAntiphon{O Cross {\dag} exceeding blessed, which alone wast counted worthy to bear the Lord, the King of heaven, alleluia.}

\collect
\lett{O}{God,} who on this day dost gladden us with the yearly solemnity of the exaltation of the holy Cross: grant, we beseech thee; that as we have known the mystery of thy Son on earth, so we may attain unto the rewards of his redemption in heaven. Through the same


\festivalLL{Second Class Double}{15 September}{Seven Sorrows of the Blessed Virgin Mary}
\feastday{Seven Sorrows B.V.M.}
\feastdate{15}{September}

\begin{paracol}{2}
\sloppy

\firstEvensong
\psalmAntiphon{Whither {\dag} is thy beloved gone, O thou fairest among women? whither is thy beloved turned aside? that we may seek him with thee.}

\firstEvensongHymn
\hymnLett{N}{ow} let the darkling eve

\hymnSecondLine{Mount suddenly on high,}

\hymnThirdLine{The sun affrighted reave}

\begin{hangparas}{1.25em}{1}
His splendours from the sky,

While I in silence grieve

O'er the mocked agony

And the divine catastrophe.\\

2. Grief-drench'd, thou dost appear

With heart of adamant,

O Mother; and dost hear

The Great Hierophant,

Upon his wooden bier

Locked in the arms of Death,

Utter in groans his parting breath.\\

3. What lookest thou upon,

Mangled and bruised and torn?

Ah, `tis the very Son

Thy yearning breast hath borne!

Surely, each breaking moan

And each deep-mouthed wound

Its fellow in thy heart hath found!\\

4. Surely, the taunts and woes,

The scourge, the dripping thorn,

The spitting and the blows,

The gall, the lance, the scorn;

Surely, each torment throws

A poison-dart at thee,

Crushed by their varied tyranny.\\

5. Yet thou with patient mien

Beneath his Cross dost stand,

Nobler in this, I ween,

Than all the martyr-band:

A thousand deaths, O Queen,

Upon thy spirit lie,

Yet thou, O marvel! dost not die.\\

6. O Holy Trinity,

Let earth and heaven raise

Their song of laud to thee

The while my spirit prays:

When evil comes to me,

The strength do thou impart

That erst upheld the Virgin’s heart! Amen.\\
\end{hangparas}

℣. Pray for us, O Queen of Martyrs.

℟. Who didst stand by the Cross of Jesus.

\firstEvensongAntiphon{Look not {\dag} upon me, because I am black, because the sun hath looked upon me: my mother's children were angry with me.}


\switchcolumn

\mattins
\psalmAntiphon{Look away from me; {\dag} I will weep bitterly, labour not to comfort me.}

\mattinsInvitatoryHymn
\hymnLett{W}{hat} a flood of tears and sorrow,

\hymnSecondLine{What an agony of pain,}

\begin{hangparas}{1.25em}{1}
Overwhelm'd the mourning Mother,

When her Son, so lately slain,

From the blood-stain'd Cross was lower'd

To her loving arms again.\\

2. With her tears she bath'd his body,

Bruis'd and stricken yet most sweet;

Underneath his heart the spear-wound;

Pierced hands and crimson'd feet:

Then, in deep maternal anguish,

Kiss'd each wound-print, as was meet.\\

3. By these tears of thine, O Mother,

By the death of thy dear Son,

By his cruel wounds empurpled,

By the pain surpass'd by none,

Move our spirits to compassion

Till our grief with thine is one.\\

4. To the Father, Son, and Spirit,

To the Trinity most high,

Co-eternal and co-equal,

Everlasting glory be;

Praise unending, honour, blessing,

Now and through eternity. Amen.
\end{hangparas}

\mattinsOfficeHymn
\hymnLett{O}{God,} in whom all grace doth dwell,

\hymnSecondLine{Grant us the grace to ponder well}

\begin{hangparas}{1.25em}{1}
The Virgin Mother’s sorrows sev'n,

The cruel wounds to Jesus giv'n.\\

2. O may the tears which Mary pour'd

Gain for us pardon of the Lord;

The holy tears which in their worth

Excel all penances of earth.\\

3. And may the contemplation sore

Of those five wounds which Jesus bore

With all the Virgin’s sorrows be

Our joy throughout eternity.\\

4. Glory to thee, O Lord, we give,

Who died to make thy servants live;

Whom with the Father we adore

And Holy Spirit evermore. Amen.\\
\end{hangparas}

℣. O Virgin Mary, by thy sorrow’s might.

℟. Make us rejoice in heaven’s kingdom bright.

\mattinsAntiphon{Come ye, {\dag} and let us go up to the mountain of the Lord, and see if there be any sorrow like unto my sorrow.}

\fussy	
\end{paracol}

\secondEvensong
\psalmAntiphon{He hath no form {\dag} nor comeliness; and when we shall see him, there is no beauty that we should desire him.}

\begin{rubric}
	Office Hymn \& Versicle of I Evensong.
\end{rubric}

\secondEvensongAntiphon{Sorrow oppresseth me, {\dag} my face is swollen with weeping, and on my eyelids is the shadow of death.}

\collect
\lett{O}{God,} in whose passion according to the prophecy of Simeon, a sword of sorrow did pierce the most sweet soul of the glorious Virgin Mother Mary: mercifully grant; that we, who devoutly call to mind her sorrows, may obtain the blessed effects of thy passion. Who livest.
\begin{rubric}
	 Commemoration \blackrubric{of St. Nicomedes} (p. \pageref{NicomedesCollect}).
\end{rubric}


\subbymemorial{15 September}{St. Nicomedes}{M.}{MartyrRed}
\feastday{{St. Nicomedes}}
\feastdate{15}{September}

\memorialCollect\label{NicomedesCollect}
\lett{A}{ssist,} O Lord, thy people: that as they do profit by the glorious merits of blessed Nicomedes thy Martyr, so his advocacy may at all times succour them to the obtaining of thy mercy. Through.


\subsection{Ember Wednesday in Autumn}
\feastday{Autumnal Emberday}
\begin{inhead}
    {Second Class Feria\\
Wednesday after the Exaltation of the Holy Cross (14 September)\\
On this day and the following Friday and Saturday, no Feast is kept unless it be a First or Second Class Double.}
\end{inhead}
\fancyhead[RE,LO]{Ember Wednesday}

\mattinsAntiphon{This kind {\dag} of evil spirit can come forth by nothing but by prayer and fasting.}

%\collect
\lett{W}{e} beseech thee, O Lord, that our frailty may be upheld by the healing of thy loving kindness: that what by its own nature is ready to decay, may by thy mercy be renewed. Through.


\subsection{Ember Friday in Autumn}
\feastday{Autumnal Emberday}
\begin{inhead}
    {Second Class Feria\\
Friday after Ember Wednesday}
\end{inhead}
\fancyhead[RE,LO]{Ember Friday}

\mattinsAntiphon{A woman {\dag} in the city, which was a sinner, stood at the Lord’s feet behind him, and began to wash his feet with tears, and did wipe them with the hairs of her head, and kissed his feet, and anointed them with the ointment.}

%\collect
\lett{G}{rant,} we beseech thee, Almighty God: that we, who year by year devoutly keep this holy observance, may be acceptable unto thee both in body and in soul. Through.


\subsection{Ember Saturday in Autumn}
\feastday{Autumnal Emberday}
\begin{inhead}
    {Second Class Feria\\
Saturday after Ember Friday}
\end{inhead}
\fancyhead[RE,LO]{Ember Saturday}

\mattinsAntiphon{Give light, O Lord, {\dag} to them that sit in darkness: and guide our feet into the way of peace, thou God of Israel.}

%\collect
\lett{A}{lmighty} and everlasting God, who by salutary continence bestowest healing in body and soul: we humbly entreat thy Majesty; that, thou wouldest mercifully look upon the devout prayers and fasting of thy people, and grant us help both in this life and that which is to come. Through.


\festival{Double}{16 September}{Sts. Pope Cornelius and Cyprian}{B.M.}{MartyrRed}
\feastday{{Sts. Cornelius \& Cyprian}}
\feastdate{16}{September}

\begin{rubric}
	Common of Many Martyrs (p. \pageref{CommonMartyrs}).
\end{rubric}

%\collect
\lett{P}{rotect} us, O Lord, we beseech thee, who observe the feast of thy blessed Martyrs and Bishops Cornelius and Cyprian: and grant that by their meritorious supplication we may find favour in thy sight. Through.

\begin{rubric}
    Commemoration \blackrubric{of St. Euphemia, Lucy, and Geminian} (p. \pageref{EuphemiaCollect}).
\end{rubric}


\subbymemorial{16 September}{Sts. Euphemia, Lucy, and Geminian}{M.}{MartyrRed}
\feastday{{Sts. Euphemia \&c.}}
\feastdate{16}{September}

\memorialCollect\label{EuphemiaCollect}
\lett{G}{rant,} O Lord, that our prayers in this time of our rejoicing may be brought to good effect: that as with yearly service we recall the day of the passion of thy holy Martyrs, Euphemia, Lucy, and Geminian, so we may imitate the steadfastness of their faith. Through.


\festival{Double}{19 September}{Sts. Januarius and Companions}{M.}{MartyrRed}
\feastday{{Sts. Januarius \&c.}}
\feastdate{19}{September}

\begin{rubric}
	Common of Many Martyrs (p. \pageref{CommonMartyrs}). Commemoration \blackrubric{of St. Theodore of Canterbury,} from the Common of a Confessor Bishop (p. \pageref{ConfessorBishopCollect}).
\end{rubric}


%MANUAL ADJUSTMENT:
\clearpage
\subbymemorial{19 September}{St. Theodore of Canterbury}{B.C.}{FestiveWhite}
\feastday{{St. Theodore Canterbury}}
\feastdate{19}{September}

\begin{rubric}
	Common of a Confessor Bishop (p. \pageref{ConfessorBishopCollect}).
\end{rubric}

\begin{rubric}
	If today be Saturday, the anticipated Vigil of St. Matthew is kept.
\end{rubric}


\festival{Double}{20 September}{Sts. Eustace and Companions}{M.}{MartyrRed}
\feastday{{Sts. Eustace \&c.}}
\feastdate{20}{September}

%%CHECK:
\begin{rubric}
	Common of Many Martyrs (p. \pageref{CommonMartyrsNotBishopsCollect}), the second Collect.
\end{rubric}


\festival{Vigil}{20 September}{Vigil of St. Matthew}{A.}{LentenViolet}
\feastday{{St. Matthew Vigil}}
\feastdate{20}{September}

\begin{rubric}
	Common of Vigils of the Apostles (p. \pageref{CommonVigilApostles}).
\end{rubric}


\festival{Second Class Double}{21 September}{St. Matthew}{A.}{FestiveWhite}
\feastday{St. Matthew}
\feastdate{21}{September}

\begin{rubric}
	Common of Apostles (p. \pageref{CommonApostles}).
\end{rubric}

%\collect
\lett{M}{ay} we be assisted, O Lord, by the prayers of blessed Matthew, the Apostle and Evangelist: that those things, which of ourselves we cannot obtain, may be vouchsafed unto us by his intercession. Through.


\subbymemorial{22 September}{Sts. Maurice and Companions}{M.}{MartyrRed}
\feastday{{Sts. Maurice \&c.}}
\feastdate{22}{September}

\memorialCollect
\lett{G}{rant,} we beseech thee, Almighty God: that the solemn festival of thy holy Martyrs, Maurice and his Companions, may in such wise gladden us: that, as we do lean upon their advocacy, so we may glory in their heavenly birth. Through.


\subbymemorial{23 September}{Pope St. Linus of Rome}{B.M.}{MartyrRed}
\feastday{{Pope St. Linus}}
\feastdate{23}{September}

\memorialCollect
\lett{O}{God,} who makest us glad with the yearly solemnity of blessed Linus thy Martyr and Bishop: mercifully grant; that, as we now celebrate his birthday, so we may likewise rejoice in his protection. Through.

\begin{rubric}
    Commemoration \blackrubric{of St. Thecla} (p. \pageref{TheclaCollect}).
\end{rubric}


%MANUAL ADJUSTMENT:
\clearpage
\subbymemorial{23 September}{St. Thecla}{V.M.}{MartyrRed}
\feastday{{St. Thecla}}
\feastdate{23}{September}

\memorialCollect\label{TheclaCollect}
\lett{G}{rant,} we beseech thee, Almighty God: that we, who celebrate the birthday of blessed Thecla, thy Virgin and Martyress; may both rejoice in her yearly solemnity, and likewise profit by the example of her faith. Through.


\subbymemorial{26 September}{Sts. Cyprian and Justina}{M.}{MartyrRed}
\feastday{{Sts. Cyprian \& Justina}}
\feastdate{26}{September}

\memorialCollect
\lett{O}{Lord,} let the protection of thy blessed Martyrs, Cyprian and Justina, continually defend us: forasmuch as thou failest not to look with mercy on those to whom thou dost grant the succour of thy assistance. Through.


\subbymemorial{27 September}{Sts. Cosmas and Damian}{M.}{MartyrRed}
\feastday{{Sts. Cosmas \& Damian}}
\feastdate{27}{September}

\memorialCollect
\lett{G}{rant,} we beseech thee, Almighty God: that we, who observe the heavenly birthday of thy holy Martyrs, Cosmas and Damian, may by their intercession be delivered from all evils that beset us. Through.


\subbymemorial{28 September}{Duke St. Wenceslaus}{M.}{MartyrRed}
\feastday{{St. Wenceslaus}}
\feastdate{28}{September}

\memorialCollect
\lett{O}{God,} who through the victory of martyrdom didst translate blessed Wenceslas from an earthly principality to the glory of heaven: defend us through his prayers from all adversity; and grant us likewise to rejoice in his fellowship. Through.


\festivalAO{First Class Double}{29 September}{St. Michael and All Angels}\label{Michaelmas}
\feastday{Michaelmas}
\feastdate{29}{September}

\openingSentence{And the smoke of the incense ascended up before God out of the angel's hand.}

\begin{paracol}{2}
\sloppy

\firstEvensong
\psalmAntiphon{While the Archangel {\dag} Michael contended with the dragon, a voice was heard of them that said: Salvation to our God, alleluia.}\\%MANUAL ADJUSTMENT

\firstEvensongHymn
\hymnLett{T}{hee,} O Christ, the Father's splendour,

\hymnSecondLine{Life and virtue of the heart,}

\begin{hangparas}{1.25em}{1}
In the presence of the Angels

Sing we now with tuneful art;

Meetly in alternate chorus

Bearing our responsive part.\\

2. Thus we praise with veneration

All the armies of the sky;

Chiefly him, the warrior Primate

Of celestial chivalry,

Michael, who in princely virtue

Cast Abaddon from on high.\\

3. By whose watchful care repelling,

King of everlasting grace,

Every ghostly adversary,

All things evil, all things base,

Grant us of thine only goodness

In thy Paradise a place.\\

4. Glory to the Father sing we

With resounding voices sweet,

Glory unto Christ our Saviour,

Glory to the Paraclete:

Standing forth, One God and Trinal,

Ere the ages; as is meet. Amen.\\
\end{hangparas}

℣. An Angel stood at the altar of the temple.

℟. Having in his hand a golden censer.

\firstEvensongAntiphon{While John beheld {\dag} the sacred mystery, the Archangel Michael sounded the trumpet: Forgive us, O Lord our God, that openest the book, and loosest the seals thereof, alleluia.}


\secondEvensong
\psalmAntiphon{Angels and Archangels, {\dag} Thrones and Dominions, Principalities and Powers, Virtues of the heavens, O praise the Lord of heaven, alleluia.}

\begin{rubric}
	Office Hymn of I Evensong.
\end{rubric}

℣. In the presence of the Angels I will sing praise unto thee, O my God.

℟. I will worship toward thy holy temple, and praise thy Name.

\secondEvensongAntiphon{O Prince most glorious, {\dag} Michael the Archangel, keep us in remembrance: here and everywhere, always, entreat the Son of God for us, alleluia, alleluia.}

\switchcolumn


\mattins
\psalmAntiphon{O Archangel Michael, {\dag} I have appointed thee prince over all the souls about to be received.}

\begin{rubric}
	Invitatory Hymn of I Evensong.
\end{rubric}

\mattinsOfficeHymn
\hymnLett{C}{hrist,} the fair glory of the holy Angels,

\hymnSecondLine{Thou who hast made us, thou who o’er us rulest,}

\begin{hangparas}{1.25em}{1}
Grant of thy mercy unto us thy servants

Steps up to heaven.\\

2. Send thy Archangel, Michael, to our succour;

Peacemaker blessed, may he banish from us

Striving and hatred, so that for the peaceful

All things may prosper.\\

3. Send thy Archangel, Gabriel, the mighty,

Herald of heaven; may he from us mortals

Spurn the old serpent, watching o’er the temples

Where thou art worshipp'd.\\

4. Send thy Archangel, Raphael, restorer

Of the misguided ways of men who wander,

Who at thy bidding strengthens soul and body

With thine anointing.\\

5. May the blest Mother of our God and Saviour,

May the assembly of the Saints in glory,

May the celestial companies of Angels

Ever assist us.\\

6. This he vouchsafe us, God for ever blessed,

Father eternal, Son, and Holy Spirit,

Whose is the glory which through all creation

Ever resoundeth. Amen.\\
\end{hangparas}

℣. An Angel stood at the altar of the temple.

℟. Having in his hand a golden censer.

\mattinsAntiphon{There was silence in heaven {\dag} while the dragon waged war: and Michael fought against him, and had the victory, alleluia.}

\fussy	
\end{paracol}


\collect
\lett{O}{everlasting} God, who hast ordained and constituted the services of Angels and men in a wonderful order; Mercifully grant that, as thy holy Angels always do thee service in heaven, so, by thy appointment, they may succour and defend us on earth. Through.


\festival{Greater Double}{30 September}{St. Jerome}{C.D.}{FestiveWhite}
\feastday{{St. Jerome}}
\feastdate{30}{September}

\begin{rubric}
	Common of a Confessor not a Bishop (p. \pageref{CommonConfessorNotBishop}).
\end{rubric}

%\collect
\lett{O}{God,} who for the exposition of the sacred Scriptures didst bestow upon thy Church blessed Jerome, thy Confessor and most illustrious Doctor: grant, we beseech thee; that, by the intercession of his merits, we may through thine assistance be enabled to perform those things which he taught both in word and deed. Through.

\begin{rubric}
    Commemoration \blackrubric{of St. Gregory of Armenia} (p. \pageref{ArmeniaCollect}).
\end{rubric}


%MANUAL ADJUSTMENT:
\clearpage
\subbymemorial{30 September}{St. Gregory of Armenia}{B.C.}{FestiveWhite}
\feastday{{St. Gregory Armenia}}
\feastdate{30}{September}

\memorialCollect\label{ArmeniaCollect}
\lett{O}{God,} who, through thy blessed Bishop and Confessor Gregory, didst grant that the people and King of Armenia to receive the light of true faith: likewise grant unto thy Church to rejoice in such mighty triumphs, and, by the merits and intercession of the same, to be succoured before thee. Through.


\subbymemorial{1 October}{St. Remigius (m.t.v.)}{B.C.}{FestiveWhite}
\feastday{{St. Remigius}}
\feastdate{1}{October}

\begin{rubric}
	Common of a Confessor Bishop (p. \pageref{ConfessorBishopCollect}).
\end{rubric}


\festivalAO{Greater Double}{2 October}{Holy Guardian Angels}
\feastday{Guardian Angels}
\feastdate{2}{October}

\begin{paracol}{2}
\sloppy

\firstEvensong
\psalmAntiphon{God shall give his Angels {\dag} charge over thee, to keep thee in all thy ways.}

\firstEvensongHymn
\hymnLett{T}{he} Guardians of our race, our Angel Guides we hail;

\hymnSecondLine{Our Father sendeth forth to aid our nature frail}

\begin{hangparas}{1.25em}{1}
These heav'nly friends, lest we should suffer overthrow

Through cunning of our subtle foe.\\

2. For he, who justly lost the honour once his own,

The traitor angel, rues his lost and vacant throne,

With burning envy strives to make them fall away

Whom God doth call to heav'nly day.\\

3. Then, watchful Guardian, spread thy wings and cleave the air,

Haste hither to our home committed to thy care;

Drive thence each noxious ill that might the soul infest,

Nor suffer danger here to rest.\\

4. Now to the holy Three your praise devoutly pour;

His glorious Godhead guides and governs evermore

This triple frame: to him ascribe we all our praise

Who reigns through everlasting days. Amen.\\
\end{hangparas}

℣. In the presence of the Angels I will sing praise unto thee, O my God.

℟. I will worship toward thy holy temple, and praise thy Name.

\firstEvensongAntiphon{They are all {\dag} ministering spirits, sent forth to minister to them who shall be heirs of salvation.}

\switchcolumn

\mattins
\psalmAntiphon{Their Angels {\dag} do always behold the face of my Father which is in heaven.}

\begin{rubric}
	Invitatory Hymn of I Evensong.
\end{rubric}

\mattinsOfficeHymn
\hymnLett{C}{reator} of the circling sky,

\hymnSecondLine{Who madest all by pow'r most high,}

\begin{hangparas}{1.25em}{1}
Thy Providence will never cease

To rule thy works in might and peace.\\

2. Be present when we cry to thee,

A sinful people though we be;

And as the day-dawn grows apace

Illume our minds with light of grace.\\

3. O send thine Angel thitherward

Assign'd by thee to be our guard,

That now his presence may begin

To keep us from all stain of sin.\\

4. Let him destroy that hidden snare

The eager serpent doth prepare,

Lest we be taken in the net

Before our heedless bosoms set.\\

5. At his command let every fear

Of hostile foemen disappear;

Let civil strife give way to peace,

And pestilence and famine cease.\\

6. To God the Father glory be;

For those the Saviour setteth free,

Anointed by the Holy Ghost,

Are guarded by the Angel host. Amen.\\
\end{hangparas}

℣. In the presence of the Angels I will sing praise unto thee, O my God.

℟. I will worship toward thy holy temple, and praise thy Name.

\mattinsAntiphon{And the Angel {\dag} that talked with me came again, and waked me, as a man that is wakened out of his sleep.}

\fussy	
\end{paracol}

\secondEvensong
\psalmAntiphon{Blessed be God, {\dag} who hath sent his Angel, who delivered his servants who put their trust in him.}

\begin{rubric}
	Office Hymn \& Versicle of I Evensong.
\end{rubric}

\secondEvensongAntiphon{Ye holy Angels, {\dag} our Guardians, defend us in time of battle, lest we perish in the dreadful day of judgment.}

\collect
\lett{O}{God,} who of thy ineffable providence dost vouchsafe to send thy holy Angels to be our guardians: grant unto us thy humble servants; that we may ever both be defended by their protection, and rejoice in their everlasting fellowship. Through.



%MANUAL ADJUSTMENT:
\clearpage
\festival{Double}{5 October}{Sts. Placidus and Companions}{M.}{MartyrRed}
\feastday{{Sts. Placidus \&c.}}
\feastdate{5}{October}

\begin{rubric}
	Common of Many Martyrs (p. \pageref{CommonMartyrs}).
\end{rubric}


\festivalLL{Greater Double}{7 October}{Holy Rosary of the Blessed Virgin Mary}
\feastday{Holy Rosary B.V.M.}
\feastdate{7}{October}

\begin{paracol}{2}
\sloppy

\firstEvensong
\psalmAntiphon{Mighty Virgin, {\dag} like the tower of David, whereon there hang a thousand bucklers, all shields of mighty men.}

\firstEvensongHymn
\hymnLett{T}{he} mysteries of God are shewn,

\hymnSecondLine{By heaven’s Messenger made known,}

\begin{hangparas}{1.25em}{1}
Who hails the Maid of David’s race,

God's Virgin Mother, full of grace.\\

2. The Virgin visits then her kin,

John’s mother; and the babe within

Th' enclosing womb exults to shew

That Christ is present here below.\\

3. The Word before all worlds begot

Of God the Father’s timeless thought

Is born a Child, to mortal doom,

From out the Virgin Mother’s womb.\\

4. Presented to the Lord with awe,

The Lawgiver obeys the Law:

He gives himself; the Ransomer

In poverty is ransom'd there.\\

5. The Son whose loss she mourned sore

The joyous Mother finds once more

Expounding mysteries unknown

To minds more learned than his own.\\

6. All honour, laud, and glory be,

O Jesu, Virgin-born, to thee;

All glory, as is ever meet,

To Father and to Paraclete. Amen.\\
\end{hangparas}

℣. Queen of the holy Rosary, pray for us.

℟. That we may be made worthy of the promises of Christ.

\firstEvensongAntiphon{Blessed art thou, {\dag} O Virgin Mary, Mother of God, for thou hast believed the Lord: for those things are accomplished in thee which were told thee: intercede for us unto the Lord our God.}

\switchcolumn

\secondEvensong
\psalmAntiphon{The Virgin Mary {\dag} is exalted above the choirs of Angels: and upon her head a crown of twelve stars.}

\secondEvensongHymn
\hymnLett{T}{hy} joys exultant we relate,

\hymnSecondLine{Thy wounding sorrows celebrate;}

\begin{hangparas}{1.25em}{1}
Thy glory, rob'd in endless light,

We sing, O Virgin Mother bright.\\

2. Hail, joyous Mother, blest indeed

When thou conceivedst, and in need

Didst visit; and thy Son for men

Didst bear, present, and find again.\\

3. Hail, first of martyrs; in thy pure

And anguish'd heart, thou didst endure

The agony, the scourge, the thorn,

The bitter Cross thy Son hath borne.\\

4. Hail, shining Queen, in triumphs won

For ever by thy glorious Son,

Who sent in flame the Holy Ghost,

Who thron'd thee o’er the heav'nly host.\\

5. Come, all ye nations, from this vine

Of mysteries the roses twine

In garlands for your Queen above,

The glorious Mother of fair love.\\

6. All honour, laud, and glory be,

O Jesu, Virgin-born, to thee;

All glory, as is ever meet,

To Father and to Paraclete. Amen.\\
\end{hangparas}

℣. Queen of the holy Rosary, pray for us.

℟. That we may be made worthy of the promises of Christ.

\secondEvensongAntiphon{O blessed Mother {\dag} and spotless Virgin, thou glorious Queen of the world, may all acknowledge the help of thy prayer who celebrate the solemnity of thy holy Rosary.}

\fussy	
\end{paracol}

\mattins
\psalmAntiphon{Be joyful, {\dag} Virgin Mother; Christ hath risen from the tomb.}

\begin{paracol}{2}
\sloppy

\mattinsInvitatoryHymn
\hymnLett{T}{he} Mount of Olives witnesseth

\hymnSecondLine{The awful agony of God:}

\begin{hangparas}{1.25em}{1}
His soul is sorrowful to death,

His sweat of blood bedews the sod.\\

2. And now the traitor’s work is done:

The clam'rous crowds around him surge;

Bound to pillar, O God the Son

Quivers beneath the blood-red scourge.\\

3. Lo! clad in purple soiled and worn,

Meekly the Saviour waiteth now

While wretches plait the cruel thorn

To crown with shame his royal brow.\\

4. Sweating and sighing, faint with loss

Of what hath flowed from life’s red fount,

He bears th' exceeding heavy Cross

Up to the verge of Calv'ry’s mount.\\

5. Nail'd to the wood of ancient curse,

Between two thieves the Sinless One

Still praying for his murderers,

Breathes forth his soul, and all is done!\\

6. Glory to thee, and honour meet,

Jesu, of Maiden-Mother born,

And Father and the Paraclete,

Through endless ages of the morn! Amen.	
\end{hangparas}

\switchcolumn

\mattinsOfficeHymn
\hymnLett{N}{ow} hell is vanquish'd; every chain

\hymnSecondLine{Of sin is broken; Christ again}

\begin{hangparas}{1.25em}{1}
Returning, victor over death,

The gates of heaven openeth.\\

2. We mortals saw him, till he pass'd

Into the heavens, where at last,

Partaker of God’s glory bright,

He sitteth on the Father’s right.\\

3. From thence he sheds the promis'd boon,

His Holy Spirit, on his own

In fiery tongues of love, o’erspread

Above each sad disciple’s head.\\

4. The Virgin, from the flesh set free,

Is borne beyond the stars; where she

Receives from heaven’s joyous throngs

The welcome of angelic songs.\\

5. Twice six the stars that crown her brow;

The gracious Mother reigneth now

Beside her Son’s eternal throne

O'er all creation as her own.\\

6. All honour, laud, and glory be,

O Jesu, Virgin-born, to thee;

All glory, as is ever meet,

To Father and to Paraclete. Amen.\\
\end{hangparas}

℣. God hath chosen her and preferred her.

℟. He hath made her to dwell in his tabernacle.

\mattinsAntiphon{To-day let us celebrate {\dag} devoutly the Solemnity of the holy Rosary of Mary, Mother of God, that she may intercede for us unto Jesus Christ the Lord.}

\fussy	
\end{paracol}


\collect
\lett{O}{God,} whose only-begotten Son by his life, death, and resurrection hath purchased for us the rewards of everlasting salvation: grant, we beseech thee; that we, who meditate upon these mysteries in the most sacred Rosary of the blessed Virgin Mary, may both imitate those things which they set forth, and attain unto those things which they promise. Through.
\begin{rubric}
	Commemoration \blackrubric{of St. Mark of Rome} (p. \pageref{MarkRomeCollect}).
\end{rubric}

\begin{rubric}
	Commemoration \blackrubric{of Sts. Sergius, Bacchus, Marcellus, \& Apuleius} (p. \pageref{SergiusCollect}).
\end{rubric}


%MANUAL ADJUSTMENT:
\clearpage
\subbymemorial{7 October}{Pope St. Mark of Rome}{B.C.}{FestiveWhite}
\feastday{{St. Mark Rome}}
\feastdate{7}{October}

\memorialCollect\label{MarkRomeCollect}
\lett{G}{raciously} hear our prayers, O Lord: and at the intercession of blessed Mark, thy Confessor and Bishop, mercifully grant us pardon and peace. Through.


\subbymemorial{7 October}{Sts. Sergius, Bacchus, Marcellus, and Apuleius}{M.}{MartyrRed}
\feastday{{Sts. Sergius \&c.}}
\feastdate{7}{October}

\memorialCollect\label{SergiusCollect}
\lett{M}{ay} the blessed merits of thy holy Martyrs Sergius, Bacchus, Marcellus, and Apuleius uphold us, O Lord: and ever make us fervent in thy love. Through.


\subbymemorial{9 October}{Sts. Denys, Rusticus, and Eleutherius}{M.}{MartyrRed}
\feastday{{Sts. Denys \&c.}}
\feastdate{9}{October}

\memorialCollect
\lett{O}{God,} who on this day didst strengthen blessed Denys, thy Martyr and Bishop, with the virtue of constancy in his passion, and didst vouchsafe to join unto him Rusticus and Eleutherius for the preaching of thy glory to the Gentiles: grant us, we beseech thee; by their example, to despise for love of thee the prosperity of the world, and to fear none of its adversities. Through.


\festivalLL{Second Class Double}{11 October}{Motherhood of the Blessed Virgin Mary}
\feastday{Motherhood B.V.M.}
\feastdate{11}{October}

\begin{paracol}{2}
\sloppy

%MANUAL ADJUSTMENT:
\vspace{-1\baselineskip}

\firstEvensong
\psalmAntiphon{Blessed art thou, {\dag} O Virgin Mary, who hast borne the Creator of all.}

\begin{rubric}
	Office Hymn of the Common of the Blessed Virgin Mary (p. \pageref{CommonBVM}).
\end{rubric}

℣. Blessed art thou amongst women.

℟. And blessed is the fruit of thy womb.

\firstEvensongAntiphon{Let us celebrate with joy {\dag} the Motherhood of the blessed Mary ever Virgin.}

\switchcolumn


%MANUAL ADJUSTMENT:
\vspace{-1\baselineskip}


\secondEvensong
\psalmAntiphon{Blessed art thou, {\dag} O daughter, from the Lord: for by thee we have partaken of the fruit of life.}

\begin{rubric}
	Office Hymn of the Common of the Blessed Virgin Mary (p. \pageref{CommonBVM}).
\end{rubric}

℣. Blessed art thou amongst women.

℟. And blessed is the fruit of thy womb.

\secondEvensongAntiphon{Thy motherhood, O Virgin Birthgiver of God, {\dag} announces joy unto the whole world: for from thee the Sun of Righteousness hath arisen, Christ our God.}

\fussy	
\end{paracol}

\mattins
\psalmAntiphon{Since I was humble, {\dag} I was pleasing to the Most High, and from my womb I have given birth to God and man.}

\begin{paracol}{2}
\sloppy

\mattinsInvitatoryHymn
\hymnLett{T}{he} world’s Redeemer from the earth

\hymnSecondLine{Up-bore the Virgin to the sky,}

\begin{hangparas}{1.25em}{1}
The stainless womb that gave him birth,

And thron'd her as the Queen on high.\\

2. In that white breast that knew no stain

Salvation’s hope was rob'd in clay,

The Christ that on the Cross was slain,

Whose blood has wash'd our sins away.\\

3. Let joy and hope to man be won,

And drive away all anxious fears

For Mary to her pitying Son

Will sweetly bear our prayers and tears.\\

4. The mother’s words her tender Child

Will heed and each entreaty bless;

Revere and love that mother mild,

And seek her aid in all distress.\\

5. Thou Triune God, all praise to thee,

That to the stainless bosom bore

The virginal maternity;

We sing thy glory evermore. Amen.
\end{hangparas}

\switchcolumn

\mattinsOfficeHymn
\hymnLett{S}{weet} mother of the Lord most high,

\hymnSecondLine{To thee we bow in humble prayer,}

\begin{hangparas}{1.25em}{1}
To thee from evil pow'rs we fly;

O shield and keep us in thy care.\\

2. It was to lift our fallen race

Above the curse of A-dam's crime,

The King bestow'd on thee his grace

And shap'd thy motherhood sublime.\\

3. So, Mother unto thee we pray;

Thou seest our need; thy Son entreat

That he, his anger turn'd away,

May raise our souls in mercy sweet.\\

4. All glory, Jesus, unto thee,

Born of the Virgin void of stain;

The same to Sire and Spirit be

Proclaim'd through one eternal reign. Amen.\\
\end{hangparas}

℣. The root of Jesse hath budded forth: a star hath arisen out of Jacob.

℟. The Virgin hath brought forth the Saviour: we praise thee, O our God.

\mattinsAntiphon{O holy Mary, {\dag} help thou the suffering, strengthen the faint-hearted, comfort the sorrowful; pray for the people, entreat for the clergy, intercede for all womankind vowed unto God: may all experience thy help, who celebrate thine admirable motherhood.}

\fussy	
\end{paracol}


\collect
\lett{O}{God,} who wast pleased that thy Word should take flesh of the womb of the Blessed Virgin Mary at the message of an Angel: grant to thy humble servants; that we, who believe her to be truly the Mother of God, may be aided by her intercession in thy sight. Through the same.


\subbymemorial{12 October}{St. Wilfred}{B.C.}{FestiveWhite}
\feastday{{St. Wilfred}}
\feastdate{12}{October}

\memorialCollect
\lett{O}{God,} by whose grace the blessed Bishop Wilfrid did wondrously shine forth with the glorious tokens of his merits: mercifully grant unto us; that we, who by his doctrine are taught to seek after things heavenly, may ever be defended by his advocacy. Through.


\subbymemorial{13 October}{King St. Edward (m.t.v.)}{C.}{FestiveWhite}
\feastday{{St. Edward}}
\feastdate{13}{October}

\memorialCollect
\lett{O}{God,} who didst bestow upon thy blessed Confessor King Edward the crown of everlasting glory: grant us, we beseech thee; so to venerate him on earth, that we may be enabled to reign with him in heaven. Through.


\festival{Double}{14 October}{Pope St. Callistus}{B.M.}{MartyrRed}
\feastday{{St. Callistus}}
\feastdate{14}{October}

\begin{rubric}
	Common of a Martyr (p. \pageref{CommonMartyr}).
\end{rubric}

%\collect
\lett{O}{God,} who seest that we fail by reason of our infirmity: through the examples of thy Saints mercifully restore us to the love of thee. Through.

\lettspace


%MANUAL ADJUSTMENT:
\clearpage
\festivalLL{Greater Double}{15 October}{Our Lady of Walsingham}
\feastday{Walsingham}
\feastdate{15}{October}

\openingSentence{How dreadful is this place! This is none other but the house of God, and this is the gate of heaven.}

\begin{paracol}{2}[]
\sloppy

\firstEvensong

\psalmAntiphon{Mary, Mother of the King of kings {\dag} appeared unto Richeldis, and shewed unto her the Holy House.}

\firstEvensongHymn
\hymnLett{H}{ail} Mary, ever blessed,

\hymnSecondLine{Of Walsingham the Queen!}

\hymnThirdLine{Through vision of Richeldis,}

\begin{hangparas}{1.25em}{1}
Thy favours there were seen.

When England was thy dowry,

There pilgrims bowed the knee.

At morn and noon and even,

They knelt to honour thee.\\

2. Hail Mary, ever blessed.

Thy children still delight

To tell abroad thy praises,

Thy miracles, thy might.

Still pilgrim feet are treading

Along the holy way.

Hostess of England's Nazareth,

Receive us home today!\\

3. Hail Mary, ever blessed.

The wells of water pure

Which mark thy holy places

Are signs that God doth cure

For sick of soul and body.

E’er since Richeldis' day,

They spring in benediction

Beside the Pilgrims' Way.\\

4. Hail Mary, ever blessed.

Thy name is great indeed;

For Jesus Christ our Saviour

Was in thy womb conceived.

Thy name be ever praised,

Increasing in this place,

And loud the angel's greeting:

`Hail Mary, full of grace!'\\

5. All glory, laud, and honour be

O Jesu, Virgin-born

Our worship is alone to thee

With Father and the Ghost.

Thou art our Mediator true,

Appealing for our sin;

Send unto us thine own Spirit,

Who makes us thy true kin. Amen.\\
\end{hangparas}

    ℣. O Mary, Queen of Heaven and of England, praise ye the Lord.

	℟. And bring us unto the majesty of thy Son.
	
	\firstEvensongAntiphon{O all ye Saints, {\dag} come unto the Holy Place of Walsingham, and bless ye the Lord.}


\secondEvensong

\begin{rubric}
	Office Hymn \& Versicle of I Evensong.
\end{rubric}

\psalmAntiphon{Rejoice, ye people of England, {\dag} for your land hath been sanctified by the Blessed Virgin.}

\secondEvensongAntiphon{O God, raise thou up again devotion to thy daughter: {\dag} Our blessed Lady of Walsingham.}

\switchcolumn

\mattins

\begin{rubric}
		Invitatory Hymn is the \blackrubric{Salve Regina} (p. \pageref{SalveRegina}).
\end{rubric}

\psalmAntiphon{O seek the Lord at Walsingham, {\dag} and quench thy thirst.}

\mattinsOfficeHymn

\hymnLett{W}{hen} Christ was born in Bethlehem,\par
\hymnSecondLine{The Angels worshipped and adored;}

\begin{hangparas}{1.25em}{1}
And kings and shepherds bowed the knee\par
To him their Saviour and their Lord.\\

2. For thirty years in Nazareth,\par
His Mother worshipped at the shrine\par
Of him who came in human form\par
To save the world, your Lord and mine.\\

3. For thirty years in Nazareth\par
He lived a hidden life of prayer;\par
And then went forth to work and die,\par
And every human sorrow share.\\

4. A thousand years have passed away;\par
Another Nazareth must rise\par
In England’s fair and lovely land,\par
Beneath those distant northern skies.\\

5. So Mary, God's own Mother blest,\par
Seeks out Richeldis, lady fair;\par
And bids her build at Walsingham\par
A second Nazareth of prayer.\\

6. Richeldis hastens to obey;\par
And soon the builder’s work is done:\par
The Shrine of England’s Nazareth\par
Proclaims the Mother and the Son.\\

7. But times did come when wicked hands\par
Were laid upon that holy place;\par
And men who knew not Mary's love\par
Did Walsingham's fair Shrine deface.\\

8. Four hundred years again have passed;\par
Once more the Mother of our Lord\par
Is shrined and honoured where her Son\par
Upon his Altar is adored.\\

9. So when you come to Walsingham,\par
Remember this, and mark it well:\par
It is to Nazareth you come,\par
Where Jesus, Mary, Joseph dwell.\\

10. All glory, laud, and honour be\par
To Jesu, Virgin-born, to thee;\par
All glory, as is ever meet,\par
To Father and to Paraclete. Amen.\\
\end{hangparas}

℣. Blessed is the holy Virgin Mary, and most worthy of all praise.

℟. Through her hath arisen the Sun of Justice, Christ our God, by whom we are saved and redeemed.

\mattinsAntiphon{Wicked men in their lust defied the Church, {\dag} and razed the Shrine of Our Lady.}
\end{paracol}


\collect
\lett{O}{God,} who in the blessed Virgin Mary didst make a fit dwelling-place for thy Son: grant, we beseech thee, that we who honour her shrine at Walsingham may also become temples of thy Holy Ghost. Through the same.


\festival{Second Class Double}{18 October}{St. Luke}{E.}{FestiveWhite}
\feastday{St. Luke}
\feastdate{18}{October}

\begin{rubric}
	Common of Apostles (p. \pageref{CommonApostles}).
\end{rubric}

%\collect
\lett{W}{e} beseech thee, O Lord: that thy holy Evangelist Luke may intercede for us: who for the honour of thy name continually bare in his own body the mortification of the Cross. Through.


\subbymemorial{19 October}{St. Frideswide}{V.}{FestiveWhite}
\feastday{{St. Frideswide}}
\feastdate{19}{October}

\memorialCollect
\lett{A}{lmighty} and everlasting God, the author of virtue and lover of virginity: grant us, we beseech thee; that, like as thy Virgin Frideswide was pleasing unto thee by the merit of chastity of life; so through her merits we may find favour in thy sight. Through.


\subbymemorial{21 October}{St. Hilarion}{Ab.}{FestiveWhite}
\feastday{{St. Hilarion}}
\feastdate{21}{October}

\begin{rubric}
	Common of a Confessor not a Bishop (p. \pageref{AbbotCollect}), Commemoration \blackrubric{of Sts. Ursula \& Companions,} from the Common of Virgins (p. \pageref{VirginMartyressesCollect}).
\end{rubric}


\subbymemorial{21 October}{Sts. Ursula and Companions}{V.M.}{MartyrRed}
\feastday{{Sts. Ursula \&c.}}
\feastdate{21}{October}

\begin{rubric}
	Common of Virgins (p. \pageref{VirginMartyressesCollect}).
\end{rubric}


%MANUAL ADJUSTMENT:
\clearpage
\festivalAO{Greater Double}{24 October}{St. Raphael the Archangel}
\feastday{St. Raphael}
\feastdate{24}{October}

\begin{rubric}
	Hymns \& Versicles \blackrubric{of Michaelmas} (p. \pageref{Michaelmas}).
\end{rubric}

\firstEvensong
\psalmAntiphon{The Angel {\dag} Raphael was sent unto Tobit and Sara, that he might heal them.}

\firstEvensongAntiphon{I am Raphael {\dag} the Angel, that stand before the Holy One: wherefore bless ye God for ever, and confess ye all his great and wonderful doings.}

\mattins
\psalmAntiphon{Be of good courage {\dag} O Tobit; for the time is nigh at hand that God will heal thee.}

\mattinsAntiphon{I am Raphael {\dag} the Angel, that stand before the Holy One: wherefore bless ye God for ever, and confess ye all his great and wonderful doings.}

\secondEvensong
\psalmAntiphon{Bless ye the God of heaven, {\dag} and praise him in the sight of all that magnify him and praise him for the merciful things which he hath done unto you.}

\secondEvensongAntiphon{O Prince most glorious, {\dag} Raphael the Archangel, keep us in remembrance: here and everywhere, always, entreat the Son of God for us.}

\collect
\lett{O}{God,} who didst give blessed Raphael the Archangel unto thy servant Tobias for a companion on his way: grant to us thy servants; that we may ever be guarded by his protection and strengthened by his help. Through.


\subbymemorial{25 October}{Sts. Chrysanthus and Daria}{M.}{MartyrRed}
\feastday{{\small Sts. Chrysanthus \& Daria}}
\feastdate{25}{October}

\memorialCollect
\lett{W}{e} beseech thee, O Lord, that the prayers of thy blessed Martyrs Chrysanthus and Daria may assist us: that, as we venerate them with our homage, so we may ever feel their loving help. Through.


%MANUAL ADJUSTMENT:
\clearpage
\subbymemorial{26 October}{Pope St. Evaristus of Rome}{B.M.}{MartyrRed}
\feastday{{St. Evaristus}}
\feastdate{26}{October}

\begin{rubric}
	Common of a Martyr (p. \pageref{CommonMartyrBishopCollect}).
\end{rubric}

\begin{rubric}
	If to-day be Saturday, the anticipated Vigil of Sts. Simon and Jude is kept (p. \pageref{SimonJudeVigil}).
\end{rubric}


\festival{Vigil}{27 October}{Vigil of Sts. Simon and Jude}{A.}{LentenViolet}\label{SimonJudeVigil}
\feastday{{Sts. Simon \& Jude Vigil}}
\feastdate{27}{October}

%\collect
\lett{G}{rant,} we beseech thee, Almighty God: that as we prevent the glorious birthday of thine Apostles Simon and Jude; so they may prevent us in the sight of thy Majesty, for the obtaining of thy blessings. Through.


\festival{Second Class Double}{28 October}{Sts. Simon and Jude}{A.}{FestiveWhite}
\feastday{Sts. Simon \& Jude}
\feastdate{28}{October}

\begin{rubric}
	Common of Apostles (p. \pageref{CommonApostles}).
\end{rubric}

%\collect
\lett{A}{lmighty} God, who hast built thy Church upon the foundation of the Apostles and Prophets, Jesus Christ himself being the head corner-stone; Grant us so to be joined together in unity of spirit by their doctrine, that we may be made an holy temple acceptable unto thee. Through.


\festivalLL{First Class Double}{Last Sunday in October}{Our Lord Jesus Christ the King}
\feastday{Christ the King}
%MANUAL ADJUSTMENT:
\fancyhead[RE,LO]{}

\begin{paracol}{2}
\sloppy

%MANUAL ADJUSTMENT:
\vspace{-1\baselineskip}

\firstEvensong
\psalmAntiphon{He shall be called {\dag} the Peaceful One, and his throne shall be established for ever and ever.}

\firstEvensongHymn
\hymnLett{O}{Lord} of ages, thee we sing;

\hymnSecondLine{We hail thee as the nation's King,}

\begin{hangparas}{1.25em}{1}
O Christ, our only Judge thou art,

Thou Searcher of the mind and heart.\\

2. Though evil shouting mobs maintain:

We will not have this Christ to reign,

Not such is our triumphant cry

Who hail thee King of kings most high.\\

3. O Christ, thou Prince of peace, subdue

All rebel thoughts; our minds renew;

And gather into thy one fold

The wand'rers whom thy love doth hold.\\

4. For this thou hangedst on the Tree

With arms outspread in fervent plea;

And shewedst forth thy burning heart

Transfixed by the cruel dart.\\

5. For this thy pierced side doth shed,

Beneath the forms of wine and bread,

Thy grace upon thy sons; yet thou

Art hidden at the Altar now.\\

6. The rulers of the nations raise

To thee their meed of public praise;

Instructors, judges, thee confess;

Art, science, law, thy truth express.\\

7. Let kings be fain to dedicate

To thee the emblems of their state;

And let our homes and fatherland

Be subject to thy kind command.\\

8. All glory, Lord, to thee, whose sway

The world's dominion doth obey;

All glory, as is ever meet,

To Father and to Paraclete. Amen.\\
\end{hangparas}

℣. All power is given unto me.

℟. In heaven and in earth.

\firstEvensongAntiphon{The Lord God shall give unto him {\dag} the throne of his father David: and he shall reign over the house of Jacob for ever; and of his kingdom there shall be no end, alleluia.}

\secondEvensong
\psalmAntiphon{Living waters shall go out {\dag} from Jerusalem; and the Lord shall be King over all the earth.}

\begin{rubric}
	Office Hymn of I Evensong.
\end{rubric}

℣. His kingdom shall be increased.

℟. And of his peace there shall be no end.

\secondEvensongAntiphon{He hath on his vesture {\dag} and on his thigh a name written, King of kings, and Lord of lords: to him be glory and dominion for ever and ever.}

\collect
\lett{A}{lmighty} and everlasting God, who in thy beloved Son, the King of all, hast been pleased to make all things new: mercifully grant; that all the families of the Gentiles, dispersed by the wounds of sin, may be made subject to his most gracious governance. Who liveth.
\begin{rubric}
    Commemoration of the Sunday occurring.
\end{rubric}

\switchcolumn

%MANUAL ADJUSTMENT:
\vspace{-1\baselineskip}

\mattins
\psalmAntiphon{The God of heaven {\dag} shall set up a kingdom which shall break in pieces and consume all kingdoms, and it shall stand for ever.}

\mattinsInvitatoryHymn
\hymnLett{I}{mage} Eterne of God Most High,

\hymnSecondLine{Thou Light of Light, True God, to thee,}

\begin{hangparas}{1.25em}{1}
Redeemer, laud and glory be,

And kingly reign o’er earth and sky.\\

2. For thou alone, ere Time began,

Its Hope and central-point to be

The Father justly granted thee

To rule each nation, tribe, or clan.\\

3. O Flower of a Virgin-birth

O Head of all on earth who dwell,

O Stone that from the mountain fell

And with its vastness cover'd earth!\\

4. The race of men, condemn'd to lie

Beneath the direful tyrant’s yoke,

By thee at length the shackles broke

And claim'd the Fatherland on high.\\

5. Lawgiver, Priest, and Teacher, God;

With these the title well accords

Of King of kings and Lord of lords

Upon thy vesture writ in Blood.\\

6. With grateful hearts thy rule we bless

Who justly reignest over all:

Them only truest joys befall

Who thee as King and Lord confess.\\

7. To thee, O Jesus, ruling o’er

Earth's rulers all, be glory meet,

With Father and the Paraclete,

Throughout the ages evermore! Amen.
\end{hangparas}

\mattinsOfficeHymn
\hymnLett{N}{ow} Christ unfurls, in triumph high,

\hymnSecondLine{His glorious banner to the sky:}

\begin{hangparas}{1.25em}{1}
Ye suppliant nations, kneel and praise

The King of kings with joyful lays.\\

2. He hath not won his kingdom here

By devastation, force, or fear;

But on the Cross uplifted high

By love alone draws all men nigh.\\

3. How trebly blessed is the land

Obedient unto Christ’s command,

Which urges laws that prove the worth

Of heav'nly edicts here on earth!\\

4. No arm'd rebellion kindles there,

Peace strengthens union everywhere,

And concord smiles; upon all sides

The civil order safe abides.\\

5. There married faith is kept secure;

There rip'ning youth is ever pure;

And modest households flourish, fair

With sweet and homely virtues, there.\\

6. Pour down that long'd for light of thine

Upon us all, dear King divine;

And let the conquer'd world adore

In shining peace for evermore.\\

7. All glory, Lord, to thee, whose sway

The world’s dominion doth obey;

All glory, as is ever meet,

To Father and to Paraclete. Amen.\\
\end{hangparas}

℣. His kingdom shall be increased.

℟. And of his peace there shall be no end.

\mattinsAntiphon{Unto God and his Father {\dag} hath he made us to be a kingdom, who is the first begotten of the dead, and the Prince of the kings of the earth, alleluia.}

\fussy	
\end{paracol}


\festivalAO{Vigil}{31 October}{Hallowe'en}
\feastday{{Hallowe'en}}
\feastdate{31}{October}

%\collect
\lett{O}{Lord,} our God, multiply upon us thy grace: and grant, that by our holy profession we may follow after the gladness of them, whose glorious festival we prevent. Through.


\festivalAO{First Class Double with a Common Octave}{1 November}{All Hallows}\label{AllHallows}
\feastday{All Hallows}
\feastdate{1}{November}

\begin{paracol}{2}
\sloppy

\firstEvensong
\psalmAntiphon{I beheld a great multitude {\dag} which no man could number, of all nations, standing before the throne.}

\firstEvensongHymn
\hymnLett{O}{Christ,} Redeemer of us all,

\hymnSecondLine{Protect thy servants when they call,}

\begin{hangparas}{1.25em}{1}
And hear with reconciling care

The blessed Virgin's holy prayer.\\

2. And ye, O ever blissful throng

Of heav'nly Spirits, guardians strong,

Our past and present ills dispel,

From future peril shield us well.\\

3. Ye Prophets of the Judge ador'd,

Ye twelve Apostles of the Lord,

For us your ceaseless prayer outpour,

Salvation for our souls implore.\\

4. Martyrs of God, renown'd for aye,

Confessors rang'd in bright array,

Let all your orisons unite

To bear us to the realms of light.\\

5. O sacred Virgin choirs, may ye,

With Monks of holy ministry

And every Saint of Christ, obtain

That we his fellowship may gain.\\

6. From lands wherein thy faithful dwell

Drive far away the infidel;

So we to Christ due hymns of praise

Henceforth with eager hearts may raise.\\

7. To thee, O Father, born of none,

And thee, O sole-begotten Son,

One with the Holy Paraclete,

Be glory ever, as is meet. Amen.\\
\end{hangparas}

℣. Be glad, O ye righteous, and rejoice in the Lord.

℟. And be joyful, all ye that are true of heart.

\firstEvensongAntiphon{O ye Angels and Archangels, {\dag} O ye Thrones and Dominions, Principalities and Powers, Virtues of the heavens, Cherubim and Seraphim, O ye Patriarchs and Prophets, ye holy Doctors of the law, O all ye Apostles, ye Martyrs of Christ, ye holy Confessors, ye Virgins of the Lord, ye holy Hermits, and all Saints: offer for us your intercessions.}

\switchcolumn

\mattins
\psalmAntiphon{Thou hast redeemed us, {\dag} O Lord our God, by thy blood out of every kindred, tongue, and people, and nation; and hast made us unto our God a kingdom.}

\begin{rubric}
	Invitatory Hymn of I Evensong.
\end{rubric}

\mattinsOfficeHymn
\hymnLett{J}{esu,} who cam'st the world to save,

\hymnSecondLine{By thee redeem'd, thine aid we crave:}

\begin{hangparas}{1.25em}{1}
Mother of God, in time of need,

For contrite hearts salvation plead.\\

2. Let all the bright Angelic choirs,

The ranks of Patriarchal sires,

And Prophet Saints, devout and meek,

Forgiveness for our errors seek.\\

3. The Baptist, thy great harbinger,

With heav'n's appointed Keybearer,

And all Apostles strive to win

From God remission of our sin.\\

4. So may the sacred Martyr band,

Confessors, Priests, adoring stand,

And every Virgin chastely pray

That he would purge our guilt away.\\

5. Your suffrages, ye Monks, unite

With all the citizens of light

To speed our vows, combine your pow'rs

That life’s eternal prize be ours.\\

6. To God the Father, God the Son,

And God the Spirit, Three in One,

Laud, honour, might, and glory be

From age to age eternally. Amen.\\
\end{hangparas}

℣. Let the Saints be joyful with glory.

℟. Let them rejoice in their beds.

\mattinsAntiphon{The glorious company {\dag} of the Apostles praise thee; the goodly fellowship of the Prophets praise thee; the white-robed army of Martyrs praise thee; with one heart and voice do all the elect acknowledge thee: O blessed Trinity, one only God.}

\fussy	
\end{paracol}


\secondEvensong
\psalmAntiphon{All his Saints {\dag} shall praise him, even the people of Israel, even the people that serveth him: such honour have all his Saints.}

\begin{rubric}
	Office Hymn of I Evensong.
\end{rubric}

℣. Let the Saints be joyful with glory.

℟. Let them rejoice in their beds.

\secondEvensongAntiphon{O how glorious {\dag} is the kingdom wherein all the Saints rejoice with Christ; arrayed in white robes, they follow the Lamb whithersoever he goeth.}


\collect
\lett{A}{lmighty} and everlasting God, who in one solemnity hast vouchsafed unto us to venerate the merits of all thy Saints: we beseech thee; that, at the intercession of so great a multitude, thou wouldest bestow upon us, who entreat thee, the abundance of thy mercy. Through.

\begin{rubric}
	The Office is of All Hallows through its Octave, unless there occur a Sunday or Feast Day (Double or higher), any concurring Memorial being commemorated.
\end{rubric}


\festivalAO{Double}{2 November}{All Souls}\label{AllSouls}
\feastday{All Souls}
\feastdate{2}{November}

\begin{rubric}
	The Office of the Dead is said today with Double rite.
\end{rubric}


\subbymemorial{3 November}{St. Winifred}{V.M.}{MartyrRed}
\feastday{{St. Winifred}}
\feastdate{3}{November}

\begin{rubric}
	Common of Virgins (p. \pageref{VirginMartyressesCollect}).
\end{rubric}


\subbymemorial{4 November}{Sts. Vitalis and Agricola}{M.}{MartyrRed}
\feastday{{Sts. Vitalis \& Agricola}}
\feastdate{4}{November}

\memorialCollect
\lett{G}{rant,} we beseech thee, Almighty God: that we, who devoutly celebrate the festival of thy holy Martyrs Vitalis and Agricola, may be aided by their intercession with thee. Through.


\subbymemorialO{5 November}{St. Elizabeth, Mother of St. John Baptist}{FestiveWhite}
\feastday{{St. Elizabeth}}
\feastdate{5}{November}

\begin{rubric}
	Common of Holy Matrons (p. \pageref{CommonHolyMatronsCollect}).
\end{rubric}


\subbymemorial{7 November}{St. Willibrord}{B.C.}{FestiveWhite}
\feastday{{St. Willibrord}}
\feastdate{7}{November}

\begin{rubric}
	Common of a Confessor Bishop (p. \pageref{ConfessorBishopCollect}).
\end{rubric}


\festivalAO{Greater Double}{8 November}{Patriarchs and Prophets of the Old Law}
\feastday{{All Hallows Octave Day}}
\feastdate{8}{November}

\begin{rubric}
	Hymns, Versicles, \& Antiphons \blackrubric{of All Hallows} (p. \pageref{AllHallows}).
\end{rubric}

\lett{O}{God,} who didst show forth thy Divine Wisdom in ancient times through thy Patriarchs and Prophets: grant that we may imitate their example on earth; and enjoy their intercession in heaven. Through.

\begin{rubric}
	Commemoration \blackrubric{of All Hallows} (p. \pageref{AllHallows}) \& \blackrubric{of the Four Crowned Martyrs,} from the Common of Many Martyrs (p. \pageref{CommonMartyrsNotBishopsCollect}).
\end{rubric}


%MANUAL ADJUSTMENT:
\clearpage
\subbymemorial{8 November}{Four Crowned Martyrs}{M.}{MartyrRed}
\feastday{{4 Crowned Martyrs}}
\feastdate{8}{November}

\begin{rubric}
	Common of Many Martyrs (p. \pageref{CommonMartyrsNotBishopsCollect}).
\end{rubric}



\festivalLL{Greater Double}{9 November}{Dedication of the Basilica of St. Saviour}
\feastday{{St. Saviour Basilica}}
\feastdate{9}{November}

\begin{rubric}
	Common of a Dedication of a Church (p. \pageref{CommonDedication}).
\end{rubric}

\begin{rubric}
	Commemoration \blackrubric{of St. Theodore Tyro} (p. \pageref{TheodoreTyroCollect}).
\end{rubric}


\subbymemorial{9 November}{St. Theodore Tyro}{M.}{MartyrRed}
\feastday{{St. Theodore Tyro}}
\feastdate{9}{November}

\memorialCollect\label{TheodoreTyroCollect}
\lett{O}{God,} who dost encompass and protect us with the glorious confession of blessed Theodore, thy Martyr: grant us both to profit by his example, and to be supported by his intercession. Through.


\subbymemorial{10 November}{Sts. Tryphon, Respicius, and Nympha}{M.}{MartyrRed}
\feastday{{Sts. Tryphon \&c.}}
\feastdate{10}{November}

\memorialCollect
\lett{G}{rant,} we beseech thee, O Lord, that we, ever keeping the feast of thy holy Martyrs Tryphon, Respicius, and Nympha: may through their intercession enjoy the benefit of thy protection. Through.


\festival{Greater Double}{11 November}{St. Martin of Tours}{B.C.}{FestiveWhite}
\feastday{{St. Martin Tours}}
\feastdate{11}{November}

\begin{rubric}
	Common of a Confessor Bishop (p. \pageref{CommonConfessorBishop}).
\end{rubric}

\firstEvensong
\psalmAntiphon{O Lord, {\dag} if I am yet needful unto thy people; I refuse not to labour: thy will be done.}

\firstEvensongAntiphon{O blessed Martin, {\dag} whose righteous soul possesseth Paradise; whereat the Angels triumph, the Archangels are jubilant: thee the choir of Saints proclaimeth, the throng of Virgins inviteth, saying, Abide with us for ever.}

\mattins
\psalmAntiphon{With eyes and hands {\dag} ever raised toward heaven, his undaunted spirit ceased not from supplication, alleluia.}

\mattinsAntiphon{O blessed Martin, {\dag} whose righteous soul possesseth Paradise; whereat the Angels triumph, the Archangels are jubilant: thee the choir of Saints proclaimeth, the throng of Virgins inviteth, saying, Abide with us for ever.}


\secondEvensong
\psalmAntiphon{Martin is received, {\dag} rejoicing, into the bosom of Abraham: Martin, poor and lowly here on earth, entereth rich into heaven, and is saluted with heavenly songs.}

\secondEvensongAntiphon{O blessed Bishop, {\dag} who with all his heart loved Christ the King, and feared not the dominion of princes! O holy soul, which, although withheld from the persecutor's sword, lacked not the palm of martyrdom!}

\collect
\lett{O}{God,} who seest that we stand not in our own strength: mercifully grant; that, by the intercession of blessed Martin, thy Confessor and Bishop, we may be defended against all adversities. Through.

\begin{rubric}
    Commemoration \blackrubric{of St. Mennas,} from the Common of a Martyr (p. \pageref{CommonMartyrNotBishopCollect}).
\end{rubric}


\subbymemorial{11 November}{St. Mennas}{M.}{MartyrRed}
\feastday{{St. Mennas}}
\feastdate{11}{November}

\begin{rubric}
	Common of a Martyr (p. \pageref{CommonMartyrNotBishopCollect}).
\end{rubric}


\subbymemorial{12 November}{Pope St. Martin I of Rome}{B.M.}{MartyrRed}
\feastday{{Pope St. Martin I}}
\feastdate{12}{November}

%%CHECK:
\begin{rubric}
	Common of a Martyr (p. \pageref{CommonMartyrBishopCollect}), the second Collect.
\end{rubric}


\subbymemorial{13 November}{St. Britius of Tours}{B.C.}{FestiveWhite}
\feastday{{St. Britius Tours}}
\feastdate{13}{November}

\begin{rubric}
	Common of a Confessor Bishop (p. \pageref{ConfessorBishopCollect}).
\end{rubric}


\subbymemorial{14 November}{St. Gregory Palamas}{B.C.}{FestiveWhite}
\feastday{{St. Gregory Palamas}}
\feastdate{14}{November}

\begin{rubric}
	Common of a Confessor Bishop (p. \pageref{ConfessorBishopCollect}).
\end{rubric}


%MANUAL ADJUSTMENT:
\clearpage
\subbymemorial{17 November}{St. Gregory the Wonder-worker}{B.C.}{FestiveWhite}
\feastday{{\small St. Gregory Wonder-worker}}
\feastdate{17}{November}

\begin{rubric}
	Common of a Confessor Bishop (p. \pageref{ConfessorBishopCollect}).
\end{rubric}


\subbymemorial{17 November}{St. Hilda of Whitby}{V.}{FestiveWhite}
\feastday{{St. Hilda Whitby}}
\feastdate{17}{November}

\memorialCollect
\lett{G}{rant,} we beseech thee, Almighty God: that we, who rejoice in the yearly solemnity of blessed Hilda thy Virgin, may by her intercession be led from our old nature to newness of life. Through.


\festivalAO{Greater Double}{18 November}{Dedication of the Basilica of the Holy Apostles Peter and Paul}
\feastday{{Basilica of Peter \& Paul}}
\feastdate{18}{November}

\begin{rubric}
	Common of the Dedication of a Church (p. \pageref{CommonDedication}).
\end{rubric}


\subbymemorial{19 November}{Pope St. Pontianus of Rome}{B.M.}{MartyrRed}
\feastday{{Pope St. Pontianus}}
\feastdate{19}{November}

\begin{rubric}
	Common of a Martyr (p. \pageref{CommonMartyrBishopCollect}).
\end{rubric}


\festival{Double}{20 November}{King St. Edmund}{M.}{MartyrRed}
\feastday{{St. Edmund}}
\feastdate{20}{November}

\begin{rubric}
	Common of a Martyr (p. \pageref{CommonMartyr}).
\end{rubric}

%\collect
\lett{O}{God} of unspeakable mercy, who didst give grace to the most blessed King Edmund to overcome the enemy by dying for thy name: mercifully grant unto this thy family; that, by his intercession they may be found worthy to vanquish and destroy in themselves the temptations of their ancient foe. Through.


\festivalLL{Greater Double}{21 November}{Presentation of the Blessed Virgin Mary}
\feastday{{Presentation B.V.M.}}
\feastdate{21}{November}

\begin{rubric}
	Common of the Blessed Virgin Mary (p. \pageref{CommonBVM}).
\end{rubric}

\begin{rubric}
	\textsc{Note,} The following Antiphon is said during I \& II Evensong.
\end{rubric}

\antiphon{Mag.}{O ever blessed Mother of God, {\dag} Mary ever Virgin, temple of the Godhead, hallowed shrine of the Holy Spirit: thou only, above all others, wast acceptable to our Lord Jesus Christ, alleluia.}

\collect
\lett{O}{God,} who didst will that blessed Mary ever Virgin, the dwelling-place of Holy Ghost, should this day be presented in the temple: grant, we beseech thee; that through her intercession we may be found worthy to be presented in the temple of thy glory. Through . . . in the unity of the same.


\subbymemorial{21 November}{St. Gelasius}{B.C.}{FestiveWhite}
\feastday{{St. Gelasius}}
\feastdate{21}{November}

\begin{rubric}
	Common of a Confessor Bishop (p. \pageref{ConfessorBishopCollect}).
\end{rubric}


\subbymemorial{21 November}{St. Columbanus}{Ab.}{FestiveWhite}
\feastday{{St. Columbanus}}
\feastdate{21}{November}

\begin{rubric}
	Common of a Confessor not a Bishop (p. \pageref{ConfessorNotBishopCollect}).
\end{rubric}


\festival{Double}{22 November}{St. Cecilia}{V.M.}{MartyrRed}
\feastday{{St. Cecilia}}
\feastdate{22}{November}

\begin{rubric}
	Common of a Virgin (p. \pageref{CommonVirgin}).
\end{rubric}

\firstEvensong
\psalmAntiphon{Then did Valerian {\dag} find Cecelia in her chamber with an Angel, making her supplication.}

\firstEvensongAntiphon{It is secret, {\dag} O Valerian, what now I wish to tell thee: I have an Angel of God who loves me, and guards my body with exceeding zeal.}

\mattins
\psalmAntiphon{Cecilia {\dag} thine handmaid, O Lord, ever serveth thee as zealously as a busy bee.}

\mattinsAntiphon{When the dawn of day was breaking, {\dag} Cecilia the Martyr cried, saying: Forward, soldiers of Christ! Cast off the works of darkness, and put on the armour of light.}

\secondEvensong
\psalmAntiphon{I bless thee, {\dag} O Father of my Lord Jesus Christ: for by thy Son the fire is quenched about thine handmaid.}

\secondEvensongAntiphon{This noble maiden {\dag} alway bare the glorious Gospel of Christ in her bosom: without ceasing, she continued in prayer and supplication, and celestial converse, night and day.}

\collect
\lett{O}{God,} who makest us glad with the yearly solemnity of blessed Cecilia thy Virgin and Martyress: grant, that we do venerate her in our service, so we may follow the example of her godly conversation. Through.


\festival{Double}{23 November}{Pope St. Clement of Rome}{B.M.}{MartyrRed}
\feastday{{St. Clement}}
\feastdate{23}{November}

\begin{rubric}
	Common of a Martyr (p. \pageref{CommonMartyr}).
\end{rubric}

\firstEvensong
\psalmAntiphon{While holy Clement {\dag} prayed, the Lamb of God appeared unto him.}\\

℣. Thou hast crowned him with glory and honour, O Lord.

℟. And madest him to have dominion of the works of thy hands.

\firstEvensongAntiphon{Let us all pray {\dag} to our Lord Jesus Christ, that he would open a fountain for his Confessors.}


\mattins
\psalmAntiphon{Upon the mountain {\dag} I saw the Lamb standing, from beneath whose feet floweth a fount of living waters.}\\

℣. Let the Saints be joyful with glory.

℟. Let them rejoice in their beds.

\mattinsAntiphon{When he had started going to the sea, {\dag} the people cried out with loud voices, O Lord Jesus Christ, save him: and Clement responded, weeping, Heavenly Father, receive my spirit.}


\secondEvensong
\psalmAntiphon{All the nations {\dag} round about believed in Christ the Lord.}\\

℣. Let the Saints be joyful with glory.

℟. Let them rejoice in their beds.

\secondEvensongAntiphon{In the heavenly kingdom {\dag} rejoice the souls of the Blessed, who followed the footsteps of Christ their Master: and since for love of him they poured forth their life-blood, therefore with Christ do they exult for ever.}

\collect
\lett{O}{God,} who makest us glad with the yearly solemnity of blessed Clement thy Martyr and Bishop: mercifully grant; that as we now celebrate his birthday, so we may imitate his constancy in suffering. Through.

\begin{rubric}
    Commemoration \blackrubric{of St. Felicitas} (p. \pageref{FelicityCollect}).
\end{rubric}


\subbymemorial{23 November}{St. Felicitas}{M.}{MartyrRed}
\feastday{{St. Felicitas}}
\feastdate{23}{November}

\memorialCollect\label{FelicityCollect}
\lett{G}{rant,} we beseech thee, Almighty God: that we, celebrating the festival of blessed Felicitas, thy Martyress, may be protected by her merits and prayers. Through.


\subbymemorial{24 November}{St. Chrysogonus}{M.}{MartyrRed}
\feastday{{St. Chrysogonus}}
\feastdate{24}{November}

\memorialCollect
\lett{A}{ssist} us, O Lord, in these our supplications: that we, who acknowledge ourselves to be guilty by reason of our iniquity, may be delivered by the intercession of thy blessed Martyr Chrysogonus. Through.


\festival{Double}{25 November}{St. Catherine of Alexandria}{V.M.}{MartyrRed}
\feastday{{\small St. Catherine Alexandria}}
\feastdate{25}{November}

\begin{rubric}
	Common of a Virgin (p. \pageref{CommonVirgin}).
\end{rubric}

%\collect
\lett{O}{God,} who didst give the law to Moses on the height of Mount Sinai, and in the same place didst through thy holy Angels wondrously bestow the body of blessed Catherine, thy Virgin and Martyress: grant, we beseech thee; that by her merits and intercession we may be enabled to attain unto that mount which is Christ. Who liveth.

\subbymemorial{26 November}{St. Peter of Alexandria}{B.M.}{MartyrRed}
\feastday{{St. Peter Alexandria}}
\feastdate{26}{November}

\begin{rubric}
	Common of a Martyr (p. \pageref{CommonMartyrBishopCollect}).
\end{rubric}

